\paragraph{二原子经验谱、临界容量反演、Bernoulli 残差比与奇素数缺陷尾谱}
在 bin-fold 的两态一致渐近、critical-scale 容量极限曲线、fold-gauge 常数的 Bernoulli 递归剥离以及输出位势的奇素数单位根扭曲四阶展开已经建立之后，还可继续抽取出下列七条后果。它们分别刻画稳定层均匀经验谱的严格两点弱极限、连续容量极限曲线的分布二阶导数、目标成功率下的最优辅助比特预算、由 \((\log C_m)\) 序列恢复 Bernoulli 塔时的 Richardson 残差比、奇素数模 RH 缺陷的可和性与其 \(\ell^\alpha\) 临界指数，以及最小非平凡 prime-twist 谱半径对 \(\lambda\) 的严格二次接近律与最终单调逼近。

\begin{theorem}[bin-fold 纤维谱经自然缩放后的稳定层两点弱极限]\label{thm:xi-time-part70-binfold-uniform-empirical-two-point-law}
\leanverified{paper\_xi\_time\_part70\_binfold\_uniform\_empirical\_two\_point\_law}
定义尺度化多重度
\[
Z_m^{\mathrm{bin}}(w):=\left(\frac{\varphi}{2}\right)^m d_m^{\mathrm{bin}}(w)
\qquad (w\in X_m),
\]
以及稳定层上的均匀经验测度
\[
\nu_m:=\frac{1}{|X_m|}\sum_{w\in X_m}\delta_{Z_m^{\mathrm{bin}}(w)}.
\]
再记
\[
A_0^{(m)}:=\{w\in X_m:\ w_m=0\},
\qquad
A_1^{(m)}:=\{w\in X_m:\ w_m=1\}.
\]
则有弱收敛
\[
\nu_m\xrightarrow{\mathrm w}\nu_\infty
:=
\frac{1}{\varphi}\,\delta_1+\frac{1}{\varphi^2}\,\delta_{\varphi^{-1}}.
\]
等价地，对任意有界连续函数 \(f\)，都有
\[
\lim_{m\to\infty}\frac{1}{|X_m|}\sum_{w\in X_m}f\!\bigl(Z_m^{\mathrm{bin}}(w)\bigr)
=
\frac{1}{\varphi}f(1)+\frac{1}{\varphi^2}f(\varphi^{-1}).
\]
\end{theorem}
\begin{proof}
由定理 \ref{thm:fold-bin-two-state-asymptotic}，存在 \(\varepsilon_m\to 0\) 使得对全部 \(w\in X_m\) 一致有
\[
\bigl|Z_m^{\mathrm{bin}}(w)-\varphi^{-w_m}\bigr|\le \varepsilon_m.
\]
因此对充分大的 \(m\)，所有 \(Z_m^{\mathrm{bin}}(w)\) 都落在某个固定紧区间内。任取有界连续函数 \(f\)，其在该紧区间上一致连续。于是若定义
\[
\eta_{m,0}:=\sup_{w\in A_0^{(m)}}\bigl|f(Z_m^{\mathrm{bin}}(w))-f(1)\bigr|,
\qquad
\eta_{m,1}:=\sup_{w\in A_1^{(m)}}\bigl|f(Z_m^{\mathrm{bin}}(w))-f(\varphi^{-1})\bigr|,
\]
则 \(\eta_{m,0},\eta_{m,1}\to 0\)。从而
\[
\frac{1}{|X_m|}\sum_{w\in X_m}f\!\bigl(Z_m^{\mathrm{bin}}(w)\bigr)
=
\frac{|A_0^{(m)}|}{|X_m|}\bigl(f(1)+o(1)\bigr)
+
\frac{|A_1^{(m)}|}{|X_m|}\bigl(f(\varphi^{-1})+o(1)\bigr).
\]
再由稳定词计数
\[
|A_0^{(m)}|=F_{m+1},
\qquad
|A_1^{(m)}|=F_m,
\qquad
|X_m|=F_{m+2},
\]
以及 Fibonacci 比值极限
\[
\frac{F_{m+1}}{F_{m+2}}\to \frac{1}{\varphi},
\qquad
\frac{F_m}{F_{m+2}}\to \frac{1}{\varphi^2},
\]
便得到所示极限。证毕。
\end{proof}

\begin{theorem}[连续容量极限曲线的分布二阶导数恰为两原子测度]\label{thm:xi-time-part70-continuous-capacity-two-atom-second-derivative}
定义尺度化连续容量
\[
\widehat C_m(s):=
2^{-m}\,\mathcal C_m^{\mathrm{bin,cont}}\!\left(
s\left(\frac{2}{\varphi}\right)^m
\right)
\qquad (s\ge 0),
\]
并记其极限函数
\[
L(s):=\frac{\varphi}{\sqrt5}\min\{1,s\}+\frac{1}{\sqrt5}\min\{\varphi^{-1},s\}.
\]
则在分布意义下有精确恒等式
\[
-D^2L
=
\frac{1}{\sqrt5}\,\delta_{\varphi^{-1}}
+
\frac{\varphi}{\sqrt5}\,\delta_1.
\]
换言之，极限容量曲线只有两个不可消去的折点，其位置分别位于 \(s=\varphi^{-1}\) 与 \(s=1\)，对应原子质量分别为 \(1/\sqrt5\) 与 \(\varphi/\sqrt5\)。
\end{theorem}
\begin{proof}
这正是定理 \ref{thm:xi-foldbin-underresolution-two-atomic-curvature-measure} 中曲率测度公式的直接重写。若将
\[
L(s)=
\begin{cases}
\dfrac{\varphi^2}{\sqrt5}\,s, & 0\le s\le \varphi^{-1},\\[6pt]
\dfrac{\varphi}{\sqrt5}\,s+\dfrac{1}{\varphi\sqrt5}, & \varphi^{-1}\le s\le 1,\\[8pt]
1, & s\ge 1
\end{cases}
\]
逐段求导，则 \(L'\) 在 \(s=\varphi^{-1}\) 处的跃降为
\[
\frac{\varphi^2}{\sqrt5}-\frac{\varphi}{\sqrt5}=\frac{1}{\sqrt5},
\]
在 \(s=1\) 处的跃降为
\[
\frac{\varphi}{\sqrt5}-0=\frac{\varphi}{\sqrt5}.
\]
因此分布导数满足
\[
-D^2L
=
\frac{1}{\sqrt5}\,\delta_{\varphi^{-1}}
+
\frac{\varphi}{\sqrt5}\,\delta_1.
\]
证毕。
\end{proof}

\begin{theorem}[目标成功率下最优辅助比特预算的严格分段阈值]\label{thm:xi-time-part70-target-success-optimal-budget-threshold}
沿用定理 \ref{thm:xi-foldbin-dyadic-capacity-critical-window-limit} 的记号。对任意目标成功率 \(\sigma\in(0,1)\)，定义最小预算
\[
B_m^\star(\sigma):=
\min\bigl\{B\in\mathbb N:\ \mathrm{Succ}_m^{\mathrm{bin}}(B)\ge \sigma\bigr\}.
\]
再定义连续阈值常数
\[
s_\sigma:=
\begin{cases}
\dfrac{\sqrt5}{\varphi^2}\,\sigma,
& 0<\sigma\le \dfrac{\varphi}{\sqrt5},\\[8pt]
\dfrac{\sqrt5\,\sigma-\varphi^{-1}}{\varphi},
& \dfrac{\varphi}{\sqrt5}\le \sigma<1.
\end{cases}
\]
则有
\[
B_m^\star(\sigma)
=
m\log_2\!\left(\frac{2}{\varphi}\right)+\log_2 s_\sigma+O(1).
\]
更一般地，若整数序列 \((B_m)_{m\ge 1}\) 满足
\[
2^{B_m}\left(\frac{\varphi}{2}\right)^m\longrightarrow s\in(0,\infty),
\]
则
\[
\mathrm{Succ}_m^{\mathrm{bin}}(B_m)\longrightarrow L(s),
\]
其中 \(L\) 即定理 \ref{thm:xi-time-part70-continuous-capacity-two-atom-second-derivative} 的极限曲线；因此 \(s_\sigma\) 正是方程 \(L(s)=\sigma\) 的唯一解。
\end{theorem}
\begin{proof}
第二个结论正是定理 \ref{thm:xi-foldbin-dyadic-capacity-critical-window-limit} 与推论 \ref{cor:xi-foldbin-critical-scale-success-curve-inversion} 的内容，其中极限成功曲线即
\[
L(s)=
\begin{cases}
\dfrac{\varphi^2}{\sqrt5}\,s, & 0\le s\le \varphi^{-1},\\[6pt]
\dfrac{\varphi}{\sqrt5}\,s+\dfrac{1}{\varphi\sqrt5}, & \varphi^{-1}\le s\le 1,\\[8pt]
1, & s\ge 1.
\end{cases}
\]
因此 \(s_\sigma=L^{-1}(\sigma)\) 的显式公式恰为上式两段线性方程的解。

现证明 \(B_m^\star(\sigma)\) 的渐近公式。由于 \(\sigma\in(0,1)\)，故 \(s_\sigma\in(0,1)\)。取
\[
0<\varepsilon<\min\{s_\sigma,1-s_\sigma\}.
\]
由 \(L\) 在 \([0,1]\) 上严格递增可知
\[
L(s_\sigma-\varepsilon)<\sigma<L(s_\sigma+\varepsilon).
\]
再由 critical-scale 极限，分别取
\[
\widetilde B_m^-:=\left\lfloor \log_2\!\Bigl((s_\sigma-\varepsilon)\Bigl(\frac{2}{\varphi}\Bigr)^m\Bigr)\right\rfloor,
\qquad
\widetilde B_m^+:=\left\lceil \log_2\!\Bigl((s_\sigma+\varepsilon)\Bigl(\frac{2}{\varphi}\Bigr)^m\Bigr)\right\rceil,
\]
则当 \(m\) 充分大时有
\[
\mathrm{Succ}_m^{\mathrm{bin}}(\widetilde B_m^-)<\sigma,
\qquad
\mathrm{Succ}_m^{\mathrm{bin}}(\widetilde B_m^+)\ge \sigma.
\]
由于 \(\mathrm{Succ}_m^{\mathrm{bin}}(B)\) 关于 \(B\) 单调不减，而 \(B_m^\star(\sigma)\) 按定义是达到成功率阈值的最小整数预算，遂得
\[
\widetilde B_m^-+1\le B_m^\star(\sigma)\le \widetilde B_m^+.
\]
把上式展开后便得到
\[
\log_2\!\Bigl((s_\sigma-\varepsilon)\Bigl(\frac{2}{\varphi}\Bigr)^m\Bigr)
\le
B_m^\star(\sigma)
\le
\log_2\!\Bigl((s_\sigma+\varepsilon)\Bigl(\frac{2}{\varphi}\Bigr)^m\Bigr)+1.
\]
由于 \(\varepsilon\) 固定，端点与
\[
m\log_2\!\left(\frac{2}{\varphi}\right)+\log_2 s_\sigma
\]
之差都是有界量，故结论成立。证毕。
\end{proof}

\begin{theorem}[由 \texorpdfstring{$(\log C_m)$}{(log Cm)} 序列恢复 Bernoulli 塔时的 Richardson 残差比极限]\label{thm:xi-time-part70-logcm-richardson-residual-ratio}
记
\[
q:=\frac{\varphi}{2}\in(0,1),
\qquad
E_m:=\log C_m-\log(2\pi),
\]
\[
a_r:=
\frac{B_{2r}}{r(2r-1)}
\bigl(\varphi^{-1}+\varphi^{2r-3}\bigr)
\qquad (r\ge 1).
\]
对每个固定整数 \(R\ge 0\)，定义 \(R\) 阶 Richardson 残差
\[
\mathcal R_m^{(R)}
:=
E_m-\sum_{r=1}^{R}a_r q^{(2r-1)m}.
\]
则有严格比值极限
\[
\frac{\mathcal R_{m+1}^{(R)}}{\mathcal R_m^{(R)}}
\longrightarrow
q^{2R+1}
=
\left(\frac{\varphi}{2}\right)^{2R+1}.
\]
换言之，一旦减去前 \(R\) 层 Bernoulli 模态，剩余误差的相邻比值必然锁定到奇数次黄金幂 \((\varphi/2)^{2R+1}\)。
\end{theorem}
\begin{proof}
对任意固定 \(R\ge 0\)，将定理 \ref{thm:fold-bin-gauge-constant-stirling-bernoulli-hierarchy} 应用于 \(R+1\) 阶，得到
\[
E_m
=
\sum_{r=1}^{R+1}a_r q^{(2r-1)m}
+O\!\bigl(q^{(2R+3)m}\bigr).
\]
因此
\[
\mathcal R_m^{(R)}
=
a_{R+1}q^{(2R+1)m}
+O\!\bigl(q^{(2R+3)m}\bigr).
\]
又因为偶 Bernoulli 数满足 \(B_{2r}\neq 0\)（\(r\ge 1\)），而 \(\varphi^{-1}+\varphi^{2r-3}>0\)，故
\[
a_{R+1}\neq 0.
\]
于是
\[
\mathcal R_m^{(R)}
=
a_{R+1}q^{(2R+1)m}\bigl(1+O(q^{2m})\bigr),
\]
从而
\[
\frac{\mathcal R_{m+1}^{(R)}}{\mathcal R_m^{(R)}}
=
\frac{a_{R+1}q^{(2R+1)(m+1)}\bigl(1+O(q^{2m})\bigr)}
{a_{R+1}q^{(2R+1)m}\bigl(1+O(q^{2m})\bigr)}
\longrightarrow
q^{2R+1}.
\]
证毕。
\end{proof}

\begin{theorem}[奇素数模 RH 缺陷总质量的可和性]\label{thm:xi-time-part70-odd-prime-rh-defect-summable-mass}
记
\[
\lambda:=\varphi^2,
\qquad
\rho_{p,1}:=\rho\!\bigl(M(e^{2\pi i/p})\bigr)
\qquad
(p\ge 3\ \text{为奇素数}),
\]
并定义扭曲指数与其缺陷
\[
\theta_{p,1}:=\frac{\log \rho_{p,1}}{\log\lambda},
\qquad
\delta_p:=1-\theta_{p,1}.
\]
则有显式展开
\[
\delta_p
=
\frac{4\pi^2(6\sqrt5-5)}{250\log(\varphi^2)}\,\frac{1}{p^2}
+O\!\left(\frac{1}{p^4}\right).
\]
从而
\[
\sum_{\substack{p\ge 3\\ p\ \text{为奇素数}}}\delta_p<\infty,
\qquad
\sum_{\substack{p>P\\ p\ \text{为奇素数}}}\delta_p=O(P^{-1}).
\]
因此，尽管奇素数方向上的平方根门槛失效在频率上呈余有限出现，其缺陷幅度却形成可和尾谱。
\end{theorem}
\begin{proof}
由定理 \ref{thm:xi-output-potential-odd-prime-rh-obstruction-index-to-one}，
\[
\theta_{p,1}
=
1-\frac{6\sqrt5-5}{250\log\lambda}\Bigl(\frac{2\pi}{p}\Bigr)^2
+\frac{7+24\sqrt5}{75000\log\lambda}\Bigl(\frac{2\pi}{p}\Bigr)^4
+O(p^{-6}).
\]
将 \(\lambda=\varphi^2\) 代入并移项，便得到
\[
\delta_p
=
\frac{4\pi^2(6\sqrt5-5)}{250\log(\varphi^2)}\,\frac{1}{p^2}
+O(p^{-4}).
\]
特别地，存在常数 \(c_1,c_2>0\) 与奇素数阈值 \(p_0\)，使得对所有奇素数 \(p\ge p_0\) 都有
\[
\frac{c_1}{p^2}\le \delta_p\le \frac{c_2}{p^2}.
\]
于是
\[
\sum_{\substack{p\ge p_0\\ p\ \text{为奇素数}}}\delta_p
\ll
\sum_{n\ge p_0}\frac{1}{n^2}
<\infty,
\]
从而总和收敛。对尾和，同样有
\[
\sum_{\substack{p>P\\ p\ \text{为奇素数}}}\delta_p
\ll
\sum_{n>P}\frac{1}{n^2}
=
O(P^{-1}).
\]
证毕。
\end{proof}

\begin{corollary}[奇素数 RH 缺陷尾谱的 \texorpdfstring{$\ell^\alpha$}{ell alpha} 临界指数恰为 \texorpdfstring{$1/2$}{1/2}]\label{cor:xi-time-part70-odd-prime-rh-defect-ell-alpha-threshold}
沿用定理 \ref{thm:xi-time-part70-odd-prime-rh-defect-summable-mass} 的记号。则对任意 \(\alpha>0\)，有
\[
(\delta_p)_{p\ \text{为奇素数}}\in \ell^\alpha
\qquad\Longleftrightarrow\qquad
\alpha>\frac12.
\]
\end{corollary}
\begin{proof}
由定理 \ref{thm:xi-time-part70-odd-prime-rh-defect-summable-mass}，存在常数 \(c_1,c_2>0\) 与 \(p_0\) 使得对全部奇素数 \(p\ge p_0\) 都有
\[
c_1p^{-2}\le \delta_p\le c_2p^{-2}.
\]
于是
\[
\delta_p^\alpha\asymp p^{-2\alpha}
\qquad (p\to\infty,\ p\ \text{为奇素数}).
\]
而素数级数满足
\[
\sum_{p}p^{-s}<\infty
\qquad\Longleftrightarrow\qquad
s>1.
\]
因此
\[
\sum_{\substack{p\ge 3\\ p\ \text{为奇素数}}}\delta_p^\alpha<\infty
\qquad\Longleftrightarrow\qquad
\sum_{\substack{p\ge 3\\ p\ \text{为奇素数}}}p^{-2\alpha}<\infty
\qquad\Longleftrightarrow\qquad
2\alpha>1,
\]
即所求结论。证毕。
\end{proof}

\begin{theorem}[奇素数最小非平凡单位根扭曲的谱半径二次接近律与最终严格单调逼近]\label{thm:xi-time-part70-prime-twist-spectral-radius-quadratic-approach}
沿用推论 \ref{cor:real-input-40-collision-twist-fourth} 的记号，并记
\[
\lambda:=\rho(1)=\varphi^2,
\qquad
\rho_p:=\rho\!\left(e^{2\pi i/p}\right)
\qquad
(p\ge 3\ \text{为奇素数}).
\]
则成立：
\begin{enumerate}
  \item 有精确二次接近律
  \[
  \boxed{\;
  p^2\left(1-\frac{\rho_p}{\lambda}\right)
  \longrightarrow
  \frac{6\sqrt5-5}{250}\,(2\pi)^2.\;}
  \]
  \item 存在奇素数阈值 \(p_0\)，使得对任意奇素数 \(p<q\) 且 \(p,q\ge p_0\)，都有
  \[
  \boxed{\;
  \rho_p<\rho_q<\lambda.\;}
  \]
\end{enumerate}
换言之，最小非平凡 prime twist 的谱半径不仅以严格 \(p^{-2}\) 主项从下方逼近 \(\lambda\)，而且在充分大素数区间上按模数严格单调恢复。
\end{theorem}
\begin{proof}
由推论 \ref{cor:real-input-40-collision-twist-fourth}，当 \(t\to 0\) 时有
\[
\log\frac{\rho(t)}{\lambda}
=
-A t^2+B t^4+O(t^6),
\]
其中
\[
A:=\frac{6\sqrt5-5}{250}>0,
\qquad
B:=\frac{7+24\sqrt5}{75000}.
\]
将其指数化，得到
\[
\frac{\rho(t)}{\lambda}
=
\exp\!\bigl(-A t^2+B t^4+O(t^6)\bigr)
=
1-A t^2+O(t^4).
\]
对 \(t=t_p:=2\pi/p\) 代入，即得
\[
1-\frac{\rho_p}{\lambda}
=
A\left(\frac{2\pi}{p}\right)^2+O(p^{-4}),
\]
从而
\[
p^2\left(1-\frac{\rho_p}{\lambda}\right)
\longrightarrow
A(2\pi)^2
=
\frac{6\sqrt5-5}{250}\,(2\pi)^2,
\]
这就证明了第一条。

再令
\[
R(t):=\frac{\rho(t)}{\lambda}.
\]
由上式可得
\[
\log R(t)=-A t^2+B t^4+O(t^6).
\]
对其求导，得到
\[
\frac{R'(t)}{R(t)}
=
-2At+4Bt^3+O(t^5)
=
-2At\bigl(1+O(t^2)\bigr).
\]
因此存在 \(t_0>0\)，使得当 \(0<t\le t_0\) 时，
\[
\frac{R'(t)}{R(t)}<0.
\]
又因 \(R(t)>0\)，可知 \(R'(t)<0\)；故 \(R\) 在 \((0,t_0]\) 上严格递减。另一方面，由
\[
R(t)=1-A t^2+O(t^4)
\]
可知在同一小邻域内 \(R(t)<1\)。

现取奇素数阈值 \(p_0\) 使得
\[
\frac{2\pi}{p_0}\le t_0.
\]
若 \(p<q\) 且 \(p,q\ge p_0\)，则
\[
t_p=\frac{2\pi}{p}>\frac{2\pi}{q}=t_q,
\]
而 \(R\) 在 \((0,t_0]\) 上严格递减，遂得
\[
R(t_p)<R(t_q)<1.
\]
两边同乘以 \(\lambda\)，即
\[
\rho_p<\rho_q<\lambda.
\]
第二条得证。证毕。
\end{proof}

\noindent
上述七条结果表明，二原子塌缩并不只停留在 bin-fold 的输出实验或连续容量曲线上，而会同时外显为稳定层均匀经验谱的严格弱极限、目标成功率反演中的唯一分段阈值、\((\log C_m)\) 的奇次黄金幂残差比，以及奇素数单位根扭曲缺陷的可和尾谱与谱半径恢复律。因而，\(\{\varphi^{-1},1\}\) 这两个黄金比例尺度既支配欠分辨容量几何，也支配 Bernoulli 递归剥离后的剩余模态；而在奇素数单位根扭曲方向上，坏模数的频率与缺陷的振幅不仅被严格分离为“余有限出现”与“总质量可和”两种不同层级，而且最小非平凡扭曲的谱半径还会按严格 \(p^{-2}\) 主项最终单调地从下方逼近 \(\lambda\)。于是，稳定层二点律、资源阈值反演、Richardson 诊断与奇素数 RH 缺陷最终被压缩进同一条黄金尺度主线。

\paragraph{稳定层二态统计极限与规范体积修正的一阶账本}
在稳定层均匀经验谱的两点弱极限、统一层矩函数以及 fold-gauge 常数恢复链已经建立之后，还可把二态一致渐近对均匀层平均量与规范体积修正的支配压缩为下列四条直接后果。它们分别给出缩放退化度的量化两点律、全部缩放矩与幂和首项常数、均匀平均对数退化度的二态展开，以及 \(g_m-\kappa_m+1\) 的精确首项。

\begin{theorem}[均匀稳定类型下缩放退化度的量化两点极限律]\label{thm:xi-time-part70a-uniform-scaled-degeneracy-quantitative}
令
\[
u_m(w):=\frac{1}{|X_m|}\qquad (w\in X_m),
\]
并取 \(W_m\sim u_m\)。定义缩放退化度随机变量
\[
Y_m:=\left(\frac{\varphi}{2}\right)^m d_m^{\mathrm{bin}}(W_m).
\]
则对任意有界 Lipschitz 函数 \(f:\RR\to\RR\)，都有
\[
\mathbb E[f(Y_m)]
=
\frac{F_{m+1}}{F_{m+2}}\,f(1)
+
\frac{F_m}{F_{m+2}}\,f(\varphi^{-1})
+O\!\left(\left(\frac{\varphi}{2}\right)^m\right),
\]
其中隐含常数仅依赖于 \(f\) 的 Lipschitz 常数。特别地，\(Y_m\) 依分布收敛到两点随机变量
\[
Y_\infty=
\begin{cases}
1,&\text{概率 } \varphi^{-1},\\[3pt]
\varphi^{-1},&\text{概率 } \varphi^{-2}.
\end{cases}
\]
\end{theorem}
\begin{proof}
由定理 \ref{thm:fold-bin-two-state-asymptotic}，存在常数 \(C>0\)，使得对充分大的 \(m\) 与全部 \(w\in X_m\) 都一致有
\[
\left|
\left(\frac{\varphi}{2}\right)^m d_m^{\mathrm{bin}}(w)-\varphi^{-w_m}
\right|
\le C\left(\frac{\varphi}{2}\right)^m.
\]
因此
\[
\bigl|Y_m-\varphi^{-W_m(m)}\bigr|
\le C\left(\frac{\varphi}{2}\right)^m
\]
几乎处处成立。若记 \(L_f\) 为 \(f\) 的 Lipschitz 常数，则
\[
\bigl|f(Y_m)-f(\varphi^{-W_m(m)})\bigr|
\le
L_f C\left(\frac{\varphi}{2}\right)^m.
\]
对均匀分布取期望，得到
\[
\mathbb E[f(Y_m)]
=
\mathbb E[f(\varphi^{-W_m(m)})]
+O\!\left(\left(\frac{\varphi}{2}\right)^m\right).
\]
另一方面，
\[
\mathbb E[f(\varphi^{-W_m(m)})]
=
\frac{F_{m+1}}{F_{m+2}}\,f(1)
+
\frac{F_m}{F_{m+2}}\,f(\varphi^{-1}),
\]
因为
\[
\mathbb P(W_m(m)=0)=\frac{F_{m+1}}{F_{m+2}},
\qquad
\mathbb P(W_m(m)=1)=\frac{F_m}{F_{m+2}}.
\]
这就得到所示定量公式。再用
\[
\frac{F_{m+1}}{F_{m+2}}\to\varphi^{-1},
\qquad
\frac{F_m}{F_{m+2}}\to\varphi^{-2},
\]
即可得到所述两点极限分布。证毕。
\end{proof}

\begin{corollary}[统一缩放矩极限与 \texorpdfstring{$S_q^{\mathrm{bin}}(m)$}{Sqbin(m)} 的首项常数]\label{cor:xi-time-part70a-scaled-moments-leading-constant}
对任意实数 \(q\ge 0\)，都有
\[
\lim_{m\to\infty}
\frac{1}{|X_m|}
\sum_{w\in X_m}
\left(
\left(\frac{\varphi}{2}\right)^m d_m^{\mathrm{bin}}(w)
\right)^q
=
\frac{1}{\varphi}+\varphi^{-q-2}.
\]
等价地，若记
\[
S_q^{\mathrm{bin}}(m):=\sum_{w\in X_m}\bigl(d_m^{\mathrm{bin}}(w)\bigr)^q,
\]
则有精确首项渐近
\[
\boxed{\;
S_q^{\mathrm{bin}}(m)
\sim
\frac{\varphi+\varphi^{-q}}{\sqrt5}
\left(\frac{2^q}{\varphi^{q-1}}\right)^m.\;}
\]
\end{corollary}
\begin{proof}
对固定 \(q\ge 0\)，由上一条定理作用于函数 \(f_q(y):=y^q\) 即可。事实上，由定理 \ref{thm:xi-time-part70a-uniform-scaled-degeneracy-quantitative} 的证明，对充分大的 \(m\) 与几乎处处的样本点都有
\[
Y_m\in\left[\frac{\varphi^{-1}}{2},2\right],
\]
故 \(f_q\) 在该固定紧区间上 Lipschitz。于是
\[
\frac{1}{|X_m|}
\sum_{w\in X_m}
\left(
\left(\frac{\varphi}{2}\right)^m d_m^{\mathrm{bin}}(w)
\right)^q
\to
\frac{1}{\varphi}+\varphi^{-q-2}.
\]
将上式改写为
\[
S_q^{\mathrm{bin}}(m)
=
|X_m|
\left(\frac{2}{\varphi}\right)^{mq}
\left(\frac{1}{\varphi}+\varphi^{-q-2}+o(1)\right),
\]
再用 Binet 公式
\[
|X_m|=F_{m+2}\sim \frac{\varphi^{m+2}}{\sqrt5},
\]
便得到
\[
S_q^{\mathrm{bin}}(m)
\sim
\frac{\varphi^{m+2}}{\sqrt5}
\left(\frac{2}{\varphi}\right)^{mq}
\left(\varphi^{-1}+\varphi^{-q-2}\right)
=
\frac{\varphi+\varphi^{-q}}{\sqrt5}
\left(\frac{2^q}{\varphi^{q-1}}\right)^m.
\]
证毕。
\end{proof}

\begin{theorem}[均匀稳定层平均对数退化度的二态展开]\label{thm:xi-time-part70a-uniform-average-logdeg-two-state}
定义均匀稳定层平均对数退化度
\[
\overline{\ell}_m
:=
\frac{1}{|X_m|}\sum_{w\in X_m}\log d_m^{\mathrm{bin}}(w).
\]
则有
\[
\boxed{\;
\overline{\ell}_m
=
m\log\!\Bigl(\frac{2}{\varphi}\Bigr)
-\frac{F_m}{F_{m+2}}\log\varphi
+O\!\left(\left(\frac{\varphi}{2}\right)^m\right).\;}
\]
因而
\[
\boxed{\;
\overline{\ell}_m
=
m\log\!\Bigl(\frac{2}{\varphi}\Bigr)
-\varphi^{-2}\log\varphi
+O\!\left(\left(\frac{\varphi}{2}\right)^m\right).\;}
\]
\end{theorem}
\begin{proof}
由定理 \ref{thm:fold-bin-two-state-asymptotic} 的一致对数展开，
\[
\log d_m^{\mathrm{bin}}(w)
=
m\log\!\Bigl(\frac{2}{\varphi}\Bigr)
-w_m\log\varphi
+O\!\left(\left(\frac{\varphi}{2}\right)^m\right)
\qquad (w\in X_m),
\]
且误差在 \(w\) 上一致。对全部 \(w\in X_m\) 求平均，得到
\[
\overline{\ell}_m
=
m\log\!\Bigl(\frac{2}{\varphi}\Bigr)
-\frac{1}{|X_m|}\sum_{w\in X_m}w_m\,\log\varphi
+O\!\left(\left(\frac{\varphi}{2}\right)^m\right).
\]
而
\[
\sum_{w\in X_m}w_m=|A_1|=F_m,
\qquad
|X_m|=F_{m+2},
\]
故
\[
\frac{1}{|X_m|}\sum_{w\in X_m}w_m=\frac{F_m}{F_{m+2}}.
\]
代回即得第一式。再用 Binet 公式给出的比值极限
\[
\frac{F_m}{F_{m+2}}=\varphi^{-2}+O(\varphi^{-2m}),
\]
便得到第二式。证毕。
\end{proof}
\leanverified{paper\_xi\_time\_part70a\_uniform\_average\_logdeg\_two\_state}

\leanverified{paper\_xi\_time\_part70a\_gauge\_fiber\_gap\_first\_order}
\begin{theorem}[规范体积与纤维信息差的一阶首项]\label{thm:xi-time-part70a-gauge-fiber-gap-first-order}
定义
\[
\mu_m(w):=\frac{d_m^{\mathrm{bin}}(w)}{2^m},
\qquad
\kappa_m:=\sum_{w\in X_m}\mu_m(w)\log d_m^{\mathrm{bin}}(w),
\qquad
g_m:=2^{-m}\log|G_m^{\mathrm{bin}}|,
\]
并沿用推论 \ref{cor:fold-bin-recover-2pi} 的规范常数
\[
C_m
:=
\exp\!\left(
\frac{2^{m+1}}{|X_m|}\bigl(g_m-\kappa_m+1\bigr)-\overline{\ell}_m
\right).
\]
则有
\[
\boxed{\;
\frac{2^{m+1}}{|X_m|}\bigl(g_m-\kappa_m+1\bigr)
=
m\log\!\Bigl(\frac{2}{\varphi}\Bigr)
-\frac{F_m}{F_{m+2}}\log\varphi
+\log(2\pi)
+O\!\left(\left(\frac{\varphi}{2}\right)^m\right).\;}
\]
因此
\[
\boxed{\;
g_m-\kappa_m+1
=
\frac{\varphi^2}{2\sqrt5}
\Bigl(
m\log\!\Bigl(\frac{2}{\varphi}\Bigr)
-\varphi^{-2}\log\varphi
+\log(2\pi)
\Bigr)
\left(\frac{\varphi}{2}\right)^m
+O\!\left(\left(\frac{\varphi}{2}\right)^{2m}\right).\;}
\]
\end{theorem}
\begin{proof}
由 \(C_m\) 的定义可得精确恒等式
\[
\frac{2^{m+1}}{|X_m|}\bigl(g_m-\kappa_m+1\bigr)=\log C_m+\overline{\ell}_m.
\]
另一方面，推论 \ref{cor:fold-bin-recover-2pi} 给出
\[
\log C_m=\log(2\pi)+O\!\left(\left(\frac{\varphi}{2}\right)^m\right),
\]
而定理 \ref{thm:xi-time-part70a-uniform-average-logdeg-two-state} 已证明
\[
\overline{\ell}_m
=
m\log\!\Bigl(\frac{2}{\varphi}\Bigr)
-\frac{F_m}{F_{m+2}}\log\varphi
+O\!\left(\left(\frac{\varphi}{2}\right)^m\right).
\]
相加便得到第一条公式。

再由 Binet 公式，
\[
\frac{|X_m|}{2^{m+1}}
=
\frac{F_{m+2}}{2^{m+1}}
=
\frac{\varphi^2}{2\sqrt5}\left(\frac{\varphi}{2}\right)^m
+O\!\bigl((2\varphi)^{-m}\bigr),
\]
且
\[
\frac{F_m}{F_{m+2}}=\varphi^{-2}+O(\varphi^{-2m}).
\]
用这两式乘第一条公式，得到
\[
g_m-\kappa_m+1
=
\frac{\varphi^2}{2\sqrt5}
\Bigl(
m\log\!\Bigl(\frac{2}{\varphi}\Bigr)
-\varphi^{-2}\log\varphi
+\log(2\pi)
\Bigr)
\left(\frac{\varphi}{2}\right)^m
+O\!\bigl(m(2\varphi)^{-m}\bigr)
+O\!\left(\left(\frac{\varphi}{2}\right)^{2m}\right).
\]
由于
\[
\frac{1}{2\varphi}<\left(\frac{\varphi}{2}\right)^2,
\]
故
\[
m(2\varphi)^{-m}=O\!\left(\left(\frac{\varphi}{2}\right)^{2m}\right).
\]
于是第二条公式成立。证毕。
\end{proof}

\noindent
由此，bin-fold 的两态一致渐近不但控制稳定层均匀经验谱与全部缩放矩，而且把均匀平均的对数退化度与规范体积修正的首项常数同时固定在同一末位二态账本之上。于是，均匀层二点律、幂和首项、平均对数势与规范体积差额最终被收束为同一条可显式求值的黄金比例主线。

\endinput


\paragraph{bin-fold 退化度的全实矩闭式与 Bernoulli 塔的逆矩驱动}
在稳定层均匀经验谱的两点极限、全部非负缩放矩的首项常数以及 fold-gauge 常数的 Stirling--Bernoulli 层级已经建立之后，还可把 bin-fold 退化度的实参数矩谱与 Bernoulli 模态系数进一步压缩为下列两条直接接口。

\begin{theorem}[bin-fold 退化度的全实矩闭式]\label{thm:xi-time-part70aa-binfold-all-real-moments-closed}
对每个 \(m\ge 1\)，记
\[
A_0^{(m)}:=\{w\in X_m:\ w_m=0\},
\qquad
A_1^{(m)}:=\{w\in X_m:\ w_m=1\},
\]
并定义稳定层均匀实矩
\[
M_s^{\mathrm{bin}}(m)
:=
\frac{1}{|X_m|}
\sum_{w\in X_m}\bigl(d_m^{\mathrm{bin}}(w)\bigr)^s
\qquad (s\in\RR).
\]
则对任意固定实数 \(s\)，都有渐近公式
\[
\boxed{\;
M_s^{\mathrm{bin}}(m)
=
\left(\frac{2}{\varphi}\right)^{sm}
\left(
\varphi^{-1}+\varphi^{-s-2}
\right)
+
O_s\!\left(
\left(\frac{2}{\varphi}\right)^{sm}
\left(\frac{\varphi}{2}\right)^m
\right).\;}
\]
等价地，
\[
\boxed{\;
\lim_{m\to\infty}
\left(\frac{\varphi}{2}\right)^{sm}
M_s^{\mathrm{bin}}(m)
=
\varphi^{-1}+\varphi^{-s-2}.\;}
\]
\leanverified{paper\_xi\_time\_part70aa\_binfold\_all\_real\_moments\_closed}
\end{theorem}
\begin{proof}
记
\[
\varepsilon_m:=\left(\frac{\varphi}{2}\right)^m.
\]
由定理 \ref{thm:fold-bin-two-state-asymptotic}，存在常数 \(C>0\)，使对充分大的 \(m\) 与全部 \(w\in X_m\) 都有
\[
\left(\frac{\varphi}{2}\right)^m d_m^{\mathrm{bin}}(w)
=
\begin{cases}
1+O(\varepsilon_m), & w\in A_0^{(m)},\\[4pt]
\varphi^{-1}+O(\varepsilon_m), & w\in A_1^{(m)},
\end{cases}
\]
且误差在 \(w\) 上一致。由于两块主值 \(1\) 与 \(\varphi^{-1}\) 都严格远离 \(0\)，对任意固定 \(s\in\RR\)，幂函数 \(x\mapsto x^s\) 在这两个主值附近都可作一致一阶展开，从而
\[
\bigl(d_m^{\mathrm{bin}}(w)\bigr)^s
=
\left(\frac{2}{\varphi}\right)^{sm}
\begin{cases}
1+O_s(\varepsilon_m), & w\in A_0^{(m)},\\[4pt]
\varphi^{-s}+O_s(\varepsilon_m), & w\in A_1^{(m)},
\end{cases}
\]
其中隐含常数仅依赖于 \(s\)。

对两块分别求和并除以 \(|X_m|=F_{m+2}\)，得到
\[
M_s^{\mathrm{bin}}(m)
=
\left(\frac{2}{\varphi}\right)^{sm}
\left(
\frac{F_{m+1}+\varphi^{-s}F_m}{F_{m+2}}
+
O_s(\varepsilon_m)
\right),
\]
因为
\[
|A_0^{(m)}|=F_{m+1},
\qquad
|A_1^{(m)}|=F_m.
\]
再用 Binet 公式给出的比值渐近
\[
\frac{F_{m+1}}{F_{m+2}}=\varphi^{-1}+O\!\left(\varphi^{-2m}\right),
\qquad
\frac{F_m}{F_{m+2}}=\varphi^{-2}+O\!\left(\varphi^{-2m}\right),
\]
并注意 \(\varphi^{-2m}=O(\varepsilon_m)\)，便得到
\[
M_s^{\mathrm{bin}}(m)
=
\left(\frac{2}{\varphi}\right)^{sm}
\left(
\varphi^{-1}+\varphi^{-s-2}
+
O_s(\varepsilon_m)
\right).
\]
这就给出所示渐近式；第二条极限式只是将主项重新整理后的等价写法。证毕。
\end{proof}

\begin{corollary}[Bernoulli 级数系数的逆矩解释]\label{cor:xi-time-part70aa-bernoulli-coefficient-inverse-moment-interpretation}
对每个整数 \(r\ge 1\)，定义奇数阶逆矩
\[
I_{2r-1}(m)
:=
\frac{1}{|X_m|}
\sum_{w\in X_m}\bigl(d_m^{\mathrm{bin}}(w)\bigr)^{-(2r-1)}.
\]
则有
\[
\boxed{\;
I_{2r-1}(m)
=
\left(\frac{\varphi}{2}\right)^{(2r-1)m}
\left(
\varphi^{-1}+\varphi^{2r-3}
\right)
+
O_r\!\left(
\left(\frac{\varphi}{2}\right)^{2rm}
\right).\;}
\]
从而
\[
\boxed{\;
\lim_{m\to\infty}
\left(\frac{2}{\varphi}\right)^{(2r-1)m}
I_{2r-1}(m)
=
\varphi^{-1}+\varphi^{2r-3}.\;}
\]
另一方面，对任意固定 \(R\ge 1\)，定理 \ref{thm:fold-bin-gauge-constant-stirling-bernoulli-hierarchy} 给出
\[
\log C_m
=
\log(2\pi)
+
\sum_{r=1}^{R}
\frac{B_{2r}}{r(2r-1)}
\bigl(\varphi^{-1}+\varphi^{2r-3}\bigr)
\left(\frac{\varphi}{2}\right)^{(2r-1)m}
+
O\!\left(\left(\frac{\varphi}{2}\right)^{(2R+1)m}\right).
\]
因此，第 \(r\) 级 Bernoulli 模态系数恰可写为
\[
\frac{B_{2r}}{r(2r-1)}
\cdot
\lim_{m\to\infty}
\left(\frac{2}{\varphi}\right)^{(2r-1)m}
I_{2r-1}(m).
\]
\end{corollary}
\begin{proof}
在定理 \ref{thm:xi-time-part70aa-binfold-all-real-moments-closed} 中取
\[
s=-(2r-1)
\]
即可得到
\[
I_{2r-1}(m)=M_{-(2r-1)}^{\mathrm{bin}}(m),
\]
并由
\[
\varphi^{-(-2r+1)-2}=\varphi^{2r-3}
\]
得到所示主项常数。误差项则满足
\[
\left(\frac{2}{\varphi}\right)^{-(2r-1)m}
\left(\frac{\varphi}{2}\right)^m
=
\left(\frac{\varphi}{2}\right)^{2rm}.
\]
最后，把该极限常数与定理 \ref{thm:fold-bin-gauge-constant-stirling-bernoulli-hierarchy} 中对应的第 \(r\) 级系数逐项比较，即得最后一式。证毕。
\end{proof}

\noindent
这表明，fold-gauge 常数序列中的 Bernoulli 塔并不是脱离 bin-fold 退化度谱而附加的解析修正；相反，它正由稳定层退化度的奇数阶逆矩常数逐级驱动。

\endinput


\paragraph{末位二态分层、逆阈值几何与 Bernoulli 塔的 Mellin 编码}
在稳定层均匀经验谱的两点极限、escort 温度族的末位 Bernoulli 塌缩、中心临界容量骨架的双折点闭式以及实参数矩谱与 Bernoulli 模态系数之间的桥接已经建立之后，还可把阈值/采样读出、路径/相位重加权与 Stirling--Bernoulli 层级进一步压缩为同一条末位二态主线。下列结果分别刻画最终严格分层、分位二原子塌缩、归一化逆阈值几何、escort 熵谱的完整常数项，以及容量曲线对 Bernoulli 塔几何权重的 Mellin 编码。

\begin{theorem}[末位相位的最终严格分层]\label{thm:xi-time-part70ab-terminal-bit-eventual-strict-stratification}
\leanverified{paper\_xi\_time\_part70ab\_terminal\_bit\_eventual\_strict\_stratification}
沿用记号
\[
d_m(w):=d_m^{\mathrm{bin}}(w),\qquad
A_i^{(m)}:=\{w\in X_m:\ w_m=i\}\quad(i=0,1),
\qquad
A_m:=\left(\frac{2}{\varphi}\right)^m.
\]
则存在整数 \(m_0\)，使得对一切 \(m\ge m_0\)，都有严格分离
\[
\min_{w\in A_0^{(m)}} d_m(w)
>
\max_{v\in A_1^{(m)}} d_m(v).
\]
因此在充分高分辨率下，纤维多重度的全序被末位 \(w_m\) 强制划分为两层。
\end{theorem}
\begin{proof}
由定理 \ref{thm:fold-bin-two-state-asymptotic}，存在绝对常数 \(C>0\)，使对充分大的 \(m\) 与全部 \(x\in X_m\) 都有
\[
\bigl|d_m(x)-\varphi^{-x_m}A_m\bigr|\le C.
\]
若 \(w\in A_0^{(m)}\) 且 \(v\in A_1^{(m)}\)，则
\[
d_m(w)-d_m(v)
\ge
A_m-C-\bigl(\varphi^{-1}A_m+C\bigr)
=
\varphi^{-2}A_m-2C.
\]
由于 \(A_m\to\infty\)，右端对充分大的 \(m\) 为正，结论成立。
\end{proof}

\begin{corollary}[极值纤维的末位定位与最小尺度失配率]\label{cor:xi-time-part70ab-extremal-fibers-terminal-bit-localization}
\leanverified{paper\_xi\_time\_part70ab\_extremal\_fibers\_terminal\_bit\_localization}
定义
\[
d_{\min}(m):=\min_{w\in X_m}d_m(w),
\qquad
M_m:=\max_{w\in X_m}d_m(w).
\]
则对充分大的 \(m\)，有
\[
\arg\min_{w\in X_m} d_m(w)\subseteq A_1^{(m)},
\qquad
\arg\max_{w\in X_m} d_m(w)\subseteq A_0^{(m)}.
\]
并且
\[
d_{\min}(m)=\varphi^{-1}A_m+O(1),
\qquad
M_m=A_m+O(1),
\qquad
\frac{M_m}{d_{\min}(m)}\longrightarrow \varphi.
\]
此外，沿偶数分辨率 \(m\to\infty\) 有
\[
\frac{d_{\min}(m)}{F_{m/2}}\longrightarrow 0.
\]
\end{corollary}
\begin{proof}
由定理 \ref{thm:xi-time-part70ab-terminal-bit-eventual-strict-stratification}，对充分大的 \(m\)，最小值只能取自 \(A_1^{(m)}\)，最大值只能取自 \(A_0^{(m)}\)。另一方面，由定理 \ref{thm:fold-bin-two-state-asymptotic} 的一致渐近，
\[
d_m(w)=\varphi^{-1}A_m+O(1)\quad(w\in A_1^{(m)}),
\qquad
d_m(w)=A_m+O(1)\quad(w\in A_0^{(m)}),
\]
且误差在各块上统一，故
\[
d_{\min}(m)=\varphi^{-1}A_m+O(1),
\qquad
M_m=A_m+O(1).
\]
于是 \(M_m/d_{\min}(m)\to \varphi\)。

再沿偶数分辨率比较 \(F_{m/2}\)。由 Binet 公式，
\[
F_{m/2}\sim \frac{\varphi^{m/2}}{\sqrt5},
\qquad
d_{\min}(m)\sim \varphi^{-1}\left(\frac{2}{\varphi}\right)^m,
\]
从而
\[
\frac{d_{\min}(m)}{F_{m/2}}
\sim
\sqrt5\,\varphi^{-1}
\left(\frac{2}{\varphi^{3/2}}\right)^m.
\]
由于 \(\varphi^3=2+\sqrt5>4\)，故 \(\varphi^{3/2}>2\)，于是上式趋于 \(0\)。
\end{proof}

\begin{theorem}[中心投影的分位二原子塌缩]\label{thm:xi-time-part70ab-uniform-quantile-two-atom-collapse}
\leanverified{paper\_xi\_time\_part70ab\_uniform\_quantile\_two\_atom\_collapse}
设 \(U_m\) 服从 \(X_m\) 上的均匀分布，并定义尺度化随机变量
\[
Y_m:=\left(\frac{\varphi}{2}\right)^m d_m(U_m).
\]
再定义左连续分位函数
\[
Q_m(\alpha):=\inf\{t\ge 0:\ \mathbf P(Y_m\le t)\ge \alpha\},
\qquad
0<\alpha<1.
\]
则对每个
\[
\alpha\in(0,1)\setminus\{\varphi^{-2}\},
\]
都有极限
\[
Q_m(\alpha)\longrightarrow Q(\alpha),
\qquad
Q(\alpha)=
\begin{cases}
\varphi^{-1}, & 0<\alpha<\varphi^{-2},\\[2mm]
1, & \varphi^{-2}<\alpha<1.
\end{cases}
\]
并且在临界分位 \(\alpha=\varphi^{-2}\) 处出现奇偶振荡：
\[
Q_{2n}(\varphi^{-2})\longrightarrow 1,
\qquad
Q_{2n+1}(\varphi^{-2})\longrightarrow \varphi^{-1}.
\]
\end{theorem}
\begin{proof}
取
\[
0<\varepsilon<\frac{1-\varphi^{-1}}{3}.
\]
由定理 \ref{thm:fold-bin-two-state-asymptotic}，对充分大的 \(m\) 与全部 \(w\in X_m\) 都有
\[
\left|\left(\frac{\varphi}{2}\right)^m d_m(w)-\varphi^{-w_m}\right|
\le
C\left(\frac{\varphi}{2}\right)^m<\varepsilon.
\]
因此当 \(m\) 充分大时，
\[
Y_m\in[\varphi^{-1}-\varepsilon,\varphi^{-1}+\varepsilon]
\quad\text{在 }A_1^{(m)}\text{ 上},
\]
\[
Y_m\in[1-\varepsilon,1+\varepsilon]
\quad\text{在 }A_0^{(m)}\text{ 上}.
\]
故分位值只由块质量
\[
\mathbf P\bigl(Y_m\le \varphi^{-1}+\varepsilon\bigr)
=
\frac{|A_1^{(m)}|}{|X_m|}
=
\frac{F_m}{F_{m+2}}
\]
决定。若 \(\alpha<\varphi^{-2}\) 或 \(\alpha>\varphi^{-2}\)，由
\[
\frac{F_m}{F_{m+2}}\longrightarrow \varphi^{-2}
\]
知对充分大的 \(m\)，比较关系 \(\alpha\lessgtr F_m/F_{m+2}\) 不再改变，于是 \(Q_m(\alpha)\) 分别落在左原子簇与右原子簇中，令 \(\varepsilon\to 0\) 即得所示极限。

最后，由 Binet 公式可直接计算
\[
\frac{F_m}{F_{m+2}}-\varphi^{-2}
=
\frac{(-1)^{m+1}(1-\varphi^{-4})\varphi^{-m}}{\sqrt5\,F_{m+2}},
\]
其符号恰为 \((-1)^{m+1}\)。因此当 \(m\) 为奇数时，
\[
\frac{F_m}{F_{m+2}}>\varphi^{-2},
\]
故 \(Q_m(\varphi^{-2})\) 落在左原子簇；当 \(m\) 为偶数时，
\[
\frac{F_m}{F_{m+2}}<\varphi^{-2},
\]
故 \(Q_m(\varphi^{-2})\) 落在右原子簇。令 \(\varepsilon\to 0\) 即得两条子序列极限。
\end{proof}

\begin{theorem}[escort 温度族的分位相图]\label{thm:xi-time-part70ab-escort-quantile-phase-diagram}
对任意 \(q\ge 0\)，令
\[
W_m^{(q)}\sim \pi_{m,q},
\qquad
Y_m^{(q)}:=\left(\frac{\varphi}{2}\right)^m d_m\bigl(W_m^{(q)}\bigr),
\qquad
\theta(q):=\frac{1}{1+\varphi^{q+1}}.
\]
并定义左连续分位函数
\[
Q_{m,q}(\alpha):=\inf\{t\ge 0:\ \mathbf P(Y_m^{(q)}\le t)\ge \alpha\},
\qquad
0<\alpha<1.
\]
则对每个固定 \(q\ge 0\) 与每个
\[
\alpha\in(0,1)\setminus\{\theta(q)\},
\]
都有极限
\[
Q_{m,q}(\alpha)\longrightarrow Q_q(\alpha),
\qquad
Q_q(\alpha)=
\begin{cases}
\varphi^{-1}, & 0<\alpha<\theta(q),\\[2mm]
1, & \theta(q)<\alpha<1.
\end{cases}
\]
\end{theorem}
\begin{proof}
同样取
\[
0<\varepsilon<\frac{1-\varphi^{-1}}{3}.
\]
由定理 \ref{thm:fold-bin-two-state-asymptotic}，对充分大的 \(m\) 与全部 \(w\in X_m\) 都有
\[
\left|\left(\frac{\varphi}{2}\right)^m d_m(w)-\varphi^{-w_m}\right|<\varepsilon.
\]
因此当 \(m\) 充分大时，
\[
Y_m^{(q)}\in[\varphi^{-1}-\varepsilon,\varphi^{-1}+\varepsilon]
\quad\text{在事件 }W_m^{(q)}(m)=1\text{ 上},
\]
\[
Y_m^{(q)}\in[1-\varepsilon,1+\varepsilon]
\quad\text{在事件 }W_m^{(q)}(m)=0\text{ 上}.
\]
另一方面，由命题 \ref{prop:fold-bin-escort-lastbit}，
\[
\mathbf P\bigl(W_m^{(q)}(m)=1\bigr)=\theta(q)+O\!\left(\left(\frac{\varphi}{2}\right)^m\right).
\]
于是当 \(\alpha<\theta(q)\) 时，充分大的 \(m\) 上有 \(\alpha<\mathbf P(W_m^{(q)}(m)=1)\)，故 \(Q_{m,q}(\alpha)\) 落在左原子簇；当 \(\alpha>\theta(q)\) 时，则 \(Q_{m,q}(\alpha)\) 落在右原子簇。令 \(\varepsilon\to 0\) 即得结论。
\end{proof}

\begin{theorem}[中心容量逆问题的闭式解]\label{thm:xi-time-part70ab-center-capacity-inverse-threshold-closed}
\leanverified{paper\_xi\_time\_part70ab\_center\_capacity\_inverse\_threshold\_closed}
定义归一化连续容量
\[
\Psi_m(\tau):=
2^{-m}\,\mathcal C_m^{\mathrm{bin,cont}}(\tau A_m),
\qquad
\tau\ge 0,
\]
并定义达到水平 \(\rho\) 的最小归一化阈值
\[
\tau_\ast(\rho):=
\inf\left\{\tau\ge 0:\ \liminf_{m\to\infty}\Psi_m(\tau)\ge \rho\right\},
\qquad
0\le \rho\le 1.
\]
则有严格闭式
\[
\tau_\ast(\rho)=
\begin{cases}
\displaystyle \frac{\sqrt5}{\varphi^2}\,\rho,
& 0\le \rho\le \dfrac{\varphi}{\sqrt5},\\[3mm]
\displaystyle \frac{\sqrt5\,\rho-\varphi^{-1}}{\varphi},
& \dfrac{\varphi}{\sqrt5}\le \rho\le 1.
\end{cases}
\]
\end{theorem}
\begin{proof}
由定理 \ref{thm:xi-time-part70b-center-capacity-two-kink-skeleton}，
\[
\Psi_m(\tau)\longrightarrow \Psi(\tau):=
\frac{\varphi}{\sqrt5}\min\{1,\tau\}
+
\frac{1}{\sqrt5}\min\{\varphi^{-1},\tau\}.
\]
因此
\[
\tau_\ast(\rho)=\inf\{\tau\ge 0:\ \Psi(\tau)\ge \rho\}.
\]
而
\[
\Psi(\tau)=
\begin{cases}
\dfrac{\varphi^2}{\sqrt5}\tau, & 0\le \tau\le \varphi^{-1},\\[6pt]
\dfrac{\varphi}{\sqrt5}\tau+\dfrac{1}{\varphi\sqrt5}, & \varphi^{-1}\le \tau\le 1,\\[8pt]
1, & \tau\ge 1,
\end{cases}
\]
故广义逆由两段线性方程直接求得，即为所示公式。
\end{proof}

\begin{theorem}[escort 响应的逆温度阈值公式]\label{thm:xi-time-part70ab-escort-response-inverse-threshold-closed}
对固定 \(q\ge 0\)，定义 escort 归一化响应
\[
\Psi_{m,q}^{\mathrm{esc}}(\tau):=
\mathbb E_{\pi_{m,q}}
\left[
\min\!\left(1,\frac{\tau A_m}{d_m(W_m^{(q)})}\right)
\right],
\qquad
\tau\ge 0.
\]
再定义达到水平 \(\rho\) 的最小归一化阈值
\[
\tau_{q,\ast}(\rho):=
\inf\left\{\tau\ge 0:\ \liminf_{m\to\infty}\Psi_{m,q}^{\mathrm{esc}}(\tau)\ge \rho\right\},
\qquad
0\le \rho\le 1.
\]
记
\[
\rho_c(q):=\theta(q)+\varphi^{-1}\bigl(1-\theta(q)\bigr).
\]
则有闭式
\[
\tau_{q,\ast}(\rho)=
\begin{cases}
\displaystyle \frac{\rho}{1+(\varphi-1)\theta(q)},
& 0\le \rho\le \rho_c(q),\\[3mm]
\displaystyle \frac{\rho-\theta(q)}{1-\theta(q)},
& \rho_c(q)\le \rho\le 1.
\end{cases}
\]
\end{theorem}
\begin{proof}
由定理 \ref{thm:xi-time-part70b-escort-two-kink-response}，
\[
\Psi_{m,q}^{\mathrm{esc}}(\tau)\longrightarrow
\Psi_q^{\mathrm{esc}}(\tau):=
\bigl(1-\theta(q)\bigr)\min\{1,\tau\}
+
\theta(q)\min\{1,\varphi\tau\}.
\]
因此
\[
\tau_{q,\ast}(\rho)=\inf\{\tau\ge 0:\ \Psi_q^{\mathrm{esc}}(\tau)\ge \rho\}.
\]
把 \(\Psi_q^{\mathrm{esc}}\) 在 \([0,\varphi^{-1}]\) 与 \([\varphi^{-1},1]\) 上分段展开，得到
\[
\Psi_q^{\mathrm{esc}}(\tau)=
\begin{cases}
\bigl(1+(\varphi-1)\theta(q)\bigr)\tau,
& 0\le \tau\le \varphi^{-1},\\[6pt]
\theta(q)+\bigl(1-\theta(q)\bigr)\tau,
& \varphi^{-1}\le \tau\le 1,\\[6pt]
1, & \tau\ge 1.
\end{cases}
\]
其连接点函数值为
\[
\Psi_q^{\mathrm{esc}}(\varphi^{-1})
=
\theta(q)+\varphi^{-1}\bigl(1-\theta(q)\bigr)
=
\rho_c(q).
\]
故广义逆由上述两段线性函数分别求反即可。
\end{proof}

\begin{theorem}[escort 熵谱的完整常数项]\label{thm:xi-time-part70ab-escort-renyi-entropy-full-constant}
对固定 \(q\ge 0\)、\(a>0\) 且 \(a\neq 1\)，定义 escort Rényi 熵
\[
H_a(\pi_{m,q})
:=
\frac{1}{1-a}\log\sum_{w\in X_m}\pi_{m,q}(w)^a.
\]
则有渐近展开
\[
H_a(\pi_{m,q})
=
m\log\varphi
-\frac12\log 5
+
\frac{1}{1-a}
\log\!\Bigl(
\varphi^{1-a}\bigl(1-\theta(q)\bigr)^a+\theta(q)^a
\Bigr)
+o(1).
\]
当 \(a\to 1\) 时，得到 Shannon 熵常数项
\[
H(\pi_{m,q})
=
m\log\varphi
-\frac12\log 5
+
\bigl(1-\theta(q)\bigr)\log\varphi
+
h\bigl(\theta(q)\bigr)
+o(1),
\]
其中
\[
h(t):=-t\log t-(1-t)\log(1-t).
\]
\end{theorem}
\begin{proof}
记
\[
S_t(m):=\sum_{w\in X_m} d_m(w)^t.
\]
由定理 \ref{thm:xi-time-part70aa-binfold-all-real-moments-closed} 与 Binet 公式，对每个固定实数 \(t\) 都有
\[
S_t(m)
=
\frac{\varphi+\varphi^{-t}}{\sqrt5}
\bigl(2^t\varphi^{1-t}\bigr)^m
\bigl(1+o(1)\bigr).
\]
因此
\[
\sum_{w\in X_m}\pi_{m,q}(w)^a
=
\frac{S_{aq}(m)}{S_q(m)^a}
=
5^{(a-1)/2}\,
\varphi^{(1-a)m}\,
\frac{\varphi+\varphi^{-aq}}{(\varphi+\varphi^{-q})^a}
\bigl(1+o(1)\bigr).
\]
取对数并除以 \(1-a\)，得到
\[
H_a(\pi_{m,q})
=
m\log\varphi
-\frac12\log 5
+
\frac{1}{1-a}
\log\!\frac{\varphi+\varphi^{-aq}}{(\varphi+\varphi^{-q})^a}
+o(1).
\]
再由
\[
\theta(q)=\frac{\varphi^{-q}}{\varphi+\varphi^{-q}},
\qquad
1-\theta(q)=\frac{\varphi}{\varphi+\varphi^{-q}},
\]
可将最后一项重写为
\[
\frac{\varphi+\varphi^{-aq}}{(\varphi+\varphi^{-q})^a}
=
\varphi^{1-a}\bigl(1-\theta(q)\bigr)^a+\theta(q)^a.
\]
这就给出第一条公式。

当 \(a\to 1\) 时，对
\[
F_q(a):=
\log\!\Bigl(
\varphi^{1-a}\bigl(1-\theta(q)\bigr)^a+\theta(q)^a
\Bigr)
\]
应用 l'Hospital 法则，得到
\[
\lim_{a\to 1}\frac{F_q(a)}{1-a}
=
-F_q'(1)
=
\bigl(1-\theta(q)\bigr)\log\varphi+h\bigl(\theta(q)\bigr).
\]
代回即得 Shannon 常数项。
\end{proof}

\begin{theorem}[容量极限曲线的负矩 Mellin 反演]\label{thm:xi-time-part70ab-capacity-limit-negative-moment-mellin}
定义中心容量极限
\[
m_\ast(\tau):=
\lim_{m\to\infty}
\frac{\sqrt5}{\varphi^2}\,
2^{-m}\,\mathcal C_m^{\mathrm{bin,cont}}(\tau A_m),
\qquad
\tau\ge 0.
\]
则对任意实数 \(s>0\)，成立严格恒等式
\[
s(s+1)\int_0^\infty \tau^{-s-2}\bigl(\tau-m_\ast(\tau)\bigr)\,\dd\tau
=
\varphi^{-1}+\varphi^{s-2}.
\]
\end{theorem}
\begin{proof}
由定理 \ref{thm:xi-time-part70b-center-capacity-two-kink-skeleton}，
\[
m_\ast(\tau)=
\varphi^{-1}\min\{1,\tau\}
+
\varphi^{-2}\min\{\varphi^{-1},\tau\}.
\]
因而若定义双原子概率测度
\[
\mu_\ast:=\varphi^{-2}\delta_{\varphi^{-1}}+\varphi^{-1}\delta_1,
\]
则有
\[
m_\ast(\tau)=\int \min\{y,\tau\}\,\dd\mu_\ast(y).
\]
对任意固定 \(y>0\) 与 \(s>0\)，直接计算得
\[
\int_0^\infty \tau^{-s-2}\bigl(\tau-\min\{y,\tau\}\bigr)\,\dd\tau
=
\int_y^\infty \tau^{-s-2}(\tau-y)\,\dd\tau
=
\frac{y^{-s}}{s(s+1)}.
\]
两边乘以 \(s(s+1)\) 并对 \(\mu_\ast\) 积分，即得
\[
s(s+1)\int_0^\infty \tau^{-s-2}\bigl(\tau-m_\ast(\tau)\bigr)\,\dd\tau
=
\int y^{-s}\,\dd\mu_\ast(y)
=
\varphi^{-1}+\varphi^{-2}\varphi^s
=
\varphi^{-1}+\varphi^{s-2}.
\]
证毕。
\end{proof}

\begin{corollary}[Stirling--Bernoulli 几何系数的容量解释]\label{cor:xi-time-part70ab-bernoulli-hierarchy-capacity-interpretation}
\leanverified{paper\_xi\_time\_part70ab\_bernoulli\_hierarchy\_capacity\_interpretation}
对任意整数 \(r\ge 1\)，都有
\[
\varphi^{-1}+\varphi^{2r-3}
=
(2r-1)(2r)\int_0^\infty \tau^{-2r-1}\bigl(\tau-m_\ast(\tau)\bigr)\,\dd\tau.
\]
从而 Stirling--Bernoulli 层级中第 \(r\) 级系数可写为
\[
\frac{B_{2r}}{r(2r-1)}
\bigl(\varphi^{-1}+\varphi^{2r-3}\bigr)
=
\frac{2B_{2r}}{r}
\int_0^\infty \tau^{-2r-1}\bigl(\tau-m_\ast(\tau)\bigr)\,\dd\tau.
\]
特别地，对任意固定 \(R\ge 1\)，定理 \ref{thm:fold-bin-gauge-constant-stirling-bernoulli-hierarchy} 的展开可重写为
\[
\log C_m
=
\log(2\pi)
+
\sum_{r=1}^{R}
\frac{2B_{2r}}{r}
\left(
\int_0^\infty \tau^{-2r-1}\bigl(\tau-m_\ast(\tau)\bigr)\,\dd\tau
\right)
\left(\frac{\varphi}{2}\right)^{(2r-1)m}
+
O\!\left(\left(\frac{\varphi}{2}\right)^{(2R+1)m}\right).
\]
\end{corollary}
\begin{proof}
在定理 \ref{thm:xi-time-part70ab-capacity-limit-negative-moment-mellin} 中取 \(s=2r-1\) 即得第一式。再把该恒等式代入定理 \ref{thm:fold-bin-gauge-constant-stirling-bernoulli-hierarchy} 的系数表达，便得到后两式。
\end{proof}

\noindent
由此，末位 \(0/1\) 的二态骨架不只控制纤维多重度的最终排序与均匀/escort 采样下的分位相图，还把中心容量曲线的广义逆、escort 温度响应的广义逆以及 \(\log C_m\) 的 Stirling--Bernoulli 几何权重全部还原为同一组双原子数据。于是，阈值/采样轴、路径/相位轴与 Bernoulli 塔的解析层级便在同一条末位二态主线上闭合。

\endinput

\paragraph{全息地址前缀商的精确熵律与最小判别比特}
在 H.160--H.161 的 faithful \(2^L\)-进地址表示、精确柱分离与圆柱--球识别已经固定之后，有限深度地址商的可见信息量还可进一步压缩为一条完全闭式的计数律。

\begin{theorem}[\texorpdfstring{$2^L$}{2^L}-进全息地址前缀商的精确计数与最小判别比特]\label{thm:xi-time-part70ac-hologram-prefix-quotient-entropy-bits}
\leanverified{paper\_xi\_time\_part70ac\_hologram\_prefix\_quotient\_entropy\_bits}
设 \(E\) 为有限事件集，
\[
\mathrm{code}:E\hookrightarrow \NN,
\qquad
M:=\max_{e\in E}\mathrm{code}(e),
\qquad
B:=2^L>M,
\]
并记
\[
T:=T_2(E_{\min}),
\qquad
v\in T\setminus 2T,
\]
其中
\[
X(\mathbf e):=\sum_{t\ge 1}\mathrm{code}(e_t)\,B^{t-1}v
\qquad
\bigl(\mathbf e=(e_t)_{t\ge 1}\in E^{\NN}\bigr).
\]
并定义
\[
\mathcal X_{E,v}:=X(E^{\NN}).
\]
对每个 \(n\ge 1\)，定义第 \(n\) 级地址商集
\[
\mathcal A_n
:=
\left\{
X(\mathbf e)\bmod B^nT:\ \mathbf e\in E^{\NN}
\right\}
\subseteq T/B^nT.
\]
再对每个长度 \(n\) 前缀
\[
\mathbf a=(a_1,\dots,a_n)\in E^n
\]
记
\[
[\mathbf a]
:=
\left\{
\mathbf e\in E^{\NN}:\ \mathrm{pref}_n(\mathbf e)=\mathbf a
\right\},
\qquad
x_{\mathbf a}
:=
\sum_{t=1}^{n}\mathrm{code}(a_t)\,B^{t-1}v.
\]
则成立：
\begin{enumerate}
  \item 地址商集满足精确计数公式
  \[
  |\mathcal A_n|=|E|^n.
  \]
  \item 若定义深度 \(n\) 的 clopen 地址柱代数
  \[
  \mathcal B_n
  :=
  \left\{
  \bigcup_{\mathbf a\in U}X([\mathbf a]):\ U\subseteq E^n
  \right\},
  \]
  则
  \[
  |\mathcal B_n|=2^{|E|^n}.
  \]
  \item 若 \(A\) 为有限集，且读出映射
  \[
  R:\mathcal X_{E,v}\longrightarrow A
  \]
  满足对任意不同前缀 \(\mathbf a\neq \mathbf b\in E^n\) 及任意
  \[
  x\in X([\mathbf a]),
  \qquad
  y\in X([\mathbf b]),
  \]
  都有 \(R(x)\neq R(y)\)，则
  \[
  |A|\ge |E|^n.
  \]
  因而区分全部长度 \(n\) 前缀所需的最小判别比特数恰为
  \[
  \left\lceil n\log_2|E|\right\rceil.
  \]
\end{enumerate}
\end{theorem}
\begin{proof}
固定 \(n\ge 1\)。定义映射
\[
\Phi_n:E^n\longrightarrow \mathcal A_n,
\qquad
\mathbf a\longmapsto x_{\mathbf a}+B^nT.
\]
先证满射。任取 \(\mathbf e=(e_t)_{t\ge 1}\in E^{\NN}\)，并记
\[
\mathbf a:=\mathrm{pref}_n(\mathbf e)=(e_1,\dots,e_n).
\]
则
\[
X(\mathbf e)-x_{\mathbf a}
=
\sum_{t>n}\mathrm{code}(e_t)\,B^{t-1}v
\in B^nT,
\]
故
\[
X(\mathbf e)\bmod B^nT=\Phi_n(\mathbf a).
\]
因此 \(\Phi_n\) 满射。

再证单射。若
\[
\Phi_n(\mathbf a)=\Phi_n(\mathbf b),
\]
则
\[
x_{\mathbf a}-x_{\mathbf b}\in B^nT.
\]
任取
\[
\mathbf e\in[\mathbf a],
\qquad
\mathbf f\in[\mathbf b].
\]
由结论 \ref{con:xi-time-part62dgb-hologram-cylinder-ball-isometry}，
\[
X(\mathbf e)\equiv x_{\mathbf a}\pmod{B^nT},
\qquad
X(\mathbf f)\equiv x_{\mathbf b}\pmod{B^nT},
\]
从而
\[
X(\mathbf e)\equiv X(\mathbf f)\pmod{B^nT}.
\]
再由同一结论中的同余--前缀等价式，只能有
\[
\mathrm{pref}_n(\mathbf e)=\mathrm{pref}_n(\mathbf f),
\]
亦即 \(\mathbf a=\mathbf b\)。故 \(\Phi_n\) 为双射，于是
\[
|\mathcal A_n|=|E^n|=|E|^n,
\]
第 \(1\) 条得证。

对第 \(2\) 条，由定理 \ref{thm:xi-time-part62d-hologram-image-cantor-homeomorphism}，映射
\[
X:E^{\NN}\longrightarrow \mathcal X_{E,v}
\]
是拓扑同胚。长度 \(n\) 的 cylinder 集族
\[
\{[\mathbf a]:\ \mathbf a\in E^n\}
\]
在 \(E^{\NN}\) 中构成一个由 \(|E|^n\) 个原子组成的 clopen 分割，故其像
\[
\{X([\mathbf a]):\ \mathbf a\in E^n\}
\]
在 \(\mathcal X_{E,v}\) 中同样构成一个由 \(|E|^n\) 个原子组成的 clopen 分割。于是由这些原子生成的 Boolean 代数正是全部原子并的集合族，故
\[
|\mathcal B_n|=2^{|E|^n}.
\]

再证第 \(3\) 条。对每个 \(\mathbf a\in E^n\)，取一点
\[
x_{\mathbf a}^{\ast}\in X([\mathbf a]),
\]
这是可能的，因为每个 cylinder 都非空。若 \(\mathbf a\neq \mathbf b\)，则按假设
\[
R(x_{\mathbf a}^{\ast})\neq R(x_{\mathbf b}^{\ast}).
\]
因此映射
\[
E^n\longrightarrow A,
\qquad
\mathbf a\longmapsto R(x_{\mathbf a}^{\ast})
\]
是单射，故
\[
|A|\ge |E|^n.
\]
这给出输出值个数的下界。

反向上，由定理 \ref{thm:xi-time-part62d-hologram-image-cantor-homeomorphism}，\(X\) 为双射，故前缀读出
\[
R_n:\mathcal X_{E,v}\longrightarrow E^n,
\qquad
R_n(x):=\mathrm{pref}_n(X^{-1}(x))
\]
良定义，并且恰取 \(|E|^n\) 个值。于是最小所需输出值个数正好等于 \(|E|^n\)。将有限输出值编码为二进制字，最小判别比特数遂为
\[
\left\lceil \log_2|E|^n\right\rceil
=
\left\lceil n\log_2|E|\right\rceil.
\]
证毕。
\end{proof}

\noindent
因此，深度 \(n\) 的地址可见层并非只在数量级上与 \(|E|^n\) 同阶，而是恰由 \(|E|^n\) 个前缀原子严格组成；相应有限值读出的最小复杂度也被精确锁定为 \(\left\lceil n\log_2|E|\right\rceil\) 比特。于是，时间推进所带来的地址精度提升被压缩为一条完全刚性的有限计数律。

\endinput

\paragraph{bin-fold 输出分布的全阶信息几何常数层闭式}
在输出分布的 Rényi 熵率塌缩、隐藏对数多重度的二原子 Laplace 极限、规范体积密度的两尺度展开以及碰撞矩的两态主项都已建立之后，还可把 bin-fold 推前分布 \(\mu_m\) 的常数层信息几何进一步压缩为下列四条闭式结论。

\begin{theorem}[输出分布的全阶 Rényi 熵具有精确常数漂移]\label{thm:xi-time-part70ad-binfold-renyi-entropy-constant-drift}
\leanverified{paper\_xi\_time\_part70ad\_binfold\_renyi\_entropy\_constant\_drift}
记
\[
\mu_m(w):=\frac{d_m^{\mathrm{bin}}(w)}{2^m}
\qquad (w\in X_m),
\]
并对 \(q>0\)、\(q\neq 1\) 定义
\[
H_q(\mu_m):=\frac{1}{1-q}\log\sum_{w\in X_m}\mu_m(w)^q.
\]
则对每个固定 \(q>0\)、\(q\neq 1\)，都有精确渐近
\[
\boxed{\;
H_q(\mu_m)
=
m\log\varphi
+
\frac{1}{1-q}\log\!\Bigl(\frac{\varphi+\varphi^{-q}}{\sqrt5}\Bigr)
+
O_q\!\Bigl(\Bigl(\frac{\varphi}{2}\Bigr)^m\Bigr).\;}
\]
特别地，
\[
\boxed{\;
\lim_{m\to\infty}\frac{1}{m}H_q(\mu_m)=\log\varphi
\qquad (q>0,\ q\neq 1).\;}
\]
\end{theorem}
\begin{proof}
记
\[
A_0^{(m)}:=\{w\in X_m:\ w_m=0\},
\qquad
A_1^{(m)}:=\{w\in X_m:\ w_m=1\}.
\]
由定理 \ref{thm:fold-bin-two-state-asymptotic} 的一致主项，对每个固定 \(q>0\) 都存在常数 \(C_q>0\)，使得对充分大的 \(m\) 与全部 \(w\in X_m\) 都有
\[
\bigl(d_m^{\mathrm{bin}}(w)\bigr)^q
=
\left(\frac{2}{\varphi}\right)^{qm}
\begin{cases}
1+O_q\!\bigl((\varphi/2)^m\bigr), & w\in A_0^{(m)},\\[4pt]
\varphi^{-q}+O_q\!\bigl((\varphi/2)^m\bigr), & w\in A_1^{(m)}.
\end{cases}
\]
故
\[
S_q^{\mathrm{bin}}(m):=\sum_{w\in X_m}\bigl(d_m^{\mathrm{bin}}(w)\bigr)^q
=
\left(\frac{2}{\varphi}\right)^{qm}
\Bigl(
F_{m+1}+\varphi^{-q}F_m
\Bigr)
\Bigl(1+O_q\!\bigl((\varphi/2)^m\bigr)\Bigr),
\]
因为 \(|A_0^{(m)}|=F_{m+1}\)、\(|A_1^{(m)}|=F_m\)。于是
\[
\sum_{w\in X_m}\mu_m(w)^q
=
2^{-qm}S_q^{\mathrm{bin}}(m)
=
\varphi^{-qm}
\Bigl(
F_{m+1}+\varphi^{-q}F_m
\Bigr)
\Bigl(1+O_q\!\bigl((\varphi/2)^m\bigr)\Bigr).
\]
再由 Binet 公式，
\[
F_{m+1}+\varphi^{-q}F_m
=
\frac{\varphi^m}{\sqrt5}\bigl(\varphi+\varphi^{-q}\bigr)+O(\varphi^{-m}),
\]
而 \(O(\varphi^{-m})=O((\varphi/2)^m\varphi^m)\)，故
\[
\sum_{w\in X_m}\mu_m(w)^q
=
\varphi^{(1-q)m}
\frac{\varphi+\varphi^{-q}}{\sqrt5}
\Bigl(1+O_q\!\bigl((\varphi/2)^m\bigr)\Bigr).
\]
右端主因子严格为正，取对数并除以 \(1-q\) 即得
\[
H_q(\mu_m)
=
m\log\varphi
+
\frac{1}{1-q}\log\!\Bigl(\frac{\varphi+\varphi^{-q}}{\sqrt5}\Bigr)
+
O_q\!\Bigl(\Bigl(\frac{\varphi}{2}\Bigr)^m\Bigr).
\]
再除以 \(m\) 并令 \(m\to\infty\)，便得到熵率极限。证毕。
\end{proof}

\leanverified{paper\_xi\_time\_part70ad\_hidden\_logmultiplicity\_exponential\_moments}
\begin{theorem}[中心化对数多重度的全指数矩极限与全部累积量闭式]\label{thm:xi-time-part70ad-hidden-logmultiplicity-exponential-moments}
定义
\[
Y_m(w):=\log d_m^{\mathrm{bin}}(w)-m\log\!\Bigl(\frac{2}{\varphi}\Bigr),
\qquad
W_m\sim \mu_m.
\]
则对任意固定复数 \(t\in\CC\)，都有
\[
\boxed{\;
\EE_{\mu_m}\!\bigl[e^{tY_m(W_m)}\bigr]
=
\frac{\varphi+\varphi^{-(t+1)}}{\sqrt5}
+
O_t\!\Bigl(\Bigl(\frac{\varphi}{2}\Bigr)^m\Bigr).\;}
\]
因此 \(Y_m(W_m)\) 在分布上并且在全部指数矩意义下收敛到两点随机变量
\[
Y_\infty=
\begin{cases}
0,& \text{概率 } \dfrac{\varphi}{\sqrt5},\\[6pt]
-\log\varphi,& \text{概率 } \dfrac{1}{\varphi\sqrt5}.
\end{cases}
\]
若记
\[
p:=\frac{1}{\varphi\sqrt5},
\qquad
\ell:=\log\varphi,
\]
则其极限累积量生成函数为
\[
K_\infty(t)=\log\bigl((1-p)+pe^{-t\ell}\bigr),
\]
并且前四阶极限累积量满足
\[
\boxed{\;
\kappa_1^\infty=-p\ell,\qquad
\kappa_2^\infty=p(1-p)\ell^2,\qquad
\kappa_3^\infty=-p(1-p)(1-2p)\ell^3,\qquad
\kappa_4^\infty=p(1-p)(1-6p+6p^2)\ell^4.\;}
\]
\end{theorem}
\begin{proof}
记 \(B_m:=W_m(m)\in\{0,1\}\)。由定理 \ref{thm:fold-bin-two-state-asymptotic} 的一致对数展开，
\[
Y_m(w)=-w_m\log\varphi+O\!\Bigl(\Bigl(\frac{\varphi}{2}\Bigr)^m\Bigr)
\qquad (w\in X_m),
\]
于是对任意固定 \(t\in\CC\)，都有一致展开
\[
e^{tY_m(w)}
=
\varphi^{-t w_m}
\Bigl(1+O_t\!\Bigl(\Bigl(\frac{\varphi}{2}\Bigr)^m\Bigr)\Bigr).
\]
对 \(\mu_m\) 取期望，并利用命题 \ref{prop:fold-bin-escort-lastbit} 在 \(q=1\) 处的末位质量公式
\[
\mu_m(B_m=0)=\frac{\varphi}{\sqrt5}+O\!\Bigl(\Bigl(\frac{\varphi}{2}\Bigr)^m\Bigr),
\qquad
\mu_m(B_m=1)=\frac{1}{\varphi\sqrt5}+O\!\Bigl(\Bigl(\frac{\varphi}{2}\Bigr)^m\Bigr),
\]
便得到
\[
\EE_{\mu_m}\!\bigl[e^{tY_m(W_m)}\bigr]
=
\frac{\varphi}{\sqrt5}
+
\frac{1}{\varphi\sqrt5}\varphi^{-t}
+
O_t\!\Bigl(\Bigl(\frac{\varphi}{2}\Bigr)^m\Bigr)
=
\frac{\varphi+\varphi^{-(t+1)}}{\sqrt5}
+
O_t\!\Bigl(\Bigl(\frac{\varphi}{2}\Bigr)^m\Bigr).
\]
这就得到指数矩极限。

由于 \(Y_m\) 一致有界，上式同时给出所有指数矩收敛，并推出 \(Y_m(W_m)\Rightarrow Y_\infty\)，其中
\[
\PP(Y_\infty=0)=1-p=\frac{\varphi}{\sqrt5},
\qquad
\PP(Y_\infty=-\ell)=p=\frac{1}{\varphi\sqrt5}.
\]
故极限累积量生成函数即为
\[
K_\infty(t)=\log\EE[e^{tY_\infty}]
=
\log\bigl((1-p)+pe^{-t\ell}\bigr).
\]
对 \(t=0\) 逐阶求导，便得到
\[
\kappa_1^\infty=-p\ell,
\qquad
\kappa_2^\infty=p(1-p)\ell^2,
\]
\[
\kappa_3^\infty=-p(1-p)(1-2p)\ell^3,
\qquad
\kappa_4^\infty=p(1-p)(1-6p+6p^2)\ell^4.
\]
证毕。
\end{proof}

\begin{theorem}[规范群熵的精确首项与常数漂移]\label{thm:xi-time-part70ad-gauge-entropy-leading-drift}
记
\[
G_m^{\mathrm{bin}}
\cong
\prod_{w\in X_m}S_{d_m^{\mathrm{bin}}(w)},
\qquad
H_m^{\mathrm{gauge}}:=\log_2|G_m^{\mathrm{bin}}|.
\]
则其自然对数版本满足
\[
\boxed{\;
\frac{1}{2^m}\log |G_m^{\mathrm{bin}}|
=
m\log\!\Bigl(\frac{2}{\varphi}\Bigr)
-1
-\frac{\log\varphi}{\varphi\sqrt5}
+o(1).\;}
\]
等价地，以 bits 计，
\[
\boxed{\;
H_m^{\mathrm{gauge}}
=
2^m\left[
m\log_2\!\Bigl(\frac{2}{\varphi}\Bigr)
-\frac{1}{\log 2}
\Bigl(
1+\frac{\log\varphi}{\varphi\sqrt5}
\Bigr)
+o(1)
\right].\;}
\]
\end{theorem}
\begin{proof}
由定理 \ref{thm:xi-fold-hidden-logmultiplicity-kappa-gauge-two-scale-explicit}，若记
\[
g_m:=2^{-m}\log|G_m^{\mathrm{bin}}|,
\]
则有
\[
g_m
=
m\log\!\Bigl(\frac{2}{\varphi}\Bigr)
-\Bigl(1+\frac{\log\varphi}{1+\varphi^2}\Bigr)
+O\!\Bigl(\Bigl(\frac{\varphi}{2}\Bigr)^m\Bigr).
\]
再用恒等式
\[
\frac{1}{1+\varphi^2}=\frac{1}{\varphi\sqrt5},
\]
即得
\[
\frac{1}{2^m}\log |G_m^{\mathrm{bin}}|
=
m\log\!\Bigl(\frac{2}{\varphi}\Bigr)
-1
-\frac{\log\varphi}{\varphi\sqrt5}
+o(1).
\]
最后两边乘以 \(2^m/\log 2\)，便得到 bits 形式。证毕。
\end{proof}

\begin{theorem}[bin-fold 输出分布的 Tsallis 熵同样具有闭式常数层]\label{thm:xi-time-part70ad-binfold-tsallis-entropy-asymptotics}
\leanverified{paper\_xi\_time\_part70ad\_binfold\_tsallis\_entropy\_asymptotics}
对 \(q>0\)、\(q\neq 1\)，定义 Tsallis 熵
\[
T_q(\mu_m):=\frac{1-\sum_{w\in X_m}\mu_m(w)^q}{q-1}.
\]
则有如下两段渐近：
\begin{enumerate}
  \item 若 \(q>1\)，则
  \[
  \boxed{\;
  T_q(\mu_m)
  =
  \frac{1}{q-1}
  -
  \frac{1}{q-1}\,
  \varphi^{(1-q)m}
  \frac{\varphi+\varphi^{-q}}{\sqrt5}
  +
  O_q\!\Bigl(\varphi^{(1-q)m}\Bigl(\frac{\varphi}{2}\Bigr)^m\Bigr).\;}
  \]
  特别地，
  \[
  \boxed{\;
  T_q(\mu_m)\longrightarrow \frac{1}{q-1}.\;}
  \]
  \item 若 \(0<q<1\)，则
  \[
  \boxed{\;
  T_q(\mu_m)
  =
  \frac{1}{1-q}\,
  \varphi^{(1-q)m}
  \frac{\varphi+\varphi^{-q}}{\sqrt5}
  \bigl(1+o(1)\bigr).\;}
  \]
  因而
  \[
  \boxed{\;
  \lim_{m\to\infty}\frac{1}{m}\log T_q(\mu_m)
  =
  (1-q)\log\varphi.\;}
  \]
\end{enumerate}
\end{theorem}
\begin{proof}
由定理 \ref{thm:xi-time-part70ad-binfold-renyi-entropy-constant-drift} 的证明中间式，
\[
\sum_{w\in X_m}\mu_m(w)^q
=
\varphi^{(1-q)m}
\frac{\varphi+\varphi^{-q}}{\sqrt5}
\Bigl(1+O_q\!\Bigl(\Bigl(\frac{\varphi}{2}\Bigr)^m\Bigr)\Bigr).
\]
若 \(q>1\)，则 \(1-q<0\)，上式趋于 \(0\)。代入定义
\[
T_q(\mu_m)
=
\frac{1}{q-1}\left(1-\sum_{w\in X_m}\mu_m(w)^q\right),
\]
便立即得到第一段公式与极限。

若 \(0<q<1\)，则 \(1-q>0\)，故上式指数发散。于是
\[
T_q(\mu_m)
=
\frac{\sum_w\mu_m(w)^q-1}{1-q}
=
\frac{1}{1-q}\,
\varphi^{(1-q)m}
\frac{\varphi+\varphi^{-q}}{\sqrt5}
\bigl(1+o(1)\bigr),
\]
因为被减去的常数 \(1\) 相对于主项可忽略。最后取对数并除以 \(m\)，即可得到指数率
\[
\lim_{m\to\infty}\frac{1}{m}\log T_q(\mu_m)
=(1-q)\log\varphi.
\]
证毕。
\end{proof}

\noindent
由此，bin-fold 输出分布的全阶信息几何在常数层也达到闭合：Rényi 熵不只具有统一熵率，还具有显式 \(O(1)\) 漂移；中心化对数多重度不只具有一阶均值极限，而是整个指数矩与全部有限阶累积量都收敛到同一个二原子模型；规范群熵的指数体量则带有可审计的首项斜率与常数偏移；而 Tsallis 熵在 \(q>1\) 与 \(0<q<1\) 两侧的饱和值与增长率也由同一碰撞矩常数统一给出。于是，两态末位机制、碰撞矩闭式与规范体积展开最终在同一信息几何层面完全合流。

\endinput


\paragraph{中心双折点容量骨架、escort 温度重加权与 dark band 半阶尺度}
在 bin-fold 的中心临界容量极限、escort 温度族末位塌缩、尾部单跳阶梯以及六倍窗零块下界都已建立之后，还可把阈值/采样轴、相位/路径信息轴与 Fourier 侧 dark band 进一步压缩为同一组闭合接口。其共同特征是：中心层容量骨架始终只含两个不可消去折点，而温度参数只改变同一双折点骨架上的权重，不改变骨架本身。

\begin{theorem}[中心临界容量骨架的双折点闭式]\label{thm:xi-time-part70b-center-capacity-two-kink-skeleton}
定义
\[
A_m:=\left(\frac{2}{\varphi}\right)^m,
\qquad
\Psi_m(\tau):=
2^{-m}\,\mathcal C_m^{\mathrm{bin,cont}}(\tau A_m)
\qquad (\tau\ge 0).
\]
则对每个固定 \(\tau\ge 0\)，都有
\[
\Psi_m(\tau)\longrightarrow \Psi(\tau)
:=
\frac{\varphi}{\sqrt5}\min\{1,\tau\}
+
\frac{1}{\sqrt5}\min\{\varphi^{-1},\tau\}.
\]
等价地，
\[
\Psi(\tau)=
\begin{cases}
\dfrac{\varphi^2}{\sqrt5}\,\tau, & 0\le \tau\le \varphi^{-1},\\[6pt]
\dfrac{\varphi}{\sqrt5}\,\tau+\dfrac{1}{\varphi\sqrt5}, & \varphi^{-1}\le \tau\le 1,\\[8pt]
1, & \tau\ge 1.
\end{cases}
\]
特别地，中心临界容量骨架只有两个不可消去的折点，位置分别位于
\[
\tau=\varphi^{-1},
\qquad
\tau=1.
\]
\end{theorem}
\begin{proof}
这是定理 \ref{thm:xi-foldbin-dyadic-capacity-critical-window-limit} 与定理 \ref{thm:xi-time-part70-continuous-capacity-two-atom-second-derivative} 的直接改写。分段表达式由极限函数逐段展开得到，而折点位置则由分布二阶导数
\[
-D^2\Psi
=
\frac{1}{\sqrt5}\,\delta_{\varphi^{-1}}
+
\frac{\varphi}{\sqrt5}\,\delta_1
\]
唯一确定。证毕。
\end{proof}

\begin{theorem}[escort 温度族对中心双折点骨架的重加权响应]\label{thm:xi-time-part70b-escort-two-kink-response}
对任意固定 \(q\ge 0\)，令 \(W_m^{(q)}\sim \pi_{m,q}\)，并定义
\[
\Psi_{m,q}^{\mathrm{esc}}(\tau)
:=
\mathbb E_{\pi_{m,q}}
\left[
\min\!\left(
1,\frac{\tau A_m}{d_m^{\mathrm{bin}}(W_m^{(q)})}
\right)
\right]
\qquad (\tau\ge 0).
\]
则对每个固定 \(\tau\ge 0\)，极限
\[
\Psi_q^{\mathrm{esc}}(\tau):=
\lim_{m\to\infty}\Psi_{m,q}^{\mathrm{esc}}(\tau)
\]
存在，并满足闭式
\[
\Psi_q^{\mathrm{esc}}(\tau)
=
\bigl(1-\theta(q)\bigr)\min\{1,\tau\}
+
\theta(q)\min\{1,\varphi\tau\},
\qquad
\theta(q):=\frac{1}{1+\varphi^{q+1}}.
\]
因此，escort 温度参数 \(q\) 并不改变中心容量骨架的两个折点位置，而只在同一双折点框架上重新分配两条线性支的权重。
\end{theorem}
\begin{proof}
由定理 \ref{thm:fold-bin-two-state-asymptotic}，存在常数 \(C>0\)，使对充分大的 \(m\) 与全部 \(w\in X_m\) 都有
\[
\left|
\frac{d_m^{\mathrm{bin}}(w)}{A_m}-\varphi^{-w_m}
\right|
\le C\left(\frac{\varphi}{2}\right)^m.
\]
故对充分大的 \(m\)，上述比值一律落在紧区间
\[
\left[\frac{\varphi^{-1}}{2},\,2\right]
\]
内。对固定 \(\tau\ge 0\)，函数
\[
x\longmapsto \min\!\left(1,\frac{\tau}{x}\right)
\]
在该紧区间上 Lipschitz，因而存在常数 \(C_\tau>0\) 使
\[
\min\!\left(1,\frac{\tau A_m}{d_m^{\mathrm{bin}}(w)}\right)
=
\min\!\bigl(1,\tau\varphi^{w_m}\bigr)
+O\!\left(\left(\frac{\varphi}{2}\right)^m\right)
\]
在 \(w\in X_m\) 上一致成立。对 \(\pi_{m,q}\) 取期望，得到
\[
\Psi_{m,q}^{\mathrm{esc}}(\tau)
=
\mathbb E_{\pi_{m,q}}\!\bigl[\min(1,\tau\varphi^{W_m^{(q)}(m)})\bigr]
+O\!\left(\left(\frac{\varphi}{2}\right)^m\right).
\]
另一方面，由命题 \ref{prop:fold-bin-escort-lastbit}，
\[
\mathbb P\bigl(W_m^{(q)}(m)=1\bigr)
=
\theta(q)+O\!\left(\left(\frac{\varphi}{2}\right)^m\right),
\]
\[
\mathbb P\bigl(W_m^{(q)}(m)=0\bigr)
=
1-\theta(q)+O\!\left(\left(\frac{\varphi}{2}\right)^m\right).
\]
于是
\[
\Psi_{m,q}^{\mathrm{esc}}(\tau)
=
\bigl(1-\theta(q)\bigr)\min\{1,\tau\}
+
\theta(q)\min\{1,\varphi\tau\}
+O\!\left(\left(\frac{\varphi}{2}\right)^m\right).
\]
令 \(m\to\infty\)，即得所示极限公式。证毕。
\end{proof}

\begin{corollary}[中心双折点骨架上的温度斜率可辨识性]\label{cor:xi-time-part70b-escort-slope-identifiability}
沿用定理 \ref{thm:xi-time-part70b-escort-two-kink-response} 的记号。则在中间区间
\[
\tau\in (\varphi^{-1},1)
\]
上有
\[
\Psi_q^{\mathrm{esc}}(\tau)
=
\bigl(1-\theta(q)\bigr)\tau+\theta(q),
\]
从而导数恒为
\[
\frac{\dd}{\dd\tau}\Psi_q^{\mathrm{esc}}(\tau)
=
\sigma_q
:=
1-\theta(q)\in(0,1).
\]
并且仅由该斜率 \(\sigma_q\) 即可唯一恢复温度参数 \(q\)：
\[
q=\log_{\varphi}\!\left(\frac{\sigma_q}{1-\sigma_q}\right)-1.
\]
\end{corollary}
\begin{proof}
当 \(\tau\in(\varphi^{-1},1)\) 时，
\[
\min\{1,\tau\}=\tau,
\qquad
\min\{1,\varphi\tau\}=1,
\]
故定理 \ref{thm:xi-time-part70b-escort-two-kink-response} 中的极限曲线退化为
\[
\Psi_q^{\mathrm{esc}}(\tau)
=
\bigl(1-\theta(q)\bigr)\tau+\theta(q).
\]
于是导数恒为 \(1-\theta(q)\)。再由
\[
\theta(q)=\frac{1}{1+\varphi^{q+1}}
\]
可得
\[
\frac{1-\theta(q)}{\theta(q)}=\varphi^{q+1}.
\]
代入 \(\sigma_q=1-\theta(q)\) 并整理，即得所示反演公式。证毕。
\end{proof}

\begin{corollary}[固定末位层的阈值层纹由单跳尾和谱完全恢复]\label{cor:xi-time-part70b-lastbit-tail-jump-recovery}
沿用定理 \ref{thm:xi-foldbin-lastbit-tail-spectrum-single-jump-staircase} 的记号，并把其中的末位指标 \(i\) 记为 \(b\in\{0,1\}\)。定义
\[
\delta_{m,b}(V):=
D_{m,b}(V)
=
\#\bigl\{t\in \mathcal T_{m,b}: t\le 2^m-1-V\bigr\}.
\]
则对一切相容的 \(V\) 都有精确跳跃公式
\[
\delta_{m,b}(V)-\delta_{m,b}(V+1)
=
\mathbf 1_{\{\,2^m-1-V\in\mathcal T_{m,b}\,\}}.
\]
因而合法尾和集可由跳点完全恢复：
\[
\mathcal T_{m,b}
=
\bigl\{2^m-1-V:\ \delta_{m,b}(V)-\delta_{m,b}(V+1)=1\bigr\}.
\]
特别地，在固定末位相位内，每一次向下跳变都恰对应一个合法尾和被阈值切除。
\end{corollary}
\begin{proof}
这只是定理 \ref{thm:xi-foldbin-lastbit-tail-spectrum-single-jump-staircase} 的盒装恒等式与其逆向重写。证毕。
\end{proof}

\begin{theorem}[六倍窗 exact dark band 的半阶指数尺度]\label{thm:xi-time-part70b-dark-band-half-order-scale}
\leanverified{paper\_xi\_time\_part70b\_dark\_band\_half\_order\_scale}
沿用定理 \ref{thm:fold-zero-density-sparse} 与推论 \ref{cor:fold-zero-half-index-multiple6} 的记号，记
\[
\mathcal Z_m:=\{k\in \ZZ/(F_{m+2}\ZZ):\ \widehat d_m(k)=0\},
\]
并令 \(\tau_{\mathrm{div}}(n)\) 表示 \(n\) 的约数个数函数。若
\[
m+2\equiv 0\pmod 6,
\]
则有双边界
\[
F_{(m+2)/2}
\le
|\mathcal Z_m|
\le
\tau_{\mathrm{div}}(m+2)\,F_{(m+2)/2}.
\]
从而沿该六倍窗子序列，
\[
\log|\mathcal Z_m|
=
\frac{m}{2}\log\varphi+O(\log m),
\qquad
\frac{1}{m}\log|\mathcal Z_m|
\longrightarrow
\frac12\log\varphi.
\]
换言之，exact dark band 在六倍窗上具有 \(\varphi^{m/2}\) 级的半阶指数尺度。
\end{theorem}
\begin{proof}
下界正是推论 \ref{cor:fold-zero-half-index-multiple6} 的内容。上界则由定理 \ref{thm:fold-zero-density-sparse} 给出：
\[
|\mathcal Z_m|
\le
\tau_{\mathrm{div}}(m+2)\,F_{\lfloor (m+2)/2\rfloor}
=
\tau_{\mathrm{div}}(m+2)\,F_{(m+2)/2}.
\]
另一方面，Binet 公式给出
\[
\log F_n=n\log\varphi+O(1),
\]
而初等约数界
\[
\tau_{\mathrm{div}}(n)\le 2\sqrt n
\qquad (n\ge 1)
\]
蕴含
\[
\log\tau_{\mathrm{div}}(m+2)=O(\log m).
\]
把这些估计代入上下界，便得到
\[
\log|\mathcal Z_m|
=
\frac{m}{2}\log\varphi+O(\log m),
\]
并进一步推出极限
\[
\frac{1}{m}\log|\mathcal Z_m|\to \frac12\log\varphi.
\]
证毕。
\end{proof}

\noindent
上述结果把中心读出、escort 重加权、阈值层纹与 Fourier 侧零块压缩到同一条双折点--跳点主线上：在阈值/采样轴上，中心容量骨架只允许两个折点 \(\varphi^{-1}\) 与 \(1\)；在相位/路径信息轴上，温度参数 \(q\) 仅通过 \(\theta(q)\) 重新分配两支权重；在固定末位相位内，层纹由合法尾和集的单跳删除精确生成；而六倍窗上的 exact dark band 则在 Fourier 侧取得 \(\varphi^{m/2}\) 级的半阶指数尺度。因而，相应可视化中的连续着色纹理与局部波动若不改变这些离散骨架，便不能改变中心容量、尾和跳点结构与 dark-band 的主尺度。

\endinput


\paragraph{谱零点的高阶 \texorpdfstring{$\ZZ_2$}{Z2} 轨道抵消}
在谱零点的单坐标判别、余类并结构与六倍窗 half-order dark band 已经建立之后，还可进一步把单个按位翻转对合的相消推广为一个多坐标 \(\ZZ_2\)-轨道抵消机制。其结果是：一旦某个频率 \(k\) 触发若干独立坐标上的 \(-1\) 相位翻转，则整个 \(\Omega_m\) 会按自由的 \((\ZZ/2\ZZ)^r\)-轨道分块，而每个轨道上的相位和都逐块归零。

\begin{theorem}[谱零点的高阶 \texorpdfstring{$\ZZ_2$}{Z2} 轨道抵消]\label{thm:xi-time-part70c-spectrum-zero-z2-orbit-cancellation}
\leanverified{paper\_xi\_time\_part70c\_spectrum\_zero\_z2\_orbit\_cancellation}
沿用定理 \ref{thm:fold-spectrum-zero-criterion} 与推论 \ref{cor:fold-spectrum-zero-involution} 的记号，记
\[
M:=F_{m+2},
\qquad
\widehat d_m(k)=\sum_{\omega\in\Omega_m}e^{-2\pi i kN(\omega)/M}.
\]
对任意 \(k\in \ZZ/(M\ZZ)\)，定义
\[
J(k):=
\left\{
j\in\{1,\dots,m\}:\ 2kF_{j+1}\equiv M\pmod{2M}
\right\},
\]
并令
\[
\Phi_k(\omega):=e^{-2\pi i kN(\omega)/M}
\qquad
(\omega\in \Omega_m).
\]
对每个 \(j\in J(k)\)，记
\[
\iota_j:\Omega_m\to \Omega_m,
\qquad
\iota_j(\omega):=\omega\oplus e_j.
\]
再定义
\[
G_k:=\langle \iota_j:\ j\in J(k)\rangle.
\]
则成立：
\begin{enumerate}
  \item \(G_k\cong (\ZZ/2\ZZ)^{|J(k)|}\)，且其在 \(\Omega_m\) 上的作用自由；
  \item 对任意
  \[
  g=\prod_{j\in S}\iota_j\in G_k,
  \]
  都有
  \[
  \Phi_k(g\omega)=(-1)^{|S|}\Phi_k(\omega)
  \qquad
  (\omega\in\Omega_m);
  \]
  \item 若 \(J(k)\neq \varnothing\)，则对任意 \(G_k\)-轨道 \(\mathcal O\subset \Omega_m\)，都有
  \[
  \sum_{\omega\in\mathcal O}\Phi_k(\omega)=0;
  \]
  特别地，
  \[
  \widehat d_m(k)=0
  \qquad\Longleftrightarrow\qquad
  J(k)\neq \varnothing.
  \]
\end{enumerate}
\end{theorem}
\begin{proof}
由定理 \ref{thm:fold-spectrum-zero-criterion}，
\[
\widehat d_m(k)=0
\qquad\Longleftrightarrow\qquad
\exists\,j\in\{1,\dots,m\}\ \text{使}\ 2kF_{j+1}\equiv M\pmod{2M},
\]
这正等价于 \(J(k)\neq \varnothing\)。

对每个 \(j\in J(k)\)，推论 \ref{cor:fold-spectrum-zero-involution} 已给出：\(\iota_j\) 是 \(\Omega_m\) 上无固定点的对合，且
\[
\Phi_k(\iota_j\omega)=-\Phi_k(\omega).
\]
若 \(j,j'\in J(k)\) 且 \(j\neq j'\)，则 \(\iota_j,\iota_{j'}\) 只翻转不同坐标，因此彼此交换，且各自平方为恒等。于是
\[
G_k\cong (\ZZ/2\ZZ)^{|J(k)|}.
\]
任一非平凡元素 \(g=\prod_{j\in S}\iota_j\) 都会翻转至少一个坐标，故不可能固定任何 \(\omega\in\Omega_m\)；从而 \(G_k\) 在 \(\Omega_m\) 上自由作用。

再对
\[
g=\prod_{j\in S}\iota_j
\]
重复使用 \(\Phi_k(\iota_j\omega)=-\Phi_k(\omega)\)，便得
\[
\Phi_k(g\omega)=(-1)^{|S|}\Phi_k(\omega).
\]
因此当 \(J(k)\neq \varnothing\) 时，对任意轨道
\[
\mathcal O=G_k\cdot \omega
\]
都有
\[
\sum_{\eta\in\mathcal O}\Phi_k(\eta)
=
\Phi_k(\omega)\sum_{S\subseteq J(k)}(-1)^{|S|}
=
\Phi_k(\omega)(1-1)^{|J(k)|}
=0.
\]
把 \(\Omega_m\) 分解为 \(G_k\)-轨道的并，即得
\[
\widehat d_m(k)=\sum_{\omega\in\Omega_m}\Phi_k(\omega)=0.
\]
与第一步的零点判别合并，即得最后的充要条件。证毕。
\end{proof}

\endinput


\paragraph{六倍窗零点密度的精确指数率与 KL--最大纤维耦合}
在谱零点的高阶 \texorpdfstring{$(\ZZ_2)^r$}{(Z2)r} 轨道抵消已经给出逐点见证之后，还可继续沿频域--信息主线抽取出下列两条结果。它们分别把六倍窗子序列上的归一化零点密度锁定在严格的半维指数率 \(\varphi^{-m/2+o(m)}\)，并把输出分布相对于均匀分布的 KL 缺口与最大纤维超额压缩为同一条确定性控制不等式。

\begin{theorem}[六倍窗上的零点密度具有精确指数率]\label{thm:xi-time-part70c-zero-density-exact-ratio-exponent-multiple6}
\leanverified{paper\_xi\_time\_part70c\_zero\_density\_exact\_ratio\_exponent\_multiple6}
若
\[
m+2\equiv 0\pmod 6,
\qquad
d:=\frac{m+2}{2},
\qquad
M:=F_{m+2},
\qquad
\mathcal Z_m:=\{k\in \ZZ/(M\ZZ):\ \widehat d_m(k)=0\},
\]
则有极限
\[
\lim_{\substack{m\to\infty\\ m+2\equiv 0\ (6)}}
\frac1m\log\frac{|\mathcal Z_m|}{F_{m+2}}
=
-\frac12\log\varphi.
\]
等价地，在六倍窗子序列上，
\[
\frac{|\mathcal Z_m|}{F_{m+2}}
=
\varphi^{-m/2+o(m)}.
\]
\end{theorem}
\begin{proof}
下界方面，由推论 \ref{cor:fold-zero-half-index-multiple6}，
\[
|\mathcal Z_m|
\ge
F_d.
\]
因此
\[
\frac{|\mathcal Z_m|}{F_{m+2}}
\ge
\frac{F_d}{F_{m+2}}.
\]
由 Binet 公式 \(F_n=\varphi^{n+o(n)}\) 得
\[
\liminf_{\substack{m\to\infty\\ m+2\equiv 0\ (6)}}
\frac1m\log\frac{|\mathcal Z_m|}{F_{m+2}}
\ge
\lim_{m\to\infty}
\frac1m\log\frac{F_d}{F_{m+2}}
=
-\frac12\log\varphi.
\]

上界方面，由定理 \ref{thm:fold-zero-density-sparse}，
\[
\frac{|\mathcal Z_m|}{F_{m+2}}
\le
\tau(m+2)\frac{F_{\lfloor (m+2)/2\rfloor}}{F_{m+2}}.
\]
又由于
\[
\tau(n)=e^{o(n)},
\qquad
F_{\lfloor (m+2)/2\rfloor}= \varphi^{m/2+o(m)},
\qquad
F_{m+2}=\varphi^{m+o(m)},
\]
故
\[
\limsup_{\substack{m\to\infty\\ m+2\equiv 0\ (6)}}
\frac1m\log\frac{|\mathcal Z_m|}{F_{m+2}}
\le
-\frac12\log\varphi.
\]
上下界相合，即得结论。
\end{proof}

\begin{theorem}[最大纤维超额由 KL 缺口确定性控制]\label{thm:xi-time-part70c-kl-gap-controls-maxfiber-excess}
设
\[
p_m(r):=\frac{d_m(r)}{2^m},
\qquad
M:=F_{m+2},
\qquad
u(r):=\frac1M,
\]
并记
\[
M_m:=\max_r d_m(r),
\qquad
d_m^{\mathrm{avg}}:=\frac{2^m}{M}.
\]
则有确定性耦合不等式
\[
\boxed{\;
M_m-d_m^{\mathrm{avg}}
\le
2^m\,\|p_m-u\|_{\mathrm{TV}}
\le
2^m\sqrt{\frac12\,D_{\mathrm{KL}}(p_m\|u)}.\;}
\]
特别地，一旦 KL 缺口被控制，最大纤维相对于平均纤维的超额便被同步控制。
\leanverified{paper\_xi\_time\_part70c\_kl\_gap\_controls\_maxfiber\_excess}
\end{theorem}
\begin{proof}
首先由
\[
M_m=2^m\max_r p_m(r),
\qquad
d_m^{\mathrm{avg}}=\frac{2^m}{M}=2^m u(r),
\]
可得
\[
M_m-d_m^{\mathrm{avg}}
=
2^m\left(\max_r p_m(r)-\frac1M\right)
\le
2^m\max_r|p_m(r)-u(r)|.
\]
另一方面，对任意概率分布 \(p,u\) 都有
\[
\max_r|p(r)-u(r)|
\le
\frac12\sum_r|p(r)-u(r)|
=
\|p-u\|_{\mathrm{TV}},
\]
故
\[
M_m-d_m^{\mathrm{avg}}
\le
2^m\,\|p_m-u\|_{\mathrm{TV}}.
\]
最后应用 Pinsker 不等式
\[
\|p_m-u\|_{\mathrm{TV}}
\le
\sqrt{\frac12\,D_{\mathrm{KL}}(p_m\|u)},
\]
便得到所示结论。证毕。
\end{proof}

\noindent
由此，谱零点结构在三个层级上完成闭合：在单频层级，高阶 \texorpdfstring{$(\ZZ_2)^r$}{(Z2)r} 轨道抵消给出逐点见证；在密度层级，六倍窗上的归一化零点比例被锁定在 \(\varphi^{-m/2+o(m)}\)；在信息层级，KL 缺口与最大纤维超额被耦合为同一条确定性不等式。于是，轨道抵消、半维稀疏率与信息偏置便被压缩进同一条完全可审计的频域主线。

\endinput


\paragraph{末位终板刚性、临界预言机振荡与 Riesz 能量常数}
在末位相位的最终严格分层、bin-fold 规范群的最终满秩化、critical-scale 预言机容量骨架、escort 温度族的双折点响应以及六倍窗 exact dark band 的半阶指数尺度已经确立之后，还可把终板几何、商结构障碍、临界预算振荡与二阶能量常数进一步压缩为同一条后续链。在本段中统一记
\[
d_m(w):=d_m^{\mathrm{bin}}(w),
\qquad
\lambda:=\frac{2}{\varphi},
\qquad
\Omega_m:=\{0,1\}^m.
\]
下列结果分别刻画最小扇区的末位实现、大窗口上的全局非商化刚性、异常签名维数的 Fibonacci 饱和、临界比特率下的不可压平簇集带、escort 温度族的临界响应带、中心纤维场的精确二阶常数、碰撞概率与非平凡 Fourier 能量的精确主项，以及最大纤维超额的显式主系数。

\begin{theorem}[最小扇区的末位实现与终板逼近]\label{thm:xi-time-part70d-minsector-terminal-sheet-realization}
定义
\[
d_{\min}(m):=\min_{v\in X_m}d_m(v),
\qquad
A_1^{(m)}:=\{w\in X_m:\ w_m=1\},
\qquad
S_{m,\min}:=\{w\in X_m:\ d_m(w)=d_{\min}(m)\}.
\]
若“子覆盖同构”恒等式
\[
|S_{m,\min}|=|X_{m-2}|
\]
沿某个趋于无穷的偶数序列成立，则沿同一序列有
\[
\frac{|S_{m,\min}\triangle A_1^{(m)}|}{|X_m|}\longrightarrow 0.
\]
\end{theorem}
\leanverified{paper\_xi\_time\_part70d\_minsector\_terminal\_sheet\_realization}
\begin{proof}
由定理 \ref{thm:xi-time-part70ab-terminal-bit-eventual-strict-stratification}，对充分大的 \(m\) 有
\[
S_{m,\min}\subseteq A_1^{(m)}.
\]
另一方面，在假设下
\[
\frac{|S_{m,\min}|}{|X_m|}
=
\frac{|X_{m-2}|}{|X_m|}
=
\frac{F_m}{F_{m+2}}
\longrightarrow
\varphi^{-2},
\]
这也正是推论 \ref{cor:fold-bin-minsector-density-phi-minus2} 的内容。又由稳定词计数，
\[
|A_1^{(m)}|=F_m,
\qquad
|X_m|=F_{m+2},
\]
故
\[
\frac{|A_1^{(m)}|}{|X_m|}=\frac{F_m}{F_{m+2}}\longrightarrow \varphi^{-2}.
\]
由于 \(S_{m,\min}\subseteq A_1^{(m)}\) 且两者相对密度极限相同，遂有
\[
\frac{|A_1^{(m)}\setminus S_{m,\min}|}{|X_m|}\longrightarrow 0.
\]
而此时
\[
S_{m,\min}\triangle A_1^{(m)}=A_1^{(m)}\setminus S_{m,\min},
\]
故结论成立。
\end{proof}

\begin{theorem}[大窗口上的全局非商化刚性]\label{thm:xi-time-part70d-global-nonquotient-rigidity}
存在整数 \(m_0\)，使得对一切 \(m\ge m_0\)，等价关系
\[
\omega\sim_m\eta
\qquad\Longleftrightarrow\qquad
\Fold_m^{\mathrm{bin}}(\omega)=\Fold_m^{\mathrm{bin}}(\eta)
\]
既不能由任何有限群在 \(\Omega_m\) 上的自由作用之轨道划分实现，也不能来自任何有限阿贝尔群上的线性陪集划分。更准确地，不存在有限阿贝尔群 \(A_m\)、子群 \(H_m\le A_m\) 及双射
\[
\iota_m:\Omega_m\longrightarrow A_m
\]
使得
\[
\omega\sim_m\eta
\qquad\Longleftrightarrow\qquad
\iota_m(\omega)-\iota_m(\eta)\in H_m.
\]
\end{theorem}
\leanverified{paper\_xi\_time\_part70d\_global\_nonquotient\_rigidity}
\begin{proof}
记
\[
A_0^{(m)}:=\{w\in X_m:\ w_m=0\},
\qquad
A_1^{(m)}:=\{w\in X_m:\ w_m=1\}.
\]
由定理 \ref{thm:xi-time-part70ab-terminal-bit-eventual-strict-stratification}，存在 \(m_0\) 使得当 \(m\ge m_0\) 时，
\[
\min_{w\in A_0^{(m)}}d_m(w)>\max_{v\in A_1^{(m)}}d_m(v).
\]
由于 \(A_0^{(m)}\) 与 \(A_1^{(m)}\) 都非空，故对每个 \(m\ge m_0\)，纤维大小在同一窗口内至少取两个不同整数值。

若 \(\sim_m\) 来自某个有限群 \(G_m\) 的自由作用，则每个轨道大小都恒等于 \(|G_m|\)，因而所有纤维必同势；矛盾。

若 \(\sim_m\) 来自上式所述的线性陪集划分，则每个等价类都是某个陪集的原像，其大小恒等于 \(|H_m|\)，故所有纤维同样必同势；再度矛盾。于是两类实现均不可能存在。
\end{proof}

\begin{theorem}[异常签名维数的渐近饱和]\label{thm:xi-time-part70d-anomaly-signature-saturation}
\leanverified{paper\_xi\_time\_part70d\_anomaly\_signature\_saturation}
记
\[
G_m^{\mathrm{bin}}:=\prod_{w\in X_m}\mathfrak S_{d_m(w)}.
\]
则存在整数 \(m_0\)，使得对一切 \(m\ge m_0\)，都有
\[
\bigl(G_m^{\mathrm{bin}}\bigr)^{\mathrm{ab}}
\cong
(\ZZ/2\ZZ)^{|X_m|}
\cong
(\ZZ/2\ZZ)^{F_{m+2}}.
\]
因而任何保持“纤维独立置换”语义、并在 \(\FF_2\) 上完备分离 \(\bigl(G_m^{\mathrm{bin}}\bigr)^{\mathrm{ab}}\) 全部坐标的异常见证系统，其最小比特数恰为
\[
F_{m+2}.
\]
\end{theorem}
\begin{proof}
这正是定理 \ref{thm:xi-time-part65-binfold-center-vanishing-signature-stability} 的重述。该定理已证明：对充分大的 \(m\)，有
\[
\bigl(G_m^{\mathrm{bin}}\bigr)^{\mathrm{ab}}
\cong
(\ZZ/2\ZZ)^{F_{m+2}},
\qquad
\operatorname{rank}_{\FF_2}\!\bigl((G_m^{\mathrm{bin}})^{\mathrm{ab}}\bigr)=F_{m+2},
\]
并且纤维奇偶荷的最小完备秩在同一范围内严格等于 \(F_{m+2}\)。这恰好给出所述异常签名维数结论。
\end{proof}

\begin{theorem}[均匀微态输入下的临界比特率相变与簇集区间]\label{thm:xi-time-part70d-critical-bitrate-phase-cluster}
记
\[
\alpha:=\log_2\!\left(\frac{2}{\varphi}\right),
\qquad
\mathrm{Succ}_m(B):=\mathrm{Succ}_m^{\mathrm{bin}}(B).
\]
则成立：
\begin{enumerate}
  \item 若 \(\rho<\alpha\)，则
  \[
  \mathrm{Succ}_m(\lfloor \rho m\rfloor)\longrightarrow 0.
  \]
  \item 若 \(\rho>\alpha\)，则
  \[
  \mathrm{Succ}_m(\lfloor \rho m\rfloor)\longrightarrow 1.
  \]
  \item 对任意固定 \(\beta\in\RR\)，令
  \[
  B_m:=\lfloor \alpha m+\beta\rfloor.
  \]
  则序列 \(\bigl(\mathrm{Succ}_m(B_m)\bigr)_{m\ge 1}\) 不收敛，其全部簇集点恰组成闭区间
  \[
  \left[\frac{\varphi^2}{2\sqrt5},\,1\right].
  \]
\end{enumerate}
\end{theorem}
\begin{proof}
记 \(\lambda:=2/\varphi\)，则 \(\alpha=\log_2\lambda\)。

若 \(\rho<\alpha\)，则由平凡上界
\[
C_m^{\mathrm{bin}}(B)\le 2^B|X_m|
\]
得到
\[
\mathrm{Succ}_m(\lfloor \rho m\rfloor)
\le
\frac{2^{\lfloor \rho m\rfloor}F_{m+2}}{2^m}
=
\frac{2^{\lfloor \rho m\rfloor}}{\lambda^m}\cdot \frac{F_{m+2}}{\varphi^m}
\longrightarrow 0,
\]
因为 \(2^{\lfloor \rho m\rfloor}/\lambda^m\to 0\) 且 \(F_{m+2}/\varphi^m\to \varphi^2/\sqrt5\)。

若 \(\rho>\alpha\)，则 \(2^{\lfloor \rho m\rfloor}/\lambda^m\to+\infty\)。由推论 \ref{cor:xi-time-part70ab-extremal-fibers-terminal-bit-localization}，
\[
M_m:=\max_{w\in X_m}d_m(w)=\lambda^m+O(1).
\]
故对充分大的 \(m\) 有 \(2^{\lfloor \rho m\rfloor}>M_m\)，从而所有纤维都完全通过，即
\[
\mathrm{Succ}_m(\lfloor \rho m\rfloor)=1.
\]

现取临界序列 \(B_m=\lfloor \alpha m+\beta\rfloor\)。定义
\[
\theta_m:=\alpha m+\beta-\lfloor \alpha m+\beta\rfloor\in[0,1),
\qquad
\tau_m:=\frac{2^{B_m}}{\lambda^m}=2^{-\theta_m}\in[1/2,1].
\]
若 \(\log_2\varphi\in\QQ\)，则存在整数 \(p,q\ge 1\) 使 \(\varphi^q=2^p\in\QQ\)，这与 \(\varphi^q\notin\QQ\) 矛盾。故 \(\alpha=1-\log_2\varphi\) 为无理数，从而 \((\theta_m)\) 在 \([0,1]\) 上稠密，\((\tau_m)\) 的簇集点恰为整个区间 \([1/2,1]\)。

任取 \(\tau\in[1/2,1]\) 并取子列 \(m_j\) 使 \(\tau_{m_j}\to\tau\)。由定理 \ref{thm:xi-foldbin-dyadic-capacity-critical-window-limit}，
\[
\mathrm{Succ}_{m_j}(B_{m_j})
\longrightarrow
\Psi(\tau),
\]
其中
\[
\Psi(\tau)=
\begin{cases}
\dfrac{\varphi^2}{\sqrt5}\,\tau, & 0\le \tau\le \varphi^{-1},\\[6pt]
\dfrac{\varphi}{\sqrt5}\,\tau+\dfrac{1}{\varphi\sqrt5}, & \varphi^{-1}\le \tau\le 1,\\[8pt]
1, & \tau\ge 1.
\end{cases}
\]
因此 \(\bigl(\mathrm{Succ}_m(B_m)\bigr)\) 的全部簇集点恰为连续单调函数 \(\Psi\) 在 \([1/2,1]\) 上的像，即
\[
\Psi([1/2,1])
=
\left[\Psi(1/2),\,\Psi(1)\right]
=
\left[\frac{\varphi^2}{2\sqrt5},\,1\right].
\]
该区间非单点，故序列不收敛。
\end{proof}

\begin{theorem}[escort 温度族下的临界响应带]\label{thm:xi-time-part70d-escort-critical-response-band}
对任意 \(q\ge 0\)，定义 escort 加权的阈值响应
\[
\mathrm{Succ}_{m}^{(q)}(B)
:=
\sum_{w\in X_m}\pi_{m,q}(w)\min\!\left(1,\frac{2^B}{d_m(w)}\right),
\qquad
\theta(q):=\frac{1}{1+\varphi^{q+1}}.
\]
则同一临界比特率
\[
\alpha=\log_2\!\left(\frac{2}{\varphi}\right)
\]
给出 sharp transition：
\begin{enumerate}
  \item 若 \(\rho<\alpha\)，则
  \[
  \mathrm{Succ}_{m}^{(q)}(\lfloor \rho m\rfloor)\longrightarrow 0.
  \]
  \item 若 \(\rho>\alpha\)，则
  \[
  \mathrm{Succ}_{m}^{(q)}(\lfloor \rho m\rfloor)\longrightarrow 1.
  \]
  \item 若对固定 \(\beta\in\RR\) 令
  \[
  B_m=\lfloor \alpha m+\beta\rfloor,
  \]
  则序列 \(\bigl(\mathrm{Succ}_{m}^{(q)}(B_m)\bigr)_{m\ge 1}\) 不收敛，其全部簇集点恰为
  \[
  \left[\frac{1+(\varphi-1)\theta(q)}{2},\,1\right].
  \]
\end{enumerate}
\end{theorem}
\begin{proof}
若 \(\rho<\alpha\)，则 \(2^{\lfloor \rho m\rfloor}/\lambda^m\to 0\)。由定理 \ref{thm:fold-bin-two-state-asymptotic}，存在常数 \(C>0\) 使对充分大的 \(m\) 与全部 \(w\in X_m\) 都有
\[
d_m(w)\ge \lambda^m\left(\varphi^{-1}-C\Bigl(\frac{\varphi}{2}\Bigr)^m\right).
\]
于是
\[
\mathrm{Succ}_{m}^{(q)}(\lfloor \rho m\rfloor)
\le
\frac{2^{\lfloor \rho m\rfloor}}{\lambda^m\left(\varphi^{-1}-C(\varphi/2)^m\right)}
\longrightarrow 0.
\]

若 \(\rho>\alpha\)，则记
\[
M_m:=\max_{w\in X_m}d_m(w).
\]
由推论 \ref{cor:xi-time-part70ab-extremal-fibers-terminal-bit-localization}，\(M_m=\lambda^m+O(1)\)。故对充分大的 \(m\) 有 \(2^{\lfloor \rho m\rfloor}>M_m\)，于是每个纤维都完全通过，从而
\[
\mathrm{Succ}_{m}^{(q)}(\lfloor \rho m\rfloor)=1.
\]

现令 \(B_m=\lfloor \alpha m+\beta\rfloor\)，并仍记
\[
\tau_m:=\frac{2^{B_m}}{\lambda^m}\in[1/2,1].
\]
由上一条定理的证明，\((\tau_m)\) 的全部簇集点恰为 \([1/2,1]\)。

记
\[
I:=[1/2,1],
\qquad
J:=\left[\frac{\varphi^{-1}}{2},\,2\right].
\]
函数
\[
f(\tau,x):=\min\!\left(1,\frac{\tau}{x}\right)
\]
在紧集 \(I\times J\) 上联合 Lipschitz。再由定理 \ref{thm:fold-bin-two-state-asymptotic}，
\[
\left|\frac{d_m(w)}{\lambda^m}-\varphi^{-w_m}\right|
\le
C\Bigl(\frac{\varphi}{2}\Bigr)^m
\qquad (w\in X_m),
\]
故对充分大的 \(m\) 与全部 \(w\in X_m\) 都有
\[
\min\!\left(1,\frac{2^{B_m}}{d_m(w)}\right)
=
f\!\left(\tau_m,\frac{d_m(w)}{\lambda^m}\right)
=
f\!\left(\tau_m,\varphi^{-w_m}\right)
+O\!\left(\left(\frac{\varphi}{2}\right)^m\right),
\]
且误差对 \(w\) 与 \(\tau_m\in I\) 一致。对 \(\pi_{m,q}\) 取期望，并用命题 \ref{prop:fold-bin-escort-lastbit}
\[
\mathbf P_{\pi_{m,q}}(w_m=1)=\theta(q)+O\!\left(\left(\frac{\varphi}{2}\right)^m\right),
\qquad
\mathbf P_{\pi_{m,q}}(w_m=0)=1-\theta(q)+O\!\left(\left(\frac{\varphi}{2}\right)^m\right),
\]
得到
\[
\mathrm{Succ}_{m}^{(q)}(B_m)
=
\bigl(1-\theta(q)\bigr)\min\{1,\tau_m\}
+
\theta(q)\min\{1,\varphi\tau_m\}
+o(1).
\]
任取 \(\tau\in[1/2,1]\) 及子列 \(m_j\) 使 \(\tau_{m_j}\to\tau\)，便有
\[
\mathrm{Succ}_{m_j}^{(q)}(B_{m_j})
\longrightarrow
\Psi_q^{\mathrm{esc}}(\tau),
\qquad
\Psi_q^{\mathrm{esc}}(\tau)
:=
\bigl(1-\theta(q)\bigr)\min\{1,\tau\}
+
\theta(q)\min\{1,\varphi\tau\}.
\]
故全部簇集点恰为 \(\Psi_q^{\mathrm{esc}}([1/2,1])\)。由于 \(\Psi_q^{\mathrm{esc}}\) 在 \([1/2,1]\) 上连续单调，故该像为区间
\[
\left[\Psi_q^{\mathrm{esc}}(1/2),\,\Psi_q^{\mathrm{esc}}(1)\right].
\]
又因 \(1/2<\varphi^{-1}\)，故
\[
\Psi_q^{\mathrm{esc}}(1/2)
=
\frac{1-\theta(q)}{2}+\frac{\varphi\theta(q)}{2}
=
\frac{1+(\varphi-1)\theta(q)}{2},
\qquad
\Psi_q^{\mathrm{esc}}(1)=1.
\]
于是所求簇集区间即为
\[
\left[\frac{1+(\varphi-1)\theta(q)}{2},\,1\right].
\]
其长度严格为正，故序列不收敛。
\end{proof}

\begin{theorem}[中心纤维场的精确二阶常数]\label{thm:xi-time-part70d-central-fiber-second-moment}
\leanverified{paper\_xi\_time\_part70d\_central\_fiber\_second\_moment}
令 \(U_m\sim \mathrm{Unif}(X_m)\)，并定义尺度化纤维场
\[
Y_m:=\left(\frac{\varphi}{2}\right)^m d_m(U_m).
\]
则有极限
\[
\mathbb E(Y_m^2)\longrightarrow 2\varphi^{-2},
\qquad
\operatorname{Var}(Y_m)\longrightarrow \varphi^{-7}.
\]
\end{theorem}
\begin{proof}
由推论 \ref{cor:xi-time-part70a-scaled-moments-leading-constant}，当 \(q=1,2\) 时分别有
\[
\mathbb E(Y_m)
\longrightarrow
\frac{1}{\varphi}+\varphi^{-3},
\qquad
\mathbb E(Y_m^2)
\longrightarrow
\frac{1}{\varphi}+\varphi^{-4}.
\]
利用恒等式 \(\varphi^2=\varphi+1\) 与 \(\sqrt5=2\varphi-1\)，可将上式化简为
\[
\frac{1}{\varphi}+\varphi^{-4}=2\varphi^{-2},
\qquad
\frac{1}{\varphi}+\varphi^{-3}=\frac{\sqrt5}{\varphi^2}.
\]
因此
\[
\operatorname{Var}(Y_m)
=
\mathbb E(Y_m^2)-\bigl(\mathbb E(Y_m)\bigr)^2
\longrightarrow
2\varphi^{-2}-\frac{5}{\varphi^4}.
\]
再用 \(\varphi^2=\varphi+1\) 化简，得到
\[
2\varphi^{-2}-\frac{5}{\varphi^4}=\varphi^{-7}.
\]
证毕。
\end{proof}

\begin{theorem}[碰撞概率与非平凡 Fourier 能量的精确主项]\label{thm:xi-time-part70d-collision-riesz-main-terms}
沿用定理 \ref{thm:xi-time-part70d-central-fiber-second-moment} 的记号，并沿用定理 \ref{thm:xi-time-part9p-chi2-ranktwo-resonant-floor} 的 Fourier 记号。记
\[
\mathsf{Col}_m:=
\frac{1}{4^m}\sum_{w\in X_m}d_m(w)^2.
\]
则有
\[
\mathsf{Col}_m\sim \frac{2}{\sqrt5}\,\varphi^{-m},
\qquad
F_{m+2}\,\mathsf{Col}_m\longrightarrow \frac{2\varphi^2}{5}.
\]
进一步，非平凡 Fourier 能量满足
\[
4^{-m}\sum_{k\ne 0}\bigl|\widehat d_m(k)\bigr|^2
\longrightarrow
\frac{1}{5\varphi^3}.
\]
\end{theorem}
\begin{proof}
由 \(U_m\sim\mathrm{Unif}(X_m)\) 的定义，
\[
\mathsf{Col}_m
=
\frac{|X_m|}{4^m}\,\mathbb E\!\bigl(d_m(U_m)^2\bigr).
\]
又因为
\[
d_m(U_m)=\lambda^m Y_m,
\qquad
\lambda=\frac{2}{\varphi},
\]
故
\[
\mathsf{Col}_m
=
|X_m|\,\varphi^{-2m}\,\mathbb E(Y_m^2).
\]
由 \(|X_m|=F_{m+2}\) 与定理 \ref{thm:xi-time-part70d-central-fiber-second-moment}，
\[
\mathsf{Col}_m
=
F_{m+2}\varphi^{-2m}\bigl(2\varphi^{-2}+o(1)\bigr).
\]
再用 Binet 公式
\[
F_{m+2}\sim \frac{\varphi^{m+2}}{\sqrt5},
\]
便得
\[
\mathsf{Col}_m
\sim
\frac{\varphi^{m+2}}{\sqrt5}\cdot \varphi^{-2m}\cdot 2\varphi^{-2}
=
\frac{2}{\sqrt5}\,\varphi^{-m}.
\]
两边再乘以 \(F_{m+2}\) 即得
\[
F_{m+2}\,\mathsf{Col}_m\longrightarrow \frac{2\varphi^2}{5}.
\]

另一方面，由定理 \ref{thm:xi-time-part9p-chi2-ranktwo-resonant-floor} 的第一条精确恒等式，
\[
F_{m+2}\,\mathsf{Col}_m-1
=
4^{-m}\sum_{k\ne 0}\bigl|\widehat d_m(k)\bigr|^2.
\]
取极限并用
\[
\frac{2\varphi^2}{5}-1=\frac{1}{5\varphi^3},
\]
便得到
\[
4^{-m}\sum_{k\ne 0}\bigl|\widehat d_m(k)\bigr|^2
\longrightarrow
\frac{1}{5\varphi^3}.
\]
证毕。
\end{proof}

\begin{corollary}[最大纤维超额的精确主项]\label{cor:xi-time-part70d-maxfiber-excess-exact-main-term}
记
\[
\bar d_m:=\frac{2^m}{F_{m+2}},
\qquad
M_m:=\max_{w\in X_m}d_m(w).
\]
则有
\[
M_m-\bar d_m
=
\varphi^{-4}\left(\frac{2}{\varphi}\right)^m
+o\!\left(\left(\frac{2}{\varphi}\right)^m\right).
\]
\end{corollary}
\begin{proof}
由推论 \ref{cor:xi-time-part70ab-extremal-fibers-terminal-bit-localization}，
\[
M_m=\left(\frac{2}{\varphi}\right)^m+O(1).
\]
另一方面，由 Binet 公式，
\[
\bar d_m
=
\frac{2^m}{F_{m+2}}
=
\frac{\sqrt5}{\varphi^2}\left(\frac{2}{\varphi}\right)^m
+o\!\left(\left(\frac{2}{\varphi}\right)^m\right).
\]
由于 \(O(1)=o((2/\varphi)^m)\)，相减即得
\[
M_m-\bar d_m
=
\left(1-\frac{\sqrt5}{\varphi^2}\right)\left(\frac{2}{\varphi}\right)^m
+o\!\left(\left(\frac{2}{\varphi}\right)^m\right).
\]
再用 \(\sqrt5=2\varphi-1\) 与 \(\varphi-1=\varphi^{-1}\) 化简：
\[
1-\frac{\sqrt5}{\varphi^2}
=
1-\frac{2\varphi-1}{\varphi^2}
=
\frac{(\varphi-1)^2}{\varphi^2}
=
\varphi^{-4}.
\]
故结论成立。
\end{proof}

\noindent
由此，末位 \(1\) 终板不再只是最小纤维的容纳层，而在“子覆盖同构”外推下渐近地与整个最小扇区重合；与此同时，充分大窗口上的折叠关系也被彻底排除出自由商与线性陪集两类等势模型。再向临界资源侧推进，dyadic 预算与 escort 温度都在同一比特率 \(\log_2(2/\varphi)\) 处出现不可压平的簇集带，而其下边界由黄金比例与温度参数共同给出显式常数。最后，中心纤维场的二阶矩、碰撞概率、非平凡 Fourier 能量与最大纤维超额都锁定到闭式黄金常数，从而终板刚性、资源相变与频域能量三条先前分立的链被压缩为同一条可审计的 Fibonacci 后续主线。

\endinput

\paragraph{window-\texorpdfstring{$6$}{6} 抽象轨道型的最小循环实现与几何障碍的范畴分离}
在 window-\(6\) 的纤维多值性、自由作用障碍与线性陪集障碍都已确立之后，还可进一步把“抽象轨道大小多重集”与“真实立方体上的几何实现”彻底分开：就纯 \(G\)-集合而言，直方图 \(2^8,3^4,4^9\) 实际上已经能由最小的 \(12\) 阶循环群实现；因此 window-\(6\) 的真正障碍并不位于 Burnside 型轨道计数层，而位于把该轨道型忠实落到 \(\Omega_6\) 上时所必须满足的自由性与线性几何兼容条件。

\leanverified{paper\_xi\_time\_part70da\_window6\_abstract\_orbittype\_minimal\_c12}
\begin{theorem}[window-\texorpdfstring{$6$}{6} 的抽象轨道型可由最小阶循环群 \texorpdfstring{$C_{12}$}{C12} 实现]\label{thm:xi-time-part70da-window6-abstract-orbittype-minimal-c12}
设
\[
\mathfrak m:=2^8\,3^4\,4^9
\]
表示轨道大小多重集，其中大小为 \(2,3,4\) 的轨道个数分别为 \(8,4,9\)。则：
\begin{enumerate}
  \item 若某有限群 \(G\) 的某个有限 \(G\)-集合具有轨道型 \(\mathfrak m\)，则必有
  \[
  12\mid |G|.
  \]
  \item 该下界是锐的。更精确地，循环群 \(C_{12}\) 已足够：若对每个 \(d\mid 12\) 记 \(C_d\le C_{12}\) 为唯一的阶为 \(d\) 的子群，则
  \[
  Y:=8(C_{12}/C_6)\ \sqcup\ 4(C_{12}/C_4)\ \sqcup\ 9(C_{12}/C_3)
  \]
  是一个 \(64\) 点 \(C_{12}\)-集合，其轨道大小多重集恰为 \(\mathfrak m\)。
\end{enumerate}
因此，就抽象轨道算术而言，实现 window-\(6\) 直方图所需的最小群阶恰为
\[
12.
\]
\end{theorem}
\begin{proof}
若有限群 \(G\) 的某个轨道分解中出现大小为 \(2,3,4\) 的轨道，则每个轨道大小都必须整除 \(|G|\)。因此
\[
2\mid |G|,
\qquad
3\mid |G|,
\qquad
4\mid |G|,
\]
从而
\[
\operatorname{lcm}(2,3,4)=12\mid |G|.
\]
这就给出第一条与最小群阶的下界。

现取 \(G=C_{12}\)。由于循环群对每个除数都只有唯一子群，故存在唯一的阶 \(6,4,3\) 子群 \(C_6,C_4,C_3\)。相应陪集空间的基数分别为
\[
|C_{12}/C_6|=2,
\qquad
|C_{12}/C_4|=3,
\qquad
|C_{12}/C_3|=4.
\]
因此上式定义的离散并
\[
Y:=8(C_{12}/C_6)\ \sqcup\ 4(C_{12}/C_4)\ \sqcup\ 9(C_{12}/C_3)
\]
确为一个 \(C_{12}\)-集合，并且其轨道大小多重集恰为
\[
2^8\,3^4\,4^9.
\]
总点数为
\[
8\cdot 2+4\cdot 3+9\cdot 4=16+12+36=64.
\]
故 \(C_{12}\) 已实现该抽象轨道型。结合上界即知最小群阶恰为 \(12\)。证毕。
\end{proof}

\begin{corollary}[window-\texorpdfstring{$6$}{6} 的压缩障碍是几何性的，而非纯轨道计数性的]\label{cor:xi-time-part70da-window6-obstruction-geometric-not-burnside}
对 \(\Omega_6=\{0,1\}^6\) 上的 bin-fold 纤维划分，成立如下分层：
\begin{enumerate}
  \item 在抽象有限 \(G\)-集合范畴中，轨道型
  \[
  2^8\,3^4\,4^9
  \]
  已由定理 \ref{thm:xi-time-part70da-window6-abstract-orbittype-minimal-c12} 中的 \(C_{12}\)-集合实现。
  \item 在真实立方体 \(\Omega_6\) 上，该纤维划分不可能来自任何有限群的自由作用轨道分解。
  \item 同样地，该纤维划分也不可能来自任何有限阿贝尔群上的线性陪集划分。
\end{enumerate}
因此，window-\(6\) 的压缩障碍并不是 Burnside 型轨道大小不可实现性，而是一个严格更强的几何--线性兼容性障碍。
\end{corollary}
\begin{proof}
第一条正是定理 \ref{thm:xi-time-part70da-window6-abstract-orbittype-minimal-c12}。第二、三条则由推论 \ref{cor:foldbin6-degeneracy-multivalued-nonfree} 直接给出：window-\(6\) 的纤维直方图
\[
2:8,\qquad 3:4,\qquad 4:9
\]
不能来自任何自由有限群作用，也不能来自任何有限阿贝尔群上的线性陪集划分。将三条结论并列，便得到所述范畴分层。证毕。
\end{proof}

\endinput

\paragraph{整振幅、仿射压力与 Stokes 归一化的单核闭合}
在稳定层两态一致渐近、bin-fold 退化度的全实矩闭式、相对稳定层均匀基线的 Rényi 散度极限，以及 fold-gauge 常数的 Stirling--Bernoulli 奇层展开均已建立之后，还可把这些此前分散出现的接口压缩为一条复参数解析骨架：同一个整函数振幅同时控制 bin-fold 的复矩主项、输出分布的全轴压力、Rényi 极限曲线的规范延拓以及 \(\log C_m\) 尾部的 Bernoulli 采样。由此，一切归一化无关的有限窗对数观测都退化为该整振幅在负奇整数点上的读出，而剩余自由度严格只保留一个乘法常数。

\begin{theorem}[bin-fold 全复矩谱的整振幅律]\label{thm:xi-time-part70e-binfold-entire-amplitude-law}
\leanverified{paper\_xi\_time\_part70e\_binfold\_entire\_amplitude\_law}
记
\[
\varrho:=\frac{\varphi}{2}\in(0,1),
\qquad
d_m(w):=d_m^{\mathrm{bin}}(w),
\qquad
S_m(z):=\sum_{w\in X_m} d_m(w)^z
\qquad (z\in\CC),
\]
其中 \(d_m(w)^z:=\exp(z\log d_m(w))\) 取实正数上的主支。则对任意紧集 \(K\subset\CC\)，都有一致渐近
\[
S_m(z)
=
\Bigl(\frac{2}{\varphi}\Bigr)^{zm}
\Bigl(F_{m+1}+\varphi^{-z}F_m\Bigr)
\Bigl(1+O_K(\varrho^m)\Bigr)
\qquad (z\in K).
\]
从而
\[
\mathcal A_m(z):=
\varphi^{-m}\Bigl(\frac{\varphi}{2}\Bigr)^{zm}S_m(z)
\]
在 \(\CC\) 的每个紧集上一致收敛到
\[
\mathcal A(z):=\frac{\varphi+\varphi^{-z}}{\sqrt5},
\]
因此 \(\mathcal A\) 为整函数。
\end{theorem}
\begin{proof}
由定理 \ref{thm:fold-bin-two-state-asymptotic}，存在 \(\varepsilon_m(w)\) 满足
\[
d_m(w)=\varphi^{-w_m}\Bigl(\frac{2}{\varphi}\Bigr)^m\bigl(1+\varepsilon_m(w)\bigr),
\qquad
\sup_{w\in X_m}|\varepsilon_m(w)|=O(\varrho^m).
\]
固定紧集 \(K\subset\CC\)。对充分大的 \(m\) 有 \(\sup_{w}|\varepsilon_m(w)|<1/2\)，因而
\[
\log\bigl(1+\varepsilon_m(w)\bigr)=O(\varrho^m)
\]
在 \(w\in X_m\) 上一致成立。由于 \(K\) 有界，遂有
\[
(1+\varepsilon_m(w))^z
=
\exp\!\bigl(z\log(1+\varepsilon_m(w))\bigr)
=
1+O_K(\varrho^m)
\]
在 \(z\in K\) 与 \(w\in X_m\) 上一致成立。于是
\[
d_m(w)^z
=
\Bigl(\frac{2}{\varphi}\Bigr)^{zm}\varphi^{-zw_m}
\Bigl(1+O_K(\varrho^m)\Bigr).
\]
把 \(X_m\) 按末位分拆为
\[
A_0^{(m)}:=\{w\in X_m:w_m=0\},
\qquad
A_1^{(m)}:=\{w\in X_m:w_m=1\},
\]
并用
\[
|A_0^{(m)}|=F_{m+1},
\qquad
|A_1^{(m)}|=F_m,
\]
便得到
\[
S_m(z)
=
\Bigl(\frac{2}{\varphi}\Bigr)^{zm}
\Bigl(F_{m+1}+\varphi^{-z}F_m\Bigr)
\Bigl(1+O_K(\varrho^m)\Bigr).
\]
再乘以 \(\varphi^{-m}(\varphi/2)^{zm}\)，并用 Binet 渐近
\[
\varphi^{-m}F_{m+1}\to \frac{\varphi}{\sqrt5},
\qquad
\varphi^{-m}F_m\to \frac{1}{\sqrt5},
\]
即可得到
\[
\mathcal A_m(z)\to \frac{\varphi+\varphi^{-z}}{\sqrt5}
\]
且收敛在每个紧集上一致。极限函数由显式公式显然为整函数。证毕。
\end{proof}

\begin{corollary}[全实轴仿射压力与单谱塌缩]\label{cor:xi-time-part70e-affine-pressure-monofractal}
\leanverified{paper\_xi\_time\_part70e\_affine\_pressure\_monofractal}
记
\[
\mu_m(w):=\frac{d_m^{\mathrm{bin}}(w)}{2^m},
\qquad
\Xi_m(t):=\sum_{w\in X_m}\mu_m(w)^t
\qquad (t\in\RR).
\]
则极限
\[
\tau_\mu(t):=\lim_{m\to\infty}\frac1m\log \Xi_m(t)
\]
对全部 \(t\in\RR\) 存在，并且满足闭式
\[
\tau_\mu(t)=(1-t)\log\varphi.
\]
因而 \(\tau_\mu\) 在全实轴上仿射，其 Legendre--Fenchel 变换退化为单点谱。等价地，局部信息密度满足一致塌缩
\[
-\frac1m\log \mu_m(w)=\log\varphi+O(m^{-1})
\qquad (w\in X_m),
\]
其中误差在 \(w\) 上一致；并且相对于
\[
\log|X_m|=(m+2)\log\varphi-\frac12\log 5+o(1)
\]
的归一化局部维数亦有
\[
-\frac{\log\mu_m(w)}{\log|X_m|}=1+O(m^{-1})
\qquad (w\in X_m),
\]
其中误差同样在 \(w\) 上一致。
\end{corollary}
\begin{proof}
由定理 \ref{thm:xi-time-part70e-binfold-entire-amplitude-law} 在实轴上的结论，
\[
\Xi_m(t)=2^{-tm}S_m(t)
=
\varphi^{-tm}\Bigl(F_{m+1}+\varphi^{-t}F_m\Bigr)\Bigl(1+O_t(\varrho^m)\Bigr).
\]
又 \(F_{m+1}+\varphi^{-t}F_m=\Theta_t(\varphi^m)\)，故
\[
\frac1m\log\Xi_m(t)\to (1-t)\log\varphi.
\]
这就得到 \(\tau_\mu(t)\) 的显式公式。由于 \(\tau_\mu\) 在全轴上仿射，其 Legendre--Fenchel 变换只能在唯一斜率处非平凡，故相应多重分形谱退化为单点。

另一方面，由定理 \ref{thm:fold-bin-two-state-asymptotic}，
\[
\mu_m(w)=\frac{d_m^{\mathrm{bin}}(w)}{2^m}
=
\varphi^{-m-w_m}\bigl(1+O(\varrho^m)\bigr)
\]
在 \(w\in X_m\) 上一致成立。因此
\[
-\log\mu_m(w)=m\log\varphi+w_m\log\varphi+O(\varrho^m)
\]
一致成立，两边除以 \(m\) 即得
\[
-\frac1m\log\mu_m(w)=\log\varphi+O(m^{-1})
\]
在 \(w\) 上一致。再用 \(|X_m|=F_{m+2}\) 的 Binet 渐近
\[
\log|X_m|=(m+2)\log\varphi-\frac12\log5+o(1),
\]
即可得到最后一式。证毕。
\end{proof}

\begin{theorem}[规范化 Rényi 极限曲线与负奇点 Bernoulli 采样]\label{thm:xi-time-part70e-renyi-curve-negative-odd-sampling}
\leanverified{paper\_xi\_time\_part70e\_renyi\_curve\_negative\_odd\_sampling}
对 \(\alpha>0\)，记 \(D_\alpha^\infty\) 为定理 \ref{thm:fold-bin-renyi-divergence-limit} 中的极限 Rényi 散度，并定义规范化曲线
\[
\mathcal R(\alpha):=
5^{(\alpha-1)/2}\varphi^{2-2\alpha}
\exp\!\bigl((\alpha-1)D_\alpha^\infty\bigr).
\]
则对全部 \(\alpha>0\) 都有严格恒等式
\[
\mathcal R(\alpha)=\frac{\varphi+\varphi^{-\alpha}}{\sqrt5}=\mathcal A(\alpha),
\]
其中 \(\mathcal A\) 为定理 \ref{thm:xi-time-part70e-binfold-entire-amplitude-law} 的整振幅。因而 \(\mathcal R\) 由正实轴数据唯一延拓为整函数，并与 \(\mathcal A\) 完全一致。进一步，对任意固定整数 \(R\ge 1\)，有
\[
\log C_m-\log(2\pi)
=
\frac{\sqrt5}{\varphi^2}
\sum_{r=1}^{R}
\frac{B_{2r}}{r(2r-1)}
\mathcal R\bigl(-(2r-1)\bigr)\,
\varrho^{(2r-1)m}
+
O\!\bigl(\varrho^{(2R+1)m}\bigr).
\]
换言之，fold-gauge 常数的尾部模态恰由规范化 Rényi 极限曲线在负奇整数点的取值经 Bernoulli 变换生成。
\end{theorem}
\begin{proof}
由定理 \ref{thm:fold-bin-renyi-divergence-limit}，当 \(\alpha\neq 1\) 时有
\[
D_\alpha^\infty
=
\frac{(2\alpha-2)\log\varphi+\log(\varphi+\varphi^{-\alpha})-\alpha\log\sqrt5}{\alpha-1}.
\]
故
\[
\exp\!\bigl((\alpha-1)D_\alpha^\infty\bigr)
=
\frac{\varphi^{2\alpha-2}}{5^{\alpha/2}}
\bigl(\varphi+\varphi^{-\alpha}\bigr)
=
\frac{\varphi^{2\alpha-2}}{5^{(\alpha-1)/2}}
\cdot
\frac{\varphi+\varphi^{-\alpha}}{\sqrt5}.
\]
两边再乘以 \(5^{(\alpha-1)/2}\varphi^{2-2\alpha}\)，便得到
\[
\mathcal R(\alpha)=\frac{\varphi+\varphi^{-\alpha}}{\sqrt5}.
\]
而当 \(\alpha=1\) 时，左端按定义等于 \(1\)，右端则满足
\[
\frac{\varphi+\varphi^{-1}}{\sqrt5}=1,
\]
故同样成立。于是 \(\mathcal R\) 在正实轴上与显式函数 \((\varphi+\varphi^{-z})/\sqrt5\) 一致，因此其唯一整延拓正是定理 \ref{thm:xi-time-part70e-binfold-entire-amplitude-law} 中的 \(\mathcal A\)。

另一方面，由定理 \ref{thm:fold-bin-gauge-constant-stirling-bernoulli-hierarchy}，
\[
\log C_m-\log(2\pi)
=
\sum_{r=1}^{R}
\frac{B_{2r}}{r(2r-1)}
\bigl(\varphi^{-1}+\varphi^{2r-3}\bigr)\varrho^{(2r-1)m}
+O\!\bigl(\varrho^{(2R+1)m}\bigr).
\]
而由刚刚得到的闭式，对每个 \(r\ge 1\) 都有
\[
\frac{\sqrt5}{\varphi^2}\mathcal R\bigl(-(2r-1)\bigr)
=
\frac{\sqrt5}{\varphi^2}\cdot
\frac{\varphi+\varphi^{2r-1}}{\sqrt5}
=
\varphi^{-1}+\varphi^{2r-3}.
\]
代回即可得到所示采样公式。证毕。
\end{proof}

\begin{theorem}[零和有限窗对数观测的热力学全决定性]\label{thm:xi-time-part70e-zero-sum-window-thermo-determinacy}
\leanverified{paper\_xi\_time\_part70e\_zero\_sum\_window\_thermo\_determinacy}
令移位算子 \(T\) 作用于序列 \(f=(f_m)_{m\ge 1}\) 为
\[
(Tf)_m:=f_{m+1}.
\]
取任意多项式
\[
Q(\lambda)=\sum_{j=0}^{J}c_j\lambda^j
\]
满足零和条件
\[
Q(1)=\sum_{j=0}^{J}c_j=0.
\]
则对任意固定整数 \(R\ge 1\)，有
\[
Q(T)\log C_m
=
\frac{\sqrt5}{\varphi^2}
\sum_{r=1}^{R}
\frac{B_{2r}}{r(2r-1)}
\mathcal R\bigl(-(2r-1)\bigr)\,
Q\!\bigl(\varrho^{2r-1}\bigr)\,
\varrho^{(2r-1)m}
+
O\!\bigl(\varrho^{(2R+1)m}\bigr).
\]
因此，一切有限窗、零和的对数观测量，例如
\[
\log\frac{C_{m+1}}{C_m},
\qquad
\log C_{m+2}-2\log C_{m+1}+\log C_m,
\]
在渐近层面都完全由 \(\mathcal R\) 决定；其中规范常数 \(\log(2\pi)\) 严格消失。
\end{theorem}
\begin{proof}
由定理 \ref{thm:xi-time-part70e-renyi-curve-negative-odd-sampling}，
\[
\log C_m
=
\log(2\pi)
+
\frac{\sqrt5}{\varphi^2}
\sum_{r=1}^{R}
\frac{B_{2r}}{r(2r-1)}
\mathcal R\bigl(-(2r-1)\bigr)\,
\varrho^{(2r-1)m}
+
O\!\bigl(\varrho^{(2R+1)m}\bigr).
\]
对上式施加算子 \(Q(T)\)。首先
\[
Q(T)\log(2\pi)=Q(1)\log(2\pi)=0.
\]
其次，对任意 \(r\ge 1\) 都有
\[
T\bigl(\varrho^{(2r-1)m}\bigr)=\varrho^{2r-1}\varrho^{(2r-1)m},
\]
故
\[
Q(T)\bigl(\varrho^{(2r-1)m}\bigr)
=
Q\!\bigl(\varrho^{2r-1}\bigr)\varrho^{(2r-1)m}.
\]
逐项代回即可得到所示公式。证毕。
\end{proof}

\begin{theorem}[热力学 germ 的一维 Stokes 归一化刚性]\label{thm:xi-time-part70e-stokes-normalization-rigidity}
设 \((G_m)_{m\ge 1}\) 为一列正数，并假定存在常数 \(c\in\RR\)，使得对每个固定整数 \(R\ge 1\) 都有
\[
\log G_m
=
c+
\frac{\sqrt5}{\varphi^2}
\sum_{r=1}^{R}
\frac{B_{2r}}{r(2r-1)}
\mathcal R\bigl(-(2r-1)\bigr)\,
\varrho^{(2r-1)m}
+
O\!\bigl(\varrho^{(2R+1)m}\bigr).
\]
则
\[
\log G_m-\log C_m
=
c-\log(2\pi)+O\!\bigl(\varrho^{(2R+1)m}\bigr)
\qquad (m\to\infty),
\]
从而
\[
\frac{G_m}{C_m}\longrightarrow e^{\,c-\log(2\pi)}.
\]
亦即：一旦整函数 \(\mathcal R=\mathcal A\) 被固定，整个 fold-gauge germ 在全部归一化无关方向上都已刚性化；剩余自由度严格只保留一个乘法常数，而本文中的 \(2\pi\) 正是该唯一自由度的规范取值。
\end{theorem}
\begin{proof}
由定理 \ref{thm:xi-time-part70e-renyi-curve-negative-odd-sampling}，
\[
\log C_m
=
\log(2\pi)+
\frac{\sqrt5}{\varphi^2}
\sum_{r=1}^{R}
\frac{B_{2r}}{r(2r-1)}
\mathcal R\bigl(-(2r-1)\bigr)\,
\varrho^{(2r-1)m}
+
O\!\bigl(\varrho^{(2R+1)m}\bigr).
\]
将其与 \(\log G_m\) 的假设展开相减，级数主项逐项抵消，即得
\[
\log G_m-\log C_m
=
c-\log(2\pi)+O\!\bigl(\varrho^{(2R+1)m}\bigr).
\]
取 \(R=1\) 并令 \(m\to\infty\)，得到
\[
\log G_m-\log C_m\to c-\log(2\pi).
\]
指数化即得所示比值极限。证毕。
\end{proof}

\begin{conjecture}[连续射影谱的精确闭式与整振幅插值]\label{conj:xi-time-part70e-projective-pressure-entire-amplitude}
沿用附录 \ref{app:pom-projective-pressure-operator} 的射影转移算子 \(\mathcal L_t\)（\(t\ge 1\)）。则其主谱半径满足精确闭式
\[
\rho(\mathcal L_t)=2^t\varphi^{1-t}
\qquad (t\ge 1),
\]
等价地，连续归一化压力
\[
\log\rho(\mathcal L_t)-t\log 2=(1-t)\log\varphi
\]
在整个半轴上恰为一条仿射直线。进一步，存在一种自然归一化，使得主特征数据在 \(t\)-轴上的解析延拓恰给出定理 \ref{thm:xi-time-part70e-binfold-entire-amplitude-law} 的整振幅
\[
\mathcal A(t)=\frac{\varphi+\varphi^{-t}}{\sqrt5}.
\]
若该猜想成立，则定理 \ref{thm:xi-time-part70e-renyi-curve-negative-odd-sampling} 中的尾部采样公式便不再只是 bin-fold moment asymptotics 的结果，而可被提升为连续射影主谱在负奇整数点的 Bernoulli 采样律。
\end{conjecture}
\begin{remark}
该猜想与现有证据链完全同向。首先，定理 \ref{thm:pom-real-q-pressure-analytic-interpolation} 已把整数 moment-kernel 塔嵌入一条实解析压力曲线；定理 \ref{thm:pom-projective-operator-degenerates-to-moment-kernel} 与定理 \ref{thm:xi-projective-moment-tower-anchors-normalized-pressure} 则说明，整数节点上的主谱数据正是离散碰撞谱 \(r_q\)。其次，命题 \ref{prop:fold-bin-power-sum-asymptotic} 已经给出
\[
r_q=2^q\varphi^{1-q}
\qquad (q\in\ZZ_{\ge 2}),
\]
而定理 \ref{thm:xi-time-part70e-binfold-entire-amplitude-law} 与定理 \ref{thm:xi-time-part70e-renyi-curve-negative-odd-sampling} 又表明，同一条离散谱线在复参数方向上伴随一个刚性的整振幅 \(\mathcal A\)，并且该振幅的负奇采样恰重建 \(\log C_m\) 的 Bernoulli 奇层。因而，从离散 moment-kernel、Rényi 极限曲线与 fold-gauge 尾部采样三侧同时观察，全部可见证据都指向同一连续谱公式。
\end{remark}

\noindent
由此得到的图景是：本文框架中的 \(\pi\)-机制并不承载于压力函数的曲率，而承载于一条全轴仿射的热力学骨架之上的整函数振幅及其负奇采样。换言之，指数率本身已经被压平为单一直线；真正保留归一化信息的，是该直线上方的整振幅 \(\mathcal A=\mathcal R\)，以及由 Bernoulli 变换读出的唯一 Stokes 型常数 \(2\pi\)。于是，bin-fold 的复矩谱、Rényi 极限曲线、有限窗消模器与 fold-gauge 常数 germ 不再是并列现象，而被收束为同一条解析主线。

\endinput


\paragraph{零余类有限交截、nerve 的 flag 刚性、加权 clique--Euler 并集公式与共振整函数系数刚性}
在零点余类刚性、约数驱动的谱零点余类并结构、零点削减碰撞上界，以及共振梯整函数的 Laguerre--P\'olya 闭包、实零点结构与截断误差已经建立之后，还可继续沿同一条算术--解析主线抽取出下列七条结果。它们分别刻画零余类族任意有限交截的广义中国剩余判据、零点余类覆盖的 Helly--2 刚性与 nerve 的 flag 闭式、谱零点集合的精确加权 clique--Euler 公式、碰撞上界的约数兼容图闭式、共振整函数 \(\mathcal F_\varphi\) 的对数 Taylor 级数与倒数零矩、其偶系数的显式 Newton 递推与严格对数凹性，以及截断深度对相对误差预算的可审计控制。

\begin{theorem}[零余类族的任意有限交截由两两兼容性完全控制]\label{thm:xi-time-part71-zero-coset-finite-intersection-pairwise-compatibility}
\leanverified{paper\_xi\_time\_part71\_zero\_coset\_finite\_intersection\_pairwise\_compatibility}
设 \(M\) 为偶整数，且
\[
g_1,\dots,g_r\mid \frac{M}{2}.
\]
对每个 \(g\mid M/2\)，定义
\[
S_g(M):=
\left\{
k\in \ZZ/(M\ZZ):
\ k\equiv \frac{M}{2g}\pmod{M/g}
\right\}.
\]
则对任意有限族 \(\{g_1,\dots,g_r\}\)，都有
\[
\bigcap_{j=1}^{r}S_{g_j}(M)\neq\varnothing
\iff
2\gcd(g_i,g_j)\mid (g_j-g_i)
\qquad
(1\le i<j\le r).
\]
并且一旦交截非空，就有精确基数公式
\[
\left|
\bigcap_{j=1}^{r}S_{g_j}(M)
\right|
=
\gcd(g_1,\dots,g_r).
\]
\end{theorem}
\begin{proof}
由定理 \ref{thm:fold-zero-coset-rigidity}，
\[
S_{g_j}(M)
=
\left\{
k\in \ZZ/(M\ZZ):
\ k\equiv \frac{M}{2g_j}\pmod{M/g_j}
\right\}
\qquad (1\le j\le r).
\]
因此
\[
\bigcap_{j=1}^{r}S_{g_j}(M)\neq\varnothing
\]
当且仅当同余组
\[
k\equiv \frac{M}{2g_j}\pmod{M/g_j}
\qquad (j=1,\dots,r)
\]
可解。广义中国剩余定理断言：这一线性同余组可解，当且仅当任意两式相容。另一方面，引理 \ref{lem:fold-zero-coset-intersection} 已给出两式相容的充要条件恰为
\[
2\gcd(g_i,g_j)\mid (g_j-g_i).
\]
于是得到第一条等价关系。

现设该系统可解。由于解集模数为
\[
\operatorname{lcm}\!\left(\frac{M}{g_1},\dots,\frac{M}{g_r}\right)
=
\frac{M}{\gcd(g_1,\dots,g_r)},
\]
故在模 \(M\) 的意义下，解恰形成一个单一同余类的提升，其基数为
\[
\frac{M}{M/\gcd(g_1,\dots,g_r)}
=
\gcd(g_1,\dots,g_r).
\]
证毕。
\end{proof}

\begin{corollary}[零点余类覆盖的 Helly--2 刚性与 nerve 的 flag 闭式]\label{cor:xi-time-part71-zero-coset-nerve-flag}
沿用定理 \ref{thm:xi-time-part71-zero-coset-finite-intersection-pairwise-compatibility} 的记号。任取有限集合
\[
\mathcal G\subseteq
\left\{
g:\ g\mid \frac{M}{2}
\right\},
\]
并考虑覆盖
\[
\mathcal U_{\mathcal G}:=\{S_g(M):\ g\in\mathcal G\}.
\]
定义其 nerve 为
\[
\mathcal N(\mathcal U_{\mathcal G})
:=
\left\{
\sigma\subseteq \mathcal G:\ \bigcap_{g\in \sigma}S_g(M)\neq\varnothing
\right\}.
\]
再在顶点集 \(\mathcal G\) 上定义算术兼容图 \(\Gamma_{\mathcal G}\)：对不同的 \(g,h\in\mathcal G\)，规定
\[
g\sim h
\iff
2\gcd(g,h)\mid (h-g).
\]
则 \(\mathcal N(\mathcal U_{\mathcal G})\) 恰为 \(\Gamma_{\mathcal G}\) 的 clique 复形。特别地，\(\mathcal N(\mathcal U_{\mathcal G})\) 是一个 flag 复形：任意顶点集若两两相邻，则自动张成一个单纯形。
\leanverified{paper\_xi\_time\_part71\_zero\_coset\_nerve\_flag}
\end{corollary}
\begin{proof}
由定理 \ref{thm:xi-time-part71-zero-coset-finite-intersection-pairwise-compatibility}，对任意有限子集
\[
\sigma\subseteq \mathcal G
\]
都有
\[
\bigcap_{g\in \sigma}S_g(M)\neq\varnothing
\iff
2\gcd(g,h)\mid (h-g)
\qquad
(\forall\, g\neq h\in \sigma).
\]
右侧正表明 \(\sigma\) 在图 \(\Gamma_{\mathcal G}\) 中形成一个 clique。故
\[
\sigma\in \mathcal N(\mathcal U_{\mathcal G})
\iff
\sigma\in \mathrm{Cliq}(\Gamma_{\mathcal G}),
\]
即 nerve 恰为该兼容图的 clique 复形。于是它自动是 flag 复形。证毕。
\end{proof}

\endinput


\begin{theorem}[谱零点集合的精确加权 clique--Euler 公式]\label{thm:xi-time-part71-zero-spectrum-weighted-clique-euler}
\leanverified{paper\_xi\_time\_part71\_zero\_spectrum\_weighted\_clique\_euler}
设
\[
M:=F_{m+2},
\qquad
\mathcal Z_m:=\{k\in \ZZ/(M\ZZ):\ \widehat d_m(k)=0\},
\]
并假设 \(M\) 为偶数。定义
\[
\mathcal D_m:=
\left\{
d:\ d\mid (m+2),\ d<m+2,\ F_d\mid \frac{M}{2}
\right\}.
\]
在顶点集 \(\mathcal D_m\) 上定义兼容图 \(\Gamma_m\)：对不同顶点 \(d,e\in\mathcal D_m\)，规定
\[
d\sim e
\iff
2F_{\gcd(d,e)}\mid (F_e-F_d).
\]
若记 \(\mathrm{Cliq}(\Gamma_m)\) 为 \(\Gamma_m\) 的全部非空 clique，且对每个 clique \(C\) 记
\[
\gcd(C):=\gcd\{d:\ d\in C\},
\]
则谱零点集合满足严格公式
\[
|\mathcal Z_m|
=
\sum_{\varnothing\neq C\in \mathrm{Cliq}(\Gamma_m)}
(-1)^{|C|+1}F_{\gcd(C)}.
\]
\end{theorem}
\begin{proof}
由推论 \ref{cor:fold-zero-coset-union}，当 \(M\) 为偶数时，
\[
\mathcal Z_m
=
\bigcup_{\substack{1\le j\le m\\ g_j\mid M/2}}S_{g_j}(M),
\qquad
g_j=\gcd(F_{j+1},M)=F_{\gcd(j+1,m+2)}.
\]
去重之后，恰可写成
\[
\mathcal Z_m=\bigcup_{d\in\mathcal D_m}S_{F_d}(M).
\]

现取任意非空子集 \(I=\{d_1,\dots,d_r\}\subseteq \mathcal D_m\)。由定理 \ref{thm:xi-time-part71-zero-coset-finite-intersection-pairwise-compatibility}，
\[
\bigcap_{d\in I}S_{F_d}(M)\neq\varnothing
\]
当且仅当对所有 \(1\le i<j\le r\) 都有
\[
2\gcd(F_{d_i},F_{d_j})\mid (F_{d_j}-F_{d_i}).
\]
再由 Fibonacci 最大公因子恒等式
\[
\gcd(F_a,F_b)=F_{\gcd(a,b)},
\]
上式等价于
\[
2F_{\gcd(d_i,d_j)}\mid (F_{d_j}-F_{d_i}),
\]
也即 \(I\) 在 \(\Gamma_m\) 中形成一个 clique。

若交截非空，则再由定理 \ref{thm:xi-time-part71-zero-coset-finite-intersection-pairwise-compatibility}，
\[
\left|
\bigcap_{d\in I}S_{F_d}(M)
\right|
=
\gcd(F_{d_1},\dots,F_{d_r})
=
F_{\gcd(d_1,\dots,d_r)}
=
F_{\gcd(I)}.
\]
因此对并集
\[
\mathcal Z_m=\bigcup_{d\in\mathcal D_m}S_{F_d}(M)
\]
施行包含--排斥，便得到
\[
|\mathcal Z_m|
=
\sum_{\varnothing\neq C\in \mathrm{Cliq}(\Gamma_m)}
(-1)^{|C|+1}F_{\gcd(C)}.
\]
证毕。
\end{proof}

\begin{corollary}[碰撞上界可压缩为约数兼容图的加权 Euler 特征]\label{cor:xi-time-part71-collision-upper-bound-weighted-clique-euler}
沿用定理 \ref{thm:xi-time-part71-zero-spectrum-weighted-clique-euler} 的记号。则碰撞概率满足严格上界
\[
\mathsf{Col}_m
\le
1-\frac{1}{F_{m+2}}
\sum_{\varnothing\neq C\in \mathrm{Cliq}(\Gamma_m)}
(-1)^{|C|+1}F_{\gcd(C)}.
\]
特别地，若 \(\Gamma_m\) 不含三角形，亦即其非空 clique 仅有单点与边，则上界简化为
\[
\mathsf{Col}_m
\le
1-\frac{1}{F_{m+2}}
\left(
\sum_{d\in\mathcal D_m}F_d
-
\sum_{\{d,e\}\in E(\Gamma_m)}F_{\gcd(d,e)}
\right).
\]
\end{corollary}
\begin{proof}
由定理 \ref{thm:fold-collision-zero-reduction}，
\[
\mathsf{Col}_m\le 1-\frac{|\mathcal Z_m|}{F_{m+2}}.
\]
将定理 \ref{thm:xi-time-part71-zero-spectrum-weighted-clique-euler} 的精确公式代入，即得第一式。

若 \(\Gamma_m\) 不含三角形，则全部非空 clique 仅由单点与边组成，包含--排斥求和自动截断为两层：
\[
|\mathcal Z_m|
=
\sum_{d\in\mathcal D_m}F_d
-
\sum_{\{d,e\}\in E(\Gamma_m)}F_{\gcd(d,e)}.
\]
再代回碰撞上界即可。证毕。
\end{proof}

\begin{theorem}[共振整函数的对数 Taylor 级数与倒数零矩闭式]\label{thm:xi-time-part71-resonance-entire-log-taylor-zero-moments}
沿用定理 \ref{thm:fold-resonance-entire-lp} 的记号，记
\[
\mathcal F_\varphi(z):=\prod_{n=2}^{\infty}\cos(\pi z\varphi^{-n}).
\]
则在最近正零点半径
\[
|z|<\frac{\varphi^2}{2}
\]
内，有严格对数展开
\[
\log \mathcal F_\varphi(z)
=
-\sum_{r=1}^{\infty}\frac{\sigma_r}{r}\,z^{2r},
\qquad
\sigma_r:=
(2^{2r}-1)\zeta(2r)\frac{\varphi^{-2r}}{\varphi^{2r}-1}.
\]
等价地，正零点倒数偶矩满足
\[
\sum_{\substack{\rho\in\mathcal Z(\mathcal F_\varphi)\\ \rho>0}}\rho^{-2r}
=
(2^{2r}-1)\zeta(2r)\frac{\varphi^{-2r}}{\varphi^{2r}-1},
\]
而全零点求和式为
\[
\sum_{\rho\in\mathcal Z(\mathcal F_\varphi)}\rho^{-2r}
=
2(2^{2r}-1)\zeta(2r)\frac{\varphi^{-2r}}{\varphi^{2r}-1}.
\]
\end{theorem}
\begin{proof}
由定理 \ref{thm:fold-resonance-entire-lp}，\(\mathcal F_\varphi\) 为整函数，且
\[
\mathcal Z(\mathcal F_\varphi)
=
\left\{
\left(k+\tfrac12\right)\varphi^n:\ k\in\ZZ,\ n\ge 2
\right\}.
\]
当 \(|z|<\varphi^2/2\) 时，对所有 \(n\ge 2\) 都有
\[
|\pi z\varphi^{-n}|<\frac{\pi}{2},
\]
故可使用经典展开
\[
\log\cos x
=
-\sum_{r=1}^{\infty}\frac{(2^{2r}-1)\zeta(2r)}{r\pi^{2r}}\,x^{2r}
\qquad
(|x|<\pi/2).
\]
逐项代入 \(x=\pi z\varphi^{-n}\) 并交换求和，得到
\[
\log \mathcal F_\varphi(z)
=
\sum_{n=2}^{\infty}\log\cos(\pi z\varphi^{-n})
=
-\sum_{r=1}^{\infty}\frac{(2^{2r}-1)\zeta(2r)}{r}
\left(
\sum_{n=2}^{\infty}\varphi^{-2rn}
\right)z^{2r}.
\]
而
\[
\sum_{n=2}^{\infty}\varphi^{-2rn}
=
\frac{\varphi^{-4r}}{1-\varphi^{-2r}}
=
\frac{\varphi^{-2r}}{\varphi^{2r}-1},
\]
遂得
\[
\log \mathcal F_\varphi(z)
=
-\sum_{r=1}^{\infty}\frac{\sigma_r}{r}\,z^{2r}.
\]

另一方面，由定理 \ref{thm:fold-resonance-entire-lp} 的零点描述，
\[
\mathcal Z(\mathcal F_\varphi)\cap (0,\infty)
=
\left\{
\left(k+\tfrac12\right)\varphi^n:\ k\ge 0,\ n\ge 2
\right\}.
\]
故
\[
\sum_{\substack{\rho\in\mathcal Z(\mathcal F_\varphi)\\ \rho>0}}\rho^{-2r}
=
\sum_{n=2}^{\infty}\varphi^{-2rn}\sum_{k\ge 0}\left(k+\tfrac12\right)^{-2r}.
\]
又有
\[
\sum_{k\ge 0}\left(k+\tfrac12\right)^{-2r}
=
2^{2r}\sum_{k\ge 0}(2k+1)^{-2r}
=
(2^{2r}-1)\zeta(2r),
\]
故正零点公式成立。全零点关于原点成对对称，因此再乘以 \(2\) 即得最后一式。证毕。
\end{proof}

\begin{theorem}[共振整函数偶系数的显式 Newton 递推与严格对数凹性]\label{thm:xi-time-part71-resonance-entire-even-coefficient-newton-logconcavity}
\leanverified{paper\_xi\_time\_part71\_resonance\_entire\_even\_coefficient\_newton\_logconcavity}
将 \(\mathcal F_\varphi\) 写成偶幂级数
\[
\mathcal F_\varphi(z)=\sum_{n\ge 0}a_{2n}z^{2n},
\qquad
a_{2n+1}=0,
\]
并定义
\[
b_n:=(-1)^n a_{2n}
\qquad (n\ge 0).
\]
则对全部 \(n\ge 0\) 都有 \(b_n>0\)，并满足精确 Newton 递推
\[
n\,b_n
=
\sum_{r=1}^{n}(-1)^{r-1}\sigma_r\,b_{n-r}
\qquad (n\ge 1),
\]
其中 \(\sigma_r\) 如定理 \ref{thm:xi-time-part71-resonance-entire-log-taylor-zero-moments} 所示。进一步，序列 \((b_n)_{n\ge 0}\) 严格对数凹：
\[
b_n^2>b_{n-1}b_{n+1}
\qquad (n\ge 1).
\]
因此偶系数满足严格交替符号律
\[
a_0>0,\qquad a_2<0,\qquad a_4>0,\qquad \dots,
\]
且相邻比值 \((b_{n+1}/b_n)_{n\ge 0}\) 严格递减。
\end{theorem}
\begin{proof}
先由定理 \ref{thm:xi-time-part71-resonance-entire-log-taylor-zero-moments} 定义形式幂级数
\[
G(t):=\mathcal F_\varphi(\sqrt t)
=
\sum_{n\ge 0}a_{2n}t^n
=
\sum_{n\ge 0}(-1)^n b_n t^n.
\]
则
\[
\log G(t)
=
-\sum_{r=1}^{\infty}\frac{\sigma_r}{r}t^r.
\]
对 \(t\) 求导得
\[
\frac{G'(t)}{G(t)}
=
-\sum_{r=1}^{\infty}\sigma_r t^{r-1}.
\]
把
\[
G(t)=\sum_{n\ge 0}(-1)^n b_n t^n
\]
代入，并比较 \(t^{n-1}\) 的系数，即得
\[
n\,(-1)^n b_n
=
-\sum_{r=1}^{n}\sigma_r(-1)^{n-r}b_{n-r},
\]
等价于
\[
n\,b_n
=
\sum_{r=1}^{n}(-1)^{r-1}\sigma_r\,b_{n-r}.
\]

接着证明 \(b_n>0\) 与严格对数凹性。由定理 \ref{thm:fold-resonance-entire-lp}，
\(\mathcal F_\varphi\) 的正零点可记为
\[
0<\rho_1<\rho_2<\cdots,
\]
且由定理 \ref{thm:xi-time-part71-resonance-entire-log-taylor-zero-moments} 在 \(r=1\) 时的公式，
\[
\sum_{j\ge 1}\rho_j^{-2}<\infty.
\]
因此
\[
G(t)=\prod_{j=1}^{\infty}\left(1-\frac{t}{\rho_j^2}\right)
\]
在 \(t=0\) 邻域内成立，并且
\[
b_n=e_n(\rho_1^{-2},\rho_2^{-2},\dots)>0
\]
是正数列 \((\rho_j^{-2})_{j\ge 1}\) 的第 \(n\) 个初等对称函数。

对每个 \(N>n\)，定义截断多项式
\[
G_N(t):=\prod_{j=1}^{N}\left(1-\frac{t}{\rho_j^2}\right)
=
\sum_{s=0}^{N}(-1)^s b_s^{(N)}t^s.
\]
则 \(b_s^{(N)}\to b_s\)（\(s\) 固定）且 \(b_s^{(N)}>0\)。Newton 不等式给出
\[
\bigl(b_n^{(N)}\bigr)^2
\ge
\frac{n+1}{n}\frac{N-n+1}{N-n}\,
b_{n-1}^{(N)}b_{n+1}^{(N)}.
\]
令 \(N\to\infty\)，得到
\[
b_n^2
\ge
\frac{n+1}{n}\,b_{n-1}b_{n+1}
>
b_{n-1}b_{n+1}.
\]
这就证明了严格对数凹性。由于 \(b_n>0\)，于是
\[
a_{2n}=(-1)^n b_n
\]
严格交替。再将
\[
b_n^2>b_{n-1}b_{n+1}
\]
两边除以 \(b_{n-1}b_n>0\)，即可得到
\[
\frac{b_{n+1}}{b_n}<\frac{b_n}{b_{n-1}},
\]
故相邻比值严格递减。证毕。
\end{proof}

\begin{corollary}[共振梯截断深度的相对误差预算]\label{cor:xi-time-part71-resonance-truncation-relative-error-budget}
\leanverified{paper\_xi\_time\_part71\_resonance\_truncation\_relative\_error\_budget}
沿用命题 \ref{prop:fold-resonance-truncation-error} 的记号。对任意 \(\varepsilon>0\)，当 \(u=0\) 时有
\[
\frac{C_{\varphi,N}(0)}{C_\varphi(0)}=1
\qquad
(\forall N\ge 2).
\]
若 \(u\neq 0\)、\(N\ge 2\) 且
\[
N\ \ge\
-1+\max\!\left\{
\log_\varphi(2\pi|u|),\
\frac12\log_\varphi\!\left(
\frac{\pi^2u^2}{(1-\varphi^{-2})\log(1+\varepsilon)}
\right)
\right\},
\]
则有相对误差界
\[
1\le \frac{C_{\varphi,N}(u)}{C_\varphi(u)}\le 1+\varepsilon.
\]
\end{corollary}
\begin{proof}
当 \(u=0\) 时，
\[
C_{\varphi,N}(0)=C_\varphi(0)=1,
\]
结论平凡。以下设 \(u\neq 0\)。

由命题 \ref{prop:fold-resonance-truncation-error}，只要满足
\[
\pi|u|\varphi^{-(N+1)}\le \frac12,
\]
就有
\[
C_{\varphi,N}(u)\exp\!\left(
-\frac{\pi^2u^2\varphi^{-2(N+1)}}{1-\varphi^{-2}}
\right)
\le
C_\varphi(u)
\le
C_{\varphi,N}(u).
\]
因此
\[
1
\le
\frac{C_{\varphi,N}(u)}{C_\varphi(u)}
\le
\exp\!\left(
\frac{\pi^2u^2\varphi^{-2(N+1)}}{1-\varphi^{-2}}
\right).
\]
若进一步要求
\[
\frac{\pi^2u^2\varphi^{-2(N+1)}}{1-\varphi^{-2}}
\le
\log(1+\varepsilon),
\]
则立得
\[
\frac{C_{\varphi,N}(u)}{C_\varphi(u)}\le 1+\varepsilon.
\]
最后，把上述两条约束分别对 \(N\) 求解，便得到
\[
N\ \ge\
-1+\max\!\left\{
\log_\varphi(2\pi|u|),\
\frac12\log_\varphi\!\left(
\frac{\pi^2u^2}{(1-\varphi^{-2})\log(1+\varepsilon)}
\right)
\right\}.
\]
证毕。
\end{proof}

\noindent
上述七条结果表明，Fibonacci 模零谱与共振整函数 \(\mathcal F_\varphi\) 已经形成一条严格可逆的算术--解析双轨：一方面，谱零点集合 \(\mathcal Z_m\) 不再只是若干候选余类的松散并，而可被零余类覆盖的 flag nerve、约数兼容图的 clique 结构及其加权 Euler 特征精确决定，并进一步把碰撞上界压缩为有限图论计算；另一方面，\(\mathcal F_\varphi\) 的局部解析则不再停留于 Laguerre--P\'olya 性与 Jensen 多项式全实根性这样的抽象实根性，而会在对数 Taylor 系数、倒数零矩、偶系数 Newton 递推与截断深度预算中外显为完全可审计的 \(\zeta(2r)\)--\(\varphi\) 谱矩闭式。于是，零余类几何、有限交 nerve、碰撞上界与共振常数计算深度便被统一收束到同一条由 Fibonacci 约数格与黄金比例谱矩共同支配的闭合链条之中。

\paragraph{奇倍嵌套链、约数格指标容斥、线性零点增量与六倍窗密度指数率}
在零余类族的 \(\gcd\)-刚性交结构与谱零点集合的加权 clique--Euler 公式已经建立之后，还可继续沿约数格方向抽取出下列四条直接后果。它们分别刻画整除链上的奇倍嵌套--偶倍排斥二分律、以相容指标书写的约数格精确容斥公式、pairwise-incompatible 约数族所导致的线性零点增量下界，以及六倍窗上谱零点相对密度的精确指数率。

\begin{theorem}[零余类沿整除链的奇倍嵌套与偶倍排斥]\label{thm:xi-time-part72-zero-coset-divisibility-nested-disjoint-dichotomy}
\leanverified{paper\_xi\_time\_part72\_zero\_coset\_divisibility\_nested\_disjoint\_dichotomy}
设 \(M\) 为偶整数，并令
\[
g\mid h\mid \frac{M}{2}.
\]
则有完全二分律
\[
S_g(M)\subseteq S_h(M)
\iff
\frac{h}{g}\ \text{为奇数},
\]
\[
S_g(M)\cap S_h(M)=\varnothing
\iff
\frac{h}{g}\ \text{为偶数}.
\]
特别地，对任意整除链
\[
g_1\mid g_2\mid \cdots \mid g_r\mid \frac{M}{2},
\]
若所有相邻商
\[
\frac{g_{i+1}}{g_i}\qquad (1\le i<r)
\]
都为奇数，则
\[
S_{g_1}(M)\subseteq S_{g_2}(M)\subseteq \cdots \subseteq S_{g_r}(M),
\qquad
\left|\bigcup_{i=1}^{r}S_{g_i}(M)\right|=g_r.
\]
反之，只要存在某个 \(i\) 使 \(g_{i+1}/g_i\) 为偶数，便有
\[
S_{g_i}(M)\cap S_{g_{i+1}}(M)=\varnothing.
\]
\end{theorem}
\begin{proof}
由定理 \ref{thm:xi-fold-zero-coset-v2-intersection-rigidity}，
\[
S_g(M)\cap S_h(M)=
\begin{cases}
S_{\gcd(g,h)}(M),& v_2(g)=v_2(h),\\[2mm]
\varnothing,& v_2(g)\neq v_2(h).
\end{cases}
\]
在 \(g\mid h\) 的情形，
\[
\gcd(g,h)=g.
\]
又
\[
v_2(g)=v_2(h)
\iff
v_2(h/g)=0
\iff
\frac{h}{g}\ \text{为奇数}.
\]
因此当 \(h/g\) 为奇数时，
\[
S_g(M)\cap S_h(M)=S_g(M),
\]
从而 \(S_g(M)\subseteq S_h(M)\)；当 \(h/g\) 为偶数时，
\[
v_2(g)\neq v_2(h),
\]
故交集为空。

对整除链逐步应用上述二分律，便得到链式包含关系与相邻偶商导致的互斥结论。若全部相邻商皆为奇数，则整条链单调嵌套，故其并集等于最大项 \(S_{g_r}(M)\)，于是
\[
\left|\bigcup_{i=1}^{r}S_{g_i}(M)\right|=|S_{g_r}(M)|=g_r.
\]
证毕。
\end{proof}

\begin{theorem}[谱零点集合的约数格指标容斥公式]\label{thm:xi-time-part72-zero-spectrum-divisor-lattice-indicator-inclusion-exclusion}
\leanverified{paper\_xi\_time\_part72\_zero\_spectrum\_divisor\_lattice\_indicator\_inclusion\_exclusion}
沿用定理 \ref{thm:xi-time-part71-zero-spectrum-weighted-clique-euler} 的记号，令
\[
\mathcal D_m:=
\left\{
d:\ d\mid (m+2),\ d<m+2,\ F_d\mid \frac{F_{m+2}}{2}
\right\}.
\]
对每个非空子集 \(I\subseteq \mathcal D_m\)，定义
\[
\gcd(I):=\gcd\{d:\ d\in I\},
\]
并置相容指标
\[
\mathcal C(I)
\iff
2F_{\gcd(d,e)}\mid (F_e-F_d)
\qquad
(\forall\, d,e\in I,\ d\neq e).
\]
则谱零点集合满足完全显式的约数格容斥展开
\[
|\mathcal Z_m|
=
\sum_{\varnothing\neq I\subseteq \mathcal D_m}
(-1)^{|I|+1}\mathbf 1_{\mathcal C(I)}\,F_{\gcd(I)}.
\]
换言之，\(\mathcal Z_m\) 的基数可由 \((m+2)\) 的约数格与 Fibonacci 差值的整除判别完全决定，而无需再对频率 \(k\) 作逐点扫描。
\end{theorem}
\begin{proof}
由定理 \ref{thm:xi-time-part71-zero-spectrum-weighted-clique-euler}，
\[
|\mathcal Z_m|
=
\sum_{\varnothing\neq C\in \mathrm{Cliq}(\Gamma_m)}
(-1)^{|C|+1}F_{\gcd(C)},
\]
其中 \(\Gamma_m\) 的边关系定义为
\[
d\sim e
\iff
2F_{\gcd(d,e)}\mid (F_e-F_d).
\]
因此，对任意非空子集 \(I\subseteq \mathcal D_m\)，
\[
I\in \mathrm{Cliq}(\Gamma_m)
\iff
\mathcal C(I).
\]
于是把 clique 求和改写为对全部非空子集的求和，并以指标函数 \(\mathbf 1_{\mathcal C(I)}\) 过滤不相容子集，便得到
\[
|\mathcal Z_m|
=
\sum_{\varnothing\neq I\subseteq \mathcal D_m}
(-1)^{|I|+1}\mathbf 1_{\mathcal C(I)}\,F_{\gcd(I)}.
\]
证毕。
\end{proof}

\begin{corollary}[pairwise-incompatible 约数族给出线性零点增量下界]\label{cor:xi-time-part72-zero-spectrum-pairwise-incompatible-linear-lower-bound}
沿用定理 \ref{thm:xi-time-part72-zero-spectrum-divisor-lattice-indicator-inclusion-exclusion} 的记号。若子集
\[
\mathcal A\subseteq \mathcal D_m
\]
满足对任意不同的 \(d,e\in\mathcal A\) 都有
\[
2F_{\gcd(d,e)}\nmid (F_e-F_d),
\]
则相应余类族两两不交，从而
\[
|\mathcal Z_m|
\ge
\sum_{d\in\mathcal A}F_d.
\]
特别地，当 \(m+2\equiv 0\pmod 6\) 且
\[
d:=\frac{m+2}{2},
\]
由推论 \ref{cor:fold-zero-half-index-multiple6} 已知
\[
F_d\mid \frac{F_{m+2}}{2},
\qquad
|\mathcal Z_m|\ge F_d.
\]
若再存在 \(e\in\mathcal D_m\) 满足
\[
2F_{\gcd(d,e)}\nmid (F_e-F_d),
\]
则立刻提升为
\[
|\mathcal Z_m|\ge F_d+F_e.
\]
\end{corollary}
\begin{proof}
由定理 \ref{thm:xi-time-part71-zero-coset-finite-intersection-pairwise-compatibility}，对不同的 \(d,e\in \mathcal A\)，条件
\[
2F_{\gcd(d,e)}\nmid (F_e-F_d)
\]
等价于
\[
S_{F_d}(F_{m+2})\cap S_{F_e}(F_{m+2})=\varnothing.
\]
故族
\[
\{S_{F_d}(F_{m+2}):\ d\in\mathcal A\}
\]
两两不交。于是
\[
\left|
\bigcup_{d\in\mathcal A}S_{F_d}(F_{m+2})
\right|
=
\sum_{d\in\mathcal A}|S_{F_d}(F_{m+2})|
=
\sum_{d\in\mathcal A}F_d,
\]
其中最后一步使用定理 \ref{thm:fold-zero-coset-rigidity} 中
\[
|S_g(M)|=g.
\]
再由推论 \ref{cor:fold-zero-coset-union}，
\[
\bigcup_{d\in\mathcal A}S_{F_d}(F_{m+2})
\subseteq
\mathcal Z_m,
\]
遂得所示下界。

最后一段只需把 \(d=(m+2)/2\) 的六倍窗见证余类与一个与之 pairwise-incompatible 的 \(e\) 一并代入上述结论即可。证毕。
\end{proof}

\begin{theorem}[六倍窗上谱零点相对密度的精确指数率]\label{thm:xi-time-part72-zero-density-exact-exponential-rate}
沿用定理 \ref{thm:xi-time-part70b-dark-band-half-order-scale} 的记号。若
\[
m+2\equiv 0\pmod 6,
\qquad
M:=F_{m+2},
\]
则有极限公式
\[
\lim_{\substack{m\to\infty\\ m+2\equiv 0\!\!\!\pmod 6}}
\frac{1}{m}\log\frac{|\mathcal Z_m|}{F_{m+2}}
=
-\frac12\log\varphi.
\]
等价地，
\[
\frac{|\mathcal Z_m|}{F_{m+2}}
=
\exp\!\left(
-\frac{m}{2}\log\varphi+o(m)
\right)
\]
沿六倍窗子序列成立。
\end{theorem}
\begin{proof}
由定理 \ref{thm:xi-time-part70b-dark-band-half-order-scale}，
\[
\frac{1}{m}\log |\mathcal Z_m|
\longrightarrow
\frac12\log\varphi
\]
沿子序列 \(m+2\equiv 0\pmod 6\) 成立。另一方面，Binet 公式给出
\[
\log F_{m+2}=(m+2)\log\varphi+O(1),
\]
故
\[
\frac{1}{m}\log F_{m+2}\longrightarrow \log\varphi.
\]
两式相减即可得到
\[
\frac{1}{m}\log\frac{|\mathcal Z_m|}{F_{m+2}}
\longrightarrow
\frac12\log\varphi-\log\varphi
=
-\frac12\log\varphi.
\]
等价表述只是把极限式改写为指数渐近。证毕。
\end{proof}

\noindent
上述四条结果表明，零余类几何在有限交截判据与加权 clique--Euler 公式之外，还具有一条更细的约数格层级结构：沿整除链，余类并不会出现任意复杂的部分重叠，而只能在“奇商嵌套”与“偶商互斥”两种情形之间二分；一旦把这一二分律与相容指标容斥结合，任何 pairwise-incompatible 约数簇都会立刻转化为 \(\mathcal Z_m\) 的线性可加增量；而在六倍窗上，这些局部增量又与 half-order dark band 的双边界闭合到同一个精确指数率 \(\varphi^{-m/2}\)。于是，谱零点集合的算术几何便不仅可被精确计数，而且可在约数格上作系统的局部增量构造，并在特定同余窗上获得完全锁定的密度衰减主尺度。

\paragraph{奇商偏序下的零余类压缩生成、唯一极小覆盖与精确计数}
在定理 \ref{thm:xi-time-part71-zero-coset-finite-intersection-pairwise-compatibility} 已给出零余类族任意有限交截的完全分类、定理 \ref{thm:xi-time-part72-zero-coset-divisibility-nested-disjoint-dichotomy} 已把整除链上的交叠压缩为奇倍嵌套与偶倍排斥二分律之后，这条主线还可继续向“表示压缩”推进一步：零点集合 \(\mathcal Z_m\) 并不需要由全部候选余类来生成，而可唯一压缩到奇商偏序下的极大元族；相应的精确计数也可从全约数格容斥进一步压缩为该极大生成族上的包含--排斥闭式。

\begin{theorem}[零余类之间的全局包含关系由奇商整除完全刻画]\label{thm:xi-time-part72a-zero-coset-global-inclusion-odd-divisibility}
\leanpartial{paper\_xi\_time\_part72a\_zero\_coset\_global\_inclusion\_odd\_divisibility}{the requested Lean signature omits the paper's hypothesis that \(M\) is even; at \(M=0\), \(g=1\), and \(h=2\), both zero-cosets equal \(\{0\}\) but \(h/g=2\) is not odd}
设 \(M\) 为偶整数，且 \(g,h\mid M/2\)。则
\[
S_g(M)\subseteq S_h(M)
\iff
g\mid h\ \text{且}\ \frac{h}{g}\ \text{为奇数}.
\]
\end{theorem}
\begin{proof}
若 \(g\mid h\) 且 \(h/g\) 为奇数，则结论已由定理 \ref{thm:xi-time-part72-zero-coset-divisibility-nested-disjoint-dichotomy} 的第一条给出。

反之，设 \(S_g(M)\subseteq S_h(M)\)。则
\[
S_g(M)\cap S_h(M)=S_g(M),
\]
因而交集非空。由定理 \ref{thm:xi-time-part71-zero-coset-finite-intersection-pairwise-compatibility} 在 \(r=2\) 的情形，
\[
\left|S_g(M)\cap S_h(M)\right|=\gcd(g,h).
\]
又由定理 \ref{thm:fold-zero-coset-rigidity}，
\[
|S_g(M)|=g.
\]
于是
\[
\gcd(g,h)=|S_g(M)\cap S_h(M)|=|S_g(M)|=g,
\]
从而 \(g\mid h\)。再由交集非空判据，
\[
2\gcd(g,h)\mid (h-g),
\]
即
\[
2g\mid g\left(\frac{h}{g}-1\right).
\]
故 \(2\mid (h/g-1)\)，也就是 \(h/g\) 为奇数。证毕。
\end{proof}

\begin{theorem}[谱零点集合在奇商偏序下的唯一极小生成族]\label{thm:xi-time-part72a-zero-spectrum-maximal-odd-divisibility-compression}
\leanverified{paper\_xi\_time\_part72a\_zero\_spectrum\_maximal\_odd\_divisibility\_compression}
沿用推论 \ref{cor:fold-zero-coset-union} 与推论 \ref{cor:xi-fold-zero-set-v2-semilattice-decomposition} 的记号。记
\[
M:=F_{m+2},
\qquad
\mathcal D_m:=
\left\{
d:\ d\mid (m+2),\ d<m+2,\ F_d\mid \frac{M}{2}
\right\},
\]
\[
\mathcal G_m:=\{F_d:\ d\in \mathcal D_m\}.
\]
则
\[
\mathcal Z_m=\bigcup_{g\in \mathcal G_m}S_g(M).
\]
在 \(\mathcal G_m\) 上定义偏序
\[
g\preceq h
\iff
S_g(M)\subseteq S_h(M),
\]
等价地，由定理 \ref{thm:xi-time-part72a-zero-coset-global-inclusion-odd-divisibility}，
\[
g\preceq h
\iff
g\mid h\ \text{且}\ \frac{h}{g}\ \text{为奇数}.
\]
记 \(\mathcal M_m\subseteq \mathcal G_m\) 为全部 \(\preceq\)-极大元组成的集合。则成立：
\begin{enumerate}
  \item 谱零点集合具有压缩表示
  \[
  \mathcal Z_m=\bigcup_{g\in \mathcal M_m}S_g(M).
  \]
  \item \(\mathcal M_m\) 是满足上式的唯一极小生成族：若 \(\mathcal A\subseteq \mathcal G_m\) 也满足
  \[
  \mathcal Z_m=\bigcup_{g\in \mathcal A}S_g(M),
  \]
  则必有
  \[
  \mathcal M_m\subseteq \mathcal A.
  \]
\end{enumerate}
\end{theorem}
\begin{proof}
第一条只需删除全部非极大元。事实上，若 \(g\in \mathcal G_m\setminus \mathcal M_m\)，则存在 \(h\in \mathcal G_m\) 使 \(g\prec h\)。由定理 \ref{thm:xi-time-part72a-zero-coset-global-inclusion-odd-divisibility}，
\[
S_g(M)\subseteq S_h(M),
\]
故从并集中去掉 \(S_g(M)\) 不改变并集。对所有非极大元依次执行此操作，即得
\[
\mathcal Z_m=\bigcup_{g\in \mathcal M_m}S_g(M).
\]

现证唯一极小性。任取 \(g\in \mathcal M_m\)。若 \(h\in \mathcal G_m\setminus\{g\}\) 且
\[
S_g(M)\cap S_h(M)\neq \varnothing,
\]
则由定理 \ref{thm:xi-time-part71-zero-coset-finite-intersection-pairwise-compatibility}，
\[
S_g(M)\cap S_h(M)=S_{\gcd(g,h)}(M)
\]
且
\[
\gcd(g,h)<g,
\]
否则若 \(\gcd(g,h)=g\)，则 \(g\mid h\)，再由交非空判据可得 \(h/g\) 为奇数，从而 \(g\prec h\)，与 \(g\) 的极大性矛盾。

取 \(S_g(M)\) 的标准代表点
\[
k_g:=\frac{M}{2g}\in S_g(M).
\]
若 \(k_g\in S_h(M)\) 对某个 \(h\neq g\) 成立，则由上式知
\[
k_g\in S_g(M)\cap S_h(M)=S_{\gcd(g,h)}(M)
\]
且 \(d:=\gcd(g,h)<g\)。由于交集非空，\(g/d\) 必为奇数。于是 \(S_d(M)\subseteq S_g(M)\)，并且按定理 \ref{thm:fold-zero-coset-rigidity} 的显式描述，\(S_d(M)\) 在 \(S_g(M)\) 内对应于模 \(g/d\) 的单一余类
\[
t\equiv \frac{g/d-1}{2}\pmod{g/d}.
\]
而点 \(k_g\) 对应于参数 \(t=0\)。因为 \(g/d>1\) 且为奇数，
\[
\frac{g/d-1}{2}\not\equiv 0\pmod{g/d},
\]
故 \(k_g\notin S_d(M)\)，矛盾。

因此
\[
k_g\notin \bigcup_{h\in \mathcal G_m\setminus\{g\}}S_h(M),
\]
从而
\[
S_g(M)\nsubseteq \bigcup_{h\in \mathcal G_m\setminus\{g\}}S_h(M).
\]
于是任何生成整个 \(\mathcal Z_m\) 的子族都必须包含 \(g\)。由于 \(g\in \mathcal M_m\) 任意，遂得
\[
\mathcal M_m\subseteq \mathcal A.
\]
这就证明了唯一极小性。证毕。
\end{proof}

\begin{theorem}[极大奇商生成族上的谱零点集合精确容斥公式]\label{thm:xi-time-part72a-zero-spectrum-maximal-family-inclusion-exclusion}
沿用定理 \ref{thm:xi-time-part72a-zero-spectrum-maximal-odd-divisibility-compression} 的记号，写
\[
\mathcal M_m=\{g_1,\dots,g_r\}.
\]
则谱零点集合的基数满足精确公式
\[
|\mathcal Z_m|
=
\sum_{\varnothing\neq I\subseteq \{1,\dots,r\}}
(-1)^{|I|+1}
\mathbf 1_{\mathrm{comp}(I)}
\gcd(g_i:\ i\in I),
\]
其中相容指标 \(\mathbf 1_{\mathrm{comp}(I)}\) 定义为：当且仅当对全部 \(i,j\in I\) 都有
\[
2\gcd(g_i,g_j)\mid (g_j-g_i)
\]
时，\(\mathbf 1_{\mathrm{comp}(I)}=1\)；否则 \(\mathbf 1_{\mathrm{comp}(I)}=0\)。
\end{theorem}
\begin{proof}
由定理 \ref{thm:xi-time-part72a-zero-spectrum-maximal-odd-divisibility-compression}，
\[
\mathcal Z_m=\bigcup_{\nu=1}^{r}S_{g_\nu}(M).
\]
对该有限并应用包含--排斥，
\[
|\mathcal Z_m|
=
\sum_{\varnothing\neq I\subseteq \{1,\dots,r\}}
(-1)^{|I|+1}
\left|
\bigcap_{i\in I}S_{g_i}(M)
\right|.
\]
再由定理 \ref{thm:xi-time-part71-zero-coset-finite-intersection-pairwise-compatibility}，交截非空当且仅当上述两两相容条件成立，而一旦非空就有
\[
\left|
\bigcap_{i\in I}S_{g_i}(M)
\right|
=
\gcd(g_i:\ i\in I).
\]
代回即得所示公式。证毕。
\end{proof}

\begin{corollary}[极大生成族 pairwise-incompatible 时的精确直和公式]\label{cor:xi-time-part72a-zero-spectrum-maximal-family-pairwise-incompatible-exact-sum}
沿用定理 \ref{thm:xi-time-part72a-zero-spectrum-maximal-odd-divisibility-compression} 的记号。若极大生成族 \(\mathcal M_m=\{g_1,\dots,g_r\}\) 满足任意两元都不相容，即
\[
2\gcd(g_i,g_j)\nmid (g_j-g_i)
\qquad (i\neq j),
\]
则对应余类两两不交，从而
\[
|\mathcal Z_m|=\sum_{\nu=1}^{r}g_\nu.
\]
\end{corollary}
\begin{proof}
由定理 \ref{thm:xi-time-part71-zero-coset-finite-intersection-pairwise-compatibility}，上述条件等价于
\[
S_{g_i}(M)\cap S_{g_j}(M)=\varnothing
\qquad (i\neq j).
\]
故
\[
\mathcal Z_m=\bigsqcup_{\nu=1}^{r}S_{g_\nu}(M).
\]
再由定理 \ref{thm:fold-zero-coset-rigidity} 的 \(|S_g(M)|=g\)，可得
\[
|\mathcal Z_m|
=
\sum_{\nu=1}^{r}|S_{g_\nu}(M)|
=
\sum_{\nu=1}^{r}g_\nu.
\]
证毕。
\end{proof}

\noindent
上述结果表明，零余类族的算术几何还具有一层此前未被完全显化的“压缩表示”刚性：有限交截虽由全体候选余类的 pairwise 相容性控制，但真正生成 \(\mathcal Z_m\) 的并不需要整张约数图，而只需要奇商偏序下的极大元族；这一压缩不仅是存在性的，而且是唯一极小的。相应地，谱零点计数也可从全约数格的指标容斥，进一步压缩为极大生成族上的精确包含--排斥；而在 pairwise-incompatible 的情形下，全部交叠复杂性彻底消失，精确计数退化为一条纯粹的 Fibonacci 权重直和公式。于是，零余类交叠、包含、压缩表示与 exact counting 四者就在同一条 odd-divisibility 主线上获得了闭合。

\paragraph{pairwise-incompatible 零块族的同步信息削减}
在推论 \ref{cor:xi-time-part72-zero-spectrum-pairwise-incompatible-linear-lower-bound} 已把 pairwise-incompatible 约数族转化为 \(|\mathcal Z_m|\) 的线性可加增量，而定理 \ref{thm:xi-fold-zero-dyadic-tower-synchronous-information-gap-improvement} 已在 dyadic 塔特例下把这一增量同步传递到碰撞、KL、峰值与最大纤维偏离之后，可以看到：这一同步机制本身并不依赖 dyadic 结构，而只依赖 pairwise incompatibility 所保证的零块互不相交。

\begin{theorem}[pairwise-incompatible Fibonacci 约数族触发的同步信息削减]\label{thm:xi-time-part72b-pairwise-incompatible-synchronous-information-gap}
\leanverified{paper\_xi\_time\_part72b\_pairwise\_incompatible\_synchronous\_information\_gap}
沿用定理 \ref{thm:xi-time-part72-zero-spectrum-divisor-lattice-indicator-inclusion-exclusion} 的记号，并沿用定理 \ref{thm:fold-zero-uncertainty} 的记号。设
\[
\mathcal A\subseteq \mathcal D_m
\]
满足对任意不同的 \(d,e\in\mathcal A\)，都有
\[
2F_{\gcd(d,e)}\nmid (F_e-F_d).
\]
记
\[
L_{\mathcal A}:=\sum_{d\in\mathcal A}F_d.
\]
则
\[
|\mathcal Z_m|\ge L_{\mathcal A}.
\]
因而同时成立
\[
\mathsf{Col}_m
\le
1-\frac{L_{\mathcal A}}{F_{m+2}},
\qquad
D_{\mathrm{KL}}(p_m\|u_m)
\le
\log\!\bigl(F_{m+2}-L_{\mathcal A}\bigr),
\]
\[
S_m\le 2^m\sqrt{F_{m+2}-1-L_{\mathcal A}},
\]
\[
M_m-\bar d_m
\le
2^m\sqrt{\frac{F_{m+2}-1-L_{\mathcal A}}{F_{m+2}}}.
\]
\end{theorem}
\begin{proof}
由推论 \ref{cor:xi-time-part72-zero-spectrum-pairwise-incompatible-linear-lower-bound}，
\[
|\mathcal Z_m|\ge \sum_{d\in\mathcal A}F_d=L_{\mathcal A}.
\]
将该下界分别代入定理 \ref{thm:fold-collision-zero-reduction} 与定理
\ref{thm:fold-zero-uncertainty} 的第 \(2\)、\(3\)、\(4\) 条，便得到所示碰撞上界、KL 上界、Fourier 峰值上界以及最大纤维偏离上界。证毕。
\end{proof}

\noindent
因此，定理 \ref{thm:xi-fold-zero-dyadic-tower-synchronous-information-gap-improvement} 只是这一一般机制在 dyadic 塔族上的特例。真正驱动同步削减的并非 \(v_2\)-塔本身，而是 Fibonacci 零块之间的 pairwise incompatibility：只要若干约数对应的零块能够同时并入 \(\mathcal Z_m\) 而不发生交叠，它们所贡献的零点增量便会无损地传递到碰撞、相对熵、峰值振幅与最大纤维偏离这四个量上。

\endinput


\paragraph{二维原子层、非刚性迹系统、激活零尺度根与零温四态核心的三尺度分裂}
在真实输入 \(40\) 状态核的二维加权 \(\zeta\) 分母闭式、一参数输出位势的激活分支，以及零温四态地基层已经建立之后，还可把原子层、非刚性残余谱、局部激活层与零温核心层进一步压缩为下列若干条彼此衔接的结构后果。前两条仍对应对原子因子的中心投影式相位盲读出；紧随其后的两条把去除 rigid branches 后的全部非刚性迹压缩为单一八次 core 因子的 Newton--Cayley 闭系统，并进一步抽出 \(v=1\) 超曲面与 \((u,v)=(0,1)\) 尖点处的谱秩塌缩；随后的几条则转入相位敏感的纤维内读出，分别刻画 \(u=1\) 处的极点阶、Abel 有限部平移、激活分支的横截性、去支后的局部无碰撞窗口以及正实轴上的局部谱壁垒；最后再把局部奇性独占与零温 Perron 幅度统一到同一条谱分裂链中。

\begin{theorem}[二维加权 \texorpdfstring{$\zeta$}{zeta} 的原子因子在 ghost 层对碰撞权完全失明]\label{thm:xi-time-part73-twoparam-atomic-ghost-u-blindness}
\leanverified{paper\_xi\_time\_part73\_twoparam\_atomic\_ghost\_u\_blindness}
沿用定理 \ref{thm:real-input-40-det-uv-2plus8} 与推论 \ref{cor:real-input-40-zeta-uv-sqrtv-eigs} 的记号，并记
\[
\zeta(z;u,v)
=
\frac{1}{\Delta(z;u,v)}
=
\frac{1}{1-vz^2}\,\zeta_{\mathrm{core}}(z;u,v),
\qquad
\zeta_{\mathrm{core}}(z;u,v):=-P_8(z;u,v)^{-1}.
\]
定义原子 ghost 系数 \(a_n^{\mathrm{atom}}(u,v)\) 为
\[
\log\frac{1}{1-vz^2}
=
\sum_{n\ge 1}\frac{a_n^{\mathrm{atom}}(u,v)}{n}z^n.
\]
则对任意 \(n\ge 1\) 都有精确公式
\[
\boxed{\;
a_n^{\mathrm{atom}}(u,v)
=
\begin{cases}
2v^{n/2},& 2\mid n,\\[4pt]
0,& 2\nmid n,
\end{cases}\;}
\]
因而该原子层对碰撞参数 \(u\) 完全独立。若进一步取 \(v>0\)，则同一公式可重写为
\[
\boxed{\;
a_n^{\mathrm{atom}}(u,v)
=
(\sqrt v)^n+(-\sqrt v)^n.\;}
\]
特别地，所有奇阶 ghost/trace 审计都对因子 \((1-vz^2)^{-1}\) 完全失明：对任意 \(m\ge 0\)，其对
\[
\operatorname{Tr}(M_{u,v}^{\,2m+1})
\]
的原子贡献恒为 \(0\)。
\end{theorem}
\begin{proof}
由
\[
-\log(1-vz^2)=\sum_{k\ge 1}\frac{v^k}{k}z^{2k}
\]
并与 ghost 记号
\[
\log\frac{1}{1-vz^2}
=
\sum_{n\ge 1}\frac{a_n^{\mathrm{atom}}(u,v)}{n}z^n
\]
逐项比较系数，立得
\[
a_{2k}^{\mathrm{atom}}(u,v)=2v^k,
\qquad
a_{2k+1}^{\mathrm{atom}}(u,v)=0.
\]
这已经给出对 \(u\) 的完全独立性。

当 \(v>0\) 时，推论 \ref{cor:real-input-40-zeta-uv-sqrtv-eigs} 表明因子 \(1-vz^2\) 在谱上对应两条 rigid branches
\[
\lambda_\pm=\pm\sqrt v,
\]
故其对 \(n\)-步迹的贡献正是
\[
(\sqrt v)^n+(-\sqrt v)^n,
\]
而这与上式完全一致。于是奇数阶贡献必为零，从而任何只读取奇阶 ghost/trace 的审计函数都无法检测该原子因子。证毕。
\end{proof}

\input{subsubsec__xi-time-protocol-conclusions-part73b-nonrigid-trace-rank-collapse}

\begin{theorem}[二维加权 \texorpdfstring{$\zeta$}{zeta} 在碰撞零温缩放下坍缩为单一原子模型]\label{thm:xi-time-part73-twoparam-collision-freezing-single-atom-collapse}
沿用定理 \ref{thm:real-input-40-det-uv-2plus8} 的记号。固定 \(w>0\) 与 \(v>0\)，并令
\[
z=\sqrt{\frac{w}{u}},
\qquad
u\to+\infty.
\]
则有极限
\[
\boxed{\;
\lim_{u\to+\infty}\Delta\!\left(\sqrt{\frac{w}{u}};u,v\right)
=
1-vw.\;}
\]
更精确地，
\[
\boxed{\;
P_8\!\left(\sqrt{\frac{w}{u}};u,v\right)
=
-1+vw+O(u^{-1/2}),\;}
\]
从而
\[
\boxed{\;
\Delta\!\left(\sqrt{\frac{w}{u}};u,v\right)
=
1-vw+O(u^{-1/2}).\;}
\]
因此只要 \(vw\neq 1\)，便等价地有
\[
\boxed{\;
\zeta\!\left(\sqrt{\frac{w}{u}};u,v\right)
\longrightarrow
\frac{1}{1-vw}.\;}
\]
换言之，在碰撞零温缩放下，二维 \((u,v)\)-联合势的核心部分完全退化为长度 \(2\)、权重 \(v\) 的单原子 Euler 模型。
\end{theorem}
\begin{proof}
将定理 \ref{thm:real-input-40-det-uv-2plus8} 中 \(P_8(z;u,v)\) 的显式八次式代入
\[
z=\sqrt{\frac{w}{u}}.
\]
逐项估计可见：常数项给出 \(-1\)，而 \(z^2\) 项中的主贡献来自
\[
(uv+5v)z^2=vw+5vw\,u^{-1}.
\]
其余各项都至少带有 \(u^{-1/2}\) 的衰减；更具体地，\(z,z^3,z^5,z^7\) 项为 \(O(u^{-1/2})\)，而 \(z^4,z^6,z^8\) 项为 \(O(u^{-1})\)。故
\[
P_8\!\left(\sqrt{\frac{w}{u}};u,v\right)
=
-1+vw+O(u^{-1/2}).
\]

另一方面，
\[
vz^2-1
=
\frac{vw}{u}-1
=
-1+O(u^{-1}).
\]
于是
\[
\Delta\!\left(\sqrt{\frac{w}{u}};u,v\right)
=
\Bigl(-1+O(u^{-1})\Bigr)
\Bigl(-1+vw+O(u^{-1/2})\Bigr)
=
1-vw+O(u^{-1/2}),
\]
从而得到第一条极限。若 \(vw\neq 1\)，则极限不为零，故取倒数即得 \(\zeta\) 的极限公式。证毕。
\end{proof}

\begin{theorem}[\texorpdfstring{$u=1$}{u=1} 处 \texorpdfstring{$\zeta$}{zeta} 的双--三极点律]\label{thm:xi-time-part73-u1-zeta-double-triple-pole}
\leanverified{paper\_xi\_time\_part73\_u1\_zeta\_double\_triple\_pole}
沿用命题 \ref{prop:real-input-40-output-spectral-collisions} 的记号，并记
\[
\zeta(z,1):=\Delta(z,1)^{-1},
\qquad
\zeta_{\mathrm{core}}(z,1):=Q(z,1)^{-1}.
\]
则有精确分解
\[
\Delta(z,1)=-(z-1)^2(z+1)^3(z^2-3z+1)(z^2-z-1).
\]
因此 \(\zeta(z,1)\) 在 \(z=1\) 具有二阶极点，在 \(z=-1\) 具有三阶极点。去除原子 Witt 因子 \((1-z^2)^{-1}\) 之后，核心 \(\zeta_{\mathrm{core}}(z,1)\) 在 \(z=1\)、\(z=-1\) 分别仅余一阶与二阶极点。
\end{theorem}
\begin{proof}
由命题 \ref{prop:real-input-40-output-spectral-collisions}，
\[
Q(z,1)=(z-1)(z+1)^2(z^2-3z+1)(z^2-z-1).
\]
另一方面，
\[
\Delta(z,1)=(1-z^2)Q(z,1)=-(z-1)(z+1)Q(z,1),
\]
将两式相乘即得所示因式分解。于是 \(\zeta(z,1)=\Delta(z,1)^{-1}\) 在 \(z=1\)、\(z=-1\) 的极点阶分别为 \(2\) 与 \(3\)。而
\[
\zeta_{\mathrm{core}}(z,1)=Q(z,1)^{-1}
\]
只除去原子因子 \((1-z^2)^{-1}\)，故两处极点阶分别降低 \(1\)。证毕。
\end{proof}

\begin{theorem}[原子二周期对 Abel 有限部的精确平移]\label{thm:xi-time-part73-two-cycle-abel-finitepart-shift}
沿用定理 \ref{thm:xi-time-part65-atomic-primitive-abel-shift} 的记号。在 \(u=1\) 的同一核上，设 \(z_\star=\varphi^{-2}\)。则 primitive--core 分解满足
\[
P_{\mathrm{core}}(z)=P(z)-z^2,
\qquad
\log\mathfrak M=\log\mathfrak M_{\mathrm{core}}+z_\star^2.
\]
因此
\[
\log\mathfrak M-\log\mathfrak M_{\mathrm{core}}
=
\varphi^{-4}
=
\frac{7-3\sqrt5}{2},
\]
亦即
\[
\mathfrak M_{\mathrm{core}}=e^{-\varphi^{-4}}\mathfrak M.
\]
\end{theorem}
\begin{proof}
由定理 \ref{thm:xi-time-part65-atomic-primitive-abel-shift}，
\[
P(z,u)=uz^2+P_{\mathrm{core}}(z,u).
\]
在 \(u=1\) 处专门化即得
\[
P_{\mathrm{core}}(z)=P(z)-z^2.
\]
同一定理并给出
\[
\log\mathfrak M=\log\mathfrak M_{\mathrm{core}}+z_\star^2.
\]
又在本核上 \(z_\star=\varphi^{-2}\)，故
\[
z_\star^2=\varphi^{-4}=\frac{7-3\sqrt5}{2}.
\]
代回并指数化，便得到所示 Abel 有限部平移公式。证毕。
\end{proof}

\begin{theorem}[激活分支的横截性与二次项消失]\label{thm:xi-time-part73-activated-branch-unique-zero-scale-singularity}
沿用命题 \ref{prop:real-input-40-activated-branch-linear}、注 \ref{rem:real-input-40-activated-branch-series} 与定理 \ref{thm:xi-output-potential-activated-branch-rigidity} 的记号，并记
\[
G(\lambda,u):=\lambda^8F(\lambda^{-1},u).
\]
则在 \(u=1\) 的某邻域内存在唯一解析分支 \(\lambda_{\mathrm{act}}(u)\) 满足
\[
G(\lambda_{\mathrm{act}}(u),u)=0,
\qquad
\lambda_{\mathrm{act}}(1)=0,
\]
并且其 Taylor 展开满足
\[
\boxed{\;
\lambda_{\mathrm{act}}(u)
=(u-1)+3(u-1)^3-(u-1)^4+18(u-1)^5-22(u-1)^6
+O((u-1)^7).\;}
\]
特别地，二次项严格消失，且
\[
\lambda'_{\mathrm{act}}(1)=1,
\]
从而该分支穿过 \(\lambda=0\) 的方式是横截而非抛物切触。进一步，存在常数 \(\rho>0\) 与 \(\varepsilon>0\)，使得对一切满足 \(|u-1|<\varepsilon\) 的实参数，都有
\[
\boxed{\;
|\lambda_j(u)|\ge \rho
\qquad
(2\le j\le 8).\;}
\]
\end{theorem}
\begin{proof}
由命题 \ref{prop:real-input-40-activated-branch-linear}，
\[
\partial_\lambda G(0,1)=1\neq 0,
\qquad
\lambda_{\mathrm{act}}(1)=0.
\]
故由隐函数定理，\(G(\lambda,u)=0\) 在 \((0,1)\) 邻域内确定唯一解析分支 \(\lambda_{\mathrm{act}}(u)\)，且
\[
\lambda'_{\mathrm{act}}(1)=1.
\]
另一方面，注 \ref{rem:real-input-40-activated-branch-series} 给出
\[
\lambda_{\mathrm{act}}(u)
=(u-1)+3(u-1)^3-(u-1)^4+18(u-1)^5-22(u-1)^6
+O((u-1)^7),
\]
故二次项系数严格为 \(0\)，而线性项系数为 \(1\)，从而该分支以横截方式穿过 \(\lambda=0\)。

最后，定理 \ref{thm:xi-output-potential-activated-branch-rigidity} 已证明：在 \((\lambda,u)=(0,1)\) 的某邻域内，唯一的小根就是 \(\lambda_{\mathrm{act}}(u)\)，其余七条根分支全部远离 \(0\)。因此存在 \(r>0\) 与 \(\varepsilon_0>0\)，使得当 \(|u-1|<\varepsilon_0\) 时，除 \(\lambda_{\mathrm{act}}(u)\) 外的所有根都落在 \(|\lambda|\ge r\) 上。取 \(\rho:=r\)、\(\varepsilon:=\varepsilon_0\) 即得最后一条。证毕。
\end{proof}

\begin{theorem}[激活分支的三次归一化刚性与去支实谱的无碰撞窗口]\label{thm:xi-time-part73-activated-branch-deflation-window}
沿用定理 \ref{thm:xi-time-part73-activated-branch-unique-zero-scale-singularity} 的记号，并记
\[
I_\ast:=
\bigl(0.93283726248596492776\ldots,\ 1.0547964213679192275\ldots\bigr).
\]
定义三次归一化 defect germ
\[
\Xi(u):=\frac{\lambda_{\mathrm{act}}(u)-(u-1)}{(u-1)^3}.
\]
又定义去支因子
\[
\widetilde G(\lambda,u):=\frac{G(\lambda,u)}{\lambda-\lambda_{\mathrm{act}}(u)}.
\]
则成立：
\begin{enumerate}
  \item \(\Xi(u)\) 在 \(u=1\) 处可解析延拓，并满足
  \[
  \Xi(1)=3.
  \]
  \item \(\widetilde G\) 在 \((\lambda,u)=(0,1)\) 的某邻域内解析；
  \item 对每个实参数 \(u\in I_\ast\setminus\{1\}\)，多项式 \(\widetilde G(\cdot,u)\) 的全部实根都是单根。换言之，去除激活分支之后，审计窗口 \(I_\ast\) 内的剩余实谱不再发生双根碰撞。
\end{enumerate}
\end{theorem}
\begin{proof}
由注 \ref{rem:real-input-40-activated-branch-series}，
\[
\lambda_{\mathrm{act}}(u)-(u-1)
=
3(u-1)^3-(u-1)^4+18(u-1)^5-22(u-1)^6+O((u-1)^7).
\]
故
\[
\Xi(u)
=
3-(u-1)+18(u-1)^2-22(u-1)^3+O((u-1)^4),
\]
从而 \(\Xi\) 在 \(u=1\) 处具有可去奇性，并且
\[
\Xi(1)=3.
\]

第二条正是定理 \ref{thm:xi-output-potential-activated-branch-rigidity} 的 Weierstrass 因除结论。

最后，设某个 \(u_0\in I_\ast\setminus\{1\}\) 使 \(\widetilde G(\cdot,u_0)\) 存在实双根 \(\lambda_0\)。由于
\[
G(0,u)=u^3(1-u),
\]
且 \(u_0\neq 1\)，故 \(\lambda_0\neq 0\)。记
\[
z_0:=\lambda_0^{-1}\in \RR.
\]
在 \(\lambda\neq 0\) 上，变换 \(z=\lambda^{-1}\) 为解析双射，而
\[
G(\lambda,u)=\lambda^8F(\lambda^{-1},u)
\]
只差一个处处非零因子 \(\lambda^8\)，故 \(\lambda_0\) 为 \(\widetilde G(\cdot,u_0)\) 的实双根会强迫 \(z_0\) 成为 core 因子 \(F(z,u_0)\) 的实双根；沿一参数输出位势的记号，这等价于 \(Q(z,u_0)\) 在实轴上出现重根。可这与注 \ref{rem:real-input-40-output-potential-degeneracy-phase} 矛盾：该注指出在正实轴上除 \(u=1\) 外，全部退化点恰为
\[
0.008460\ldots,\ 0.745534\ldots,\ 0.932837\ldots,\ 1.054796\ldots,\ 2.090392\ldots,\ 15.522820\ldots,
\]
因而在开区间 \(I_\ast\setminus\{1\}\) 内不存在任何重根参数。故 \(\widetilde G(\cdot,u)\) 的实根在该窗口内全为单根。证毕。
\end{proof}

\begin{corollary}[围绕 \texorpdfstring{$u=1$}{u=1} 的局部谱病态完全由唯一激活小根承担]\label{cor:xi-time-part73-single-activated-branch-all-local-illconditioning}
沿用定理 \ref{thm:xi-time-part73-activated-branch-unique-zero-scale-singularity} 与定理 \ref{thm:xi-time-part73-activated-branch-deflation-window} 的记号。则存在 \(u=1\) 的邻域 \(\mathcal U\) 与在 \((\lambda,u)=(0,1)\) 某邻域内解析的函数 \(\widetilde G(\lambda,u)\)，使得
\[
G(\lambda,u)=\bigl(\lambda-\lambda_{\mathrm{act}}(u)\bigr)\widetilde G(\lambda,u)
\qquad(u\in\mathcal U).
\]
并且存在常数 \(c>0\)，使得对每个 \(u\in\mathcal U\)，\(\widetilde G(\cdot,u)\) 的任一零点 \(\mu(u)\) 都满足
\[
|\mu(u)|\ge c.
\]
因此，\(u=1\) 邻域内全部由小根诱导的局部谱病态都由唯一分支 \(\lambda_{\mathrm{act}}(u)\) 承担；去除该分支之后，其余七支统一恢复到 \(O(1)\) 尺度。
\end{corollary}
\begin{proof}
定理 \ref{thm:xi-time-part73-activated-branch-deflation-window} 的第二条给出
\[
\widetilde G(\lambda,u):=\frac{G(\lambda,u)}{\lambda-\lambda_{\mathrm{act}}(u)}
\]
在 \((\lambda,u)=(0,1)\) 的某邻域内解析，故所示 Weierstrass 因除成立。

另一方面，定理 \ref{thm:xi-time-part73-activated-branch-unique-zero-scale-singularity} 已证明：在 \(u=1\) 的某邻域内，除 \(\lambda_{\mathrm{act}}(u)\) 外，其余全部根都与 \(0\) 保持一致正距离。于是可取某个常数 \(c>0\)，使 \(\widetilde G(\cdot,u)\) 的所有零点统一满足
\[
|\mu(u)|\ge c.
\]
这恰表明全部局部小尺度病态只由 \(\lambda_{\mathrm{act}}(u)\) 一支承担，而去支后的剩余谱支都停留在 \(O(1)\) 距离之外。证毕。
\end{proof}

\begin{theorem}[正实轴上的谱“护城河”]\label{thm:xi-time-part73-positive-real-spectral-moat}
仍在真实输入 \(40\) 状态核的输出位势加权下，取
\[
\Delta(z,u)=(1-uz^2)Q(z,u),
\qquad
u>0.
\]
则正实参数区间
\[
1<u<4
\]
构成一个严格的谱护城河：在该区间内同时不存在显式原子分支 \(z=\pm u^{-1/2}\) 与 \(Q\)-谱的碰撞，也不存在 \(Q\)-谱在整数点 \(z=\pm1\) 的共振。更精确地，
\[
Q(\pm u^{-1/2},u)=0
\iff
u=1,
\]
且
\[
Q(1,u)=-u(u-4)(u-1)(u+1),
\qquad
Q(-1,u)=-(u-1)^2(u^2+6u-2).
\]
故在 \(u>1\) 的正实轴上，首个整数级共振恰发生于
\[
u=4,
\qquad
\lambda=1.
\]
\end{theorem}
\begin{proof}
命题 \ref{prop:real-input-40-output-spectral-collisions} 已给出
\[
Q\!\left(\frac{1}{\sqrt u},u\right)=\frac{2(\sqrt u-1)(\sqrt u+1)^2}{u^{3/2}},
\qquad
Q\!\left(-\frac{1}{\sqrt u},u\right)=\frac{2(\sqrt u-1)^2(\sqrt u+1)}{u^{3/2}},
\]
因而在 \(u>0\) 上，
\[
Q(\pm u^{-1/2},u)=0 \iff u=1.
\]
这说明在 \(1<u<4\) 内不存在原子分支与 \(Q\)-谱的碰撞。

同一命题又给出
\[
Q(1,u)=-u(u-4)(u-1)(u+1),
\qquad
Q(-1,u)=-(u-1)^2(u^2+6u-2).
\]
当 \(1<u<4\) 时，前式显然非零；而
\[
u^2+6u-2>0
\qquad (u>1),
\]
故后式亦非零。因此在 \(1<u<4\) 内不存在 \(Q\)-谱于 \(z=\pm1\) 的整数级共振。最后，\(u=4\) 时 \(Q(1,4)=0\)，从而给出首个正实整数级共振 \(\lambda=1\)。证毕。
\end{proof}

\begin{theorem}[主 Perron 根相对于原子谱支存在严格 \texorpdfstring{$\sqrt u$}{sqrt(u)} 级分离，并从上方逼近零温四态 Perron 常数]\label{thm:xi-time-part73-perron-vs-atomic-sqrtu-separation}
沿用推论 \ref{cor:real-input-40-factor-sqrtu-eigs}、命题 \ref{prop:real-input-40-pressure-freezing} 与命题 \ref{prop:real-input-40-output-mirror-branch} 的记号。记 \(\lambda(u)\) 为输出势的 Perron 正根，并令
\[
P(\theta):=\log\lambda(e^\theta),
\qquad
c^2=\rho(B_\infty),
\]
其中 \(c\approx 1.89635299351377\)、\(d\approx 0.807943317323508\)，且 \(c^2\) 是四次方程
\[
y^4-6y^3+9y^2-y-1=0
\]
的最大正根。则当 \(u\to+\infty\) 时有：
\[
\boxed{\;
\lambda(u)-\sqrt u
=
(c-1)\sqrt u+d+O(u^{-1/2}),\;}
\]
因而
\[
\boxed{\;
\lambda(u)-\sqrt u\longrightarrow +\infty.\;}
\]
进一步，
\[
\boxed{\;
\frac{\lambda(u)^2}{u}
=
\rho(B_\infty)+\frac{2cd}{\sqrt u}+O(u^{-1}),\;}
\]
特别地，由于 \(d>0\)，该比值从上方逼近 \(\rho(B_\infty)\)。
最后，若取 \(u=e^\theta\)，则
\[
\boxed{\;
P(\theta)-\frac{\theta}{2}
=
\log c+\frac{d}{c}e^{-\theta/2}+O(e^{-\theta})
\longrightarrow
\log c>0.\;}
\]
因此显式原子谱支 \(\pm\sqrt u\) 只决定零温端点的斜率 \(1/2\)，而主导自由能的幅度常数则完全由零温四态核心 \(B_\infty\) 决定。
\end{theorem}
\begin{proof}
由命题 \ref{prop:real-input-40-output-mirror-branch} 的 Puiseux 展开，
\[
\lambda(u)=c\sqrt u+d+\frac{e}{\sqrt u}+\frac{f}{u}+O(u^{-3/2}).
\]
与显式原子谱支 \(\sqrt u\) 相减，便得到
\[
\lambda(u)-\sqrt u
=
(c-1)\sqrt u+d+O(u^{-1/2}).
\]
由于 \(c>1\)，主项 \((c-1)\sqrt u\) 严格为正并趋于无穷，从而得到第一条。

再将同一展开平方：
\[
\lambda(u)^2
=
c^2u+2cd\sqrt u+O(1).
\]
两边同除以 \(u\)，并用 \(c^2=\rho(B_\infty)\)，即得
\[
\frac{\lambda(u)^2}{u}
=
\rho(B_\infty)+\frac{2cd}{\sqrt u}+O(u^{-1}).
\]
由于 \(d>0\)，该展开表明收敛从上方发生。

最后令 \(u=e^\theta\)。由
\[
\lambda(e^\theta)
=
c\,e^{\theta/2}
\left(
1+\frac{d}{c}e^{-\theta/2}+O(e^{-\theta})
\right)
\]
可得
\[
P(\theta)
=
\log\lambda(e^\theta)
=
\frac{\theta}{2}+\log c+\frac{d}{c}e^{-\theta/2}+O(e^{-\theta}),
\]
这就给出最后一条公式。证毕。
\end{proof}

\input{subsubsec__xi-time-protocol-conclusions-part73a-vzero-collapse-atomic-variable-interface}

\input{subsubsec__xi-time-protocol-conclusions-part73c-fixed-parameter-evenlength-necklace}

\noindent
以上这些结果表明，真实输入 \(40\) 状态核在二维原子层、非刚性残余谱、\(u=1\) 的局部激活层与 \(u\to+\infty\) 的零温核心层之间存在严格的多尺度分裂：一方面，二维加权 \(\zeta\) 的单原子因子在 ghost 层只向偶长度注入 \(2v^{n/2}\)，并对碰撞权 \(u\) 完全失明；去除两条 rigid branches \(\pm\sqrt v\) 之后，其余全部非刚性迹又自动坍缩为由单一八次 core 因子控制的 Newton--Cayley 递推系统，并在 \(v=1\) 超曲面与 \((u,v)=(0,1)\) 尖点上继续发生可量化的谱秩塌缩；当进一步压到输出完全抑制的切片 \(v=0\) 时，整个二维谱又严格塌缩为唯一一个长度 \(1\) 的 primitive 轨道。另一方面，在碰撞零温缩放下，整个二维核心因子进一步退化为同一原子的单参数模型 \(1-vw\)；在一参数输出位势上，同一原子 Witt 机制只在 primitive 的 \(n=2\) 层激活，并把 Abel 有限部刚性平移为 \(u z_\star(u)^2\)；再者，激活分支以无二次项的横截方式穿过 \(\lambda=0\)，其三次归一化 defect 在 \(u=1\) 处锁定为常数 \(3\)，去支后的剩余实谱则在最近退化点所围成的审计窗口内恢复为无碰撞解析族；而正实轴上的区间 \(1<u<4\) 又把原子碰撞与整数级共振共同隔离在外。最后，真正主导大 \(u\) 幅度常数的并不是显式原子谱支 \(\pm\sqrt u\)，而是零温四态核心 \(B_\infty\) 的 Perron 常数。因而在读出口径上，前述几条对应于中心投影式的相位盲读出，而其后的非刚性迹系统、局部极点、激活分支与谱壁垒则共同构成相位敏感的纤维内读出。

\input{subsubsec__xi-time-protocol-conclusions-part74-leyang-primeslice-smith-periodized-ellipse-split}

\endinput


\endinput


\endinput


\endinput
